% VI : produire un rapport — Français 1re Tech — Cours signé OnBuch+
% Programme officiel MINESEC. Compiler avec : tectonic N22-vi-produire-rapport.tex
\newcommand{\DOCMATIERE}{Français}
\newcommand{\DOCNIVEAU}{1re Tech}
\newcommand{\DOCMODULE}{Français – classe de Première technique}
\newcommand{\DOCLECON}{N22}
\newcommand{\DOCTITRE}{VI : produire un rapport}
% OnBuch+ — préambule « cours signés » ENRICHI (compilé avec tectonic / XeLaTeX)
% Système visuel « Pop » d'OnBuch+ : fond crème (jamais blanc), bordures encre
% épaisses, ombres « dures » décalées, titres Archivo Black en capitales.
% Version MATHÉMATIQUES : graphes (pgfplots/tikz), boîtes pédagogiques
% (définition, propriété, méthode, attention, le savais-tu, exercice résolu,
% exercice, TP, bilan), page de garde et signature OnBuch+ ancrée en bas.
%
% Macros attendues dans main.tex AVANT \input{preamble} :
%   \DOCMATIERE  \DOCNIVEAU  \DOCMODULE  \DOCLECON (numéro)  \DOCTITRE
\documentclass[11pt]{article}

\usepackage{fontspec}
\usepackage[french]{babel}
\usepackage[a4paper,margin=2cm,top=3.4cm,bottom=2.8cm,headheight=1.4cm,footskip=1.3cm]{geometry}
\usepackage{amsmath,amssymb,amsfonts}
\usepackage{mathtools} % \xmapsto, \coloneqq... souvent utilisés par les modèles
\usepackage{cancel} % simplifications barrées (bilans de forces)
% \abs{z} (module, valeur absolue) et \norm{u} (norme), utilisés par les modèles
\providecommand{\abs}{}\let\abs\relax\DeclarePairedDelimiter{\abs}{\lvert}{\rvert}
\providecommand{\norm}{}\let\norm\relax\DeclarePairedDelimiter{\norm}{\lVert}{\rVert}
\usepackage[table]{xcolor} % [table] : fournit aussi \rowcolors (alternance de lignes)
\usepackage{pagecolor}
\usepackage{tikz}
\usetikzlibrary{arrows.meta,positioning,calc,shapes.geometric,shapes.misc,shapes.arrows,decorations.pathmorphing,decorations.pathreplacing,decorations.markings,patterns,shadows,mindmap,trees,fit,backgrounds,matrix,intersections,angles,quotes}
\usepackage{pgfplots}
\usepackage[version=4]{mhchem} % équations nucléaires \ce{^{235}_{92}U}
\usepackage[european,siunitx]{circuitikz} % schémas électriques
\pgfplotsset{compat=1.18}
\usepgfplotslibrary{fillbetween,groupplots} % aires entre courbes (calcul intégral)
\usepackage{enumitem}
\usepackage{fancyhdr}
% [most] charge toutes les bibliothèques tcolorbox (listings, minted, raster,
% documentation...) et gonfle inutilement la mémoire TeX sur un document long
% (~300 boîtes) : on ne charge que ce qui est réellement utilisé.
\usepackage[skins,breakable]{tcolorbox}
\usepackage{graphicx}
\usepackage{colortbl}
\usepackage{booktabs}
\usepackage{tabularx}
\usepackage{multirow}
\usepackage{diagbox} % case d'en-tête diagonale des tableaux à double entrée
% Colonnes extensibles centrée / gauche / droite (C, L, R), souvent utilisées par les modèles
\newcolumntype{C}{>{\centering\arraybackslash}X}
\newcolumntype{L}{>{\raggedright\arraybackslash}X}
\newcolumntype{R}{>{\raggedleft\arraybackslash}X}
\usepackage{array}
\usepackage{multicol}
\usepackage{siunitx}
% parse-numbers=false : accepte aussi \SI{2,5 \times 10^{-3}}{...}
\sisetup{locale=FR,output-decimal-marker={,},group-minimum-digits=4,parse-numbers=false}
\DeclareSIUnit\ohm{\ensuremath{\symup{\Omega}}} % U+2126 absent de la police
\DeclareSIUnit\division{div} % réglages d'oscilloscope (ms/div, V/div)
\DeclareSIUnit\tr{tr} % tours (vitesse de rotation, ex. \SI{1500}{\tr\per\minute})
\DeclareSIUnit\molar{M} % concentration molaire (ex. \SI{3}{\molar}, \SI{1}{\micro\molar})
\DeclareSIUnit\an{an} % année (le modèle confond parfois avec \year, un compteur TeX) — ex. \SI{5}{\kilo\meter\per\an}
\DeclareSIUnit\FCFA{FCFA} % franc CFA (ex. \SI{500}{\FCFA\per\kilo\gram})
\DeclareSIUnit\degreeBrix{\textdegree Brix} % teneur en sucre d'un sirop/jus (réfractométrie)
\DeclareSIUnit\month{mois} % le modèle utilise parfois \month, un compteur TeX, comme unité de durée
\DeclareSIUnit\microliter{\micro\liter} % le modèle écrit parfois \microliter en un seul token
\DeclareSIUnit\million{millions} % dénombrement (ex. \SI{15}{\million\per\mL} spermatozoïdes)
\usepackage{qrcode}
% \degree : commande fréquemment utilisée par les modèles pour un angle ou une
% température en dehors de \SI{}{} (non fournie par les paquets chargés).
\providecommand{\degree}{^{\circ}}
% \qed (fin de démonstration), utilisé par les modèles sans amsthm
\providecommand{\qed}{\hfill$\square$}
% \barre : double barre des valeurs interdites dans un tableau de variations (tkz-tab)
\providecommand{\barre}{\ensuremath{\|}}
% environnement proof (amsthm non chargé) utilisé par les modèles
\ifdefined\proof\else\newenvironment{proof}[1][Démonstration]{\par\noindent\textit{#1.}\ }{\qed\par}\fi
\usepackage{lastpage}
\usepackage{hyperref}
\hypersetup{hidelinks}

% ============================================================
% Palette « Pop » OnBuch+ — mêmes hex que la classe P (plus_ui.dart)
% ============================================================
\definecolor{poporange}{HTML}{F59321}
\definecolor{popdark}{HTML}{E07A0C}
\definecolor{popcream}{HTML}{FFF6E8}
\definecolor{popink}{HTML}{1C1714}
\definecolor{poppurple}{HTML}{7A5AE0}
\definecolor{popgreen}{HTML}{2E9E6A}
\definecolor{poppink}{HTML}{E0457B}
\definecolor{popblue}{HTML}{2F7DE1}
\definecolor{popgold}{HTML}{C98A12}
\definecolor{popmuted}{HTML}{8A7F76}
\definecolor{popline}{HTML}{EADFD2}
\definecolor{popgray}{HTML}{8A7F76}
\colorlet{popgrey}{popgray}
% Alias tolérants (le modèle nomme parfois les couleurs différemment)
\colorlet{popwhite}{white}
\colorlet{popblack}{popink}
\colorlet{popred}{poppink}
\colorlet{popyellow}{popgold}
% Teintes claires (fonds de tableaux/encadrés) — le modèle nomme parfois
% les couleurs avec un suffixe L (ex. popblueL) sans les avoir définies.
\colorlet{poporangeL}{poporange!20}
\colorlet{popdarkL}{popdark!20}
\colorlet{poppurpleL}{poppurple!20}
\colorlet{popgreenL}{popgreen!20}
\colorlet{poppinkL}{poppink!20}
\colorlet{popblueL}{popblue!20}
\colorlet{popgoldL}{popgold!20}
\colorlet{popredL}{popred!20}
\colorlet{popgrayL}{popgray!20}
% Teintes claires pour fonds de tableaux / zones
\colorlet{popblueL}{popblue!10!white}
\colorlet{poporangeL}{poporange!14!white}
\colorlet{popgreenL}{popgreen!12!white}
\colorlet{poppurpleL}{poppurple!12!white}
\colorlet{poppinkL}{poppink!12!white}
\colorlet{popgoldL}{popgold!14!white}

\pagecolor{popcream}

% ============================================================
% Typographie
% ============================================================
\defaultfontfeatures{Ligatures=TeX}
\setmainfont{PlusJakartaSans-Medium.ttf}[Path=../../../fonts/,BoldFont=PlusJakartaSans-Bold.ttf]
% Maths en sans-serif (Fira Math) avec chiffres et lettres droites de Plus
% Jakarta, pour que les formules se fondent dans le texte.
\usepackage[math-style=ISO,bold-style=ISO]{unicode-math}
\setmathfont{FiraMath-Regular.otf}[Path=../../../fonts/]
\setmathfont{PlusJakartaSans-Medium.ttf}[Path=../../../fonts/,range={up/num,up/{Latin,latin},\mathup}]
\usepackage{icomma} % virgule décimale sans espace en mode math
% Caractères mathématiques Unicode absents des polices de texte (Plus Jakarta,
% Archivo Black) : rendus par leur équivalent mathématique, en texte comme en
% formule — sinon ils s'affichent en carré vide (ex. « Arithmétique dans ℤ »).
\usepackage{newunicodechar}
\newunicodechar{ℝ}{\ensuremath{\mathbb{R}}}
\newunicodechar{ℤ}{\ensuremath{\mathbb{Z}}}
\newunicodechar{ℂ}{\ensuremath{\mathbb{C}}}
\newunicodechar{ℕ}{\ensuremath{\mathbb{N}}}
\newunicodechar{ℚ}{\ensuremath{\mathbb{Q}}}
\newunicodechar{𝔻}{\ensuremath{\mathbb{D}}}
\newunicodechar{α}{\ensuremath{\alpha}}
\newunicodechar{β}{\ensuremath{\beta}}
\newunicodechar{γ}{\ensuremath{\gamma}}
\newunicodechar{δ}{\ensuremath{\delta}}
\newunicodechar{ε}{\ensuremath{\varepsilon}}
\newunicodechar{θ}{\ensuremath{\theta}}
\newunicodechar{λ}{\ensuremath{\lambda}}
\newunicodechar{μ}{\ensuremath{\mu}}
\newunicodechar{σ}{\ensuremath{\sigma}}
\newunicodechar{τ}{\ensuremath{\tau}}
\newunicodechar{φ}{\ensuremath{\varphi}}
\newunicodechar{ψ}{\ensuremath{\psi}}
\newunicodechar{ω}{\ensuremath{\omega}}
\newunicodechar{Σ}{\ensuremath{\Sigma}}
% majuscules produites par \MakeUppercase dans les titres (α-aminés -> Α)
\newunicodechar{Α}{\ensuremath{\alpha}}
\newunicodechar{Β}{\ensuremath{\beta}}
\newunicodechar{Γ}{\ensuremath{\Gamma}}
\newunicodechar{Δ}{\ensuremath{\Delta}}
\newunicodechar{Ε}{\ensuremath{\varepsilon}}
\newunicodechar{Θ}{\ensuremath{\Theta}}
\newunicodechar{Λ}{\ensuremath{\Lambda}}
\newunicodechar{Μ}{\ensuremath{\mu}}
\newunicodechar{Τ}{\ensuremath{\tau}}
\newunicodechar{Φ}{\ensuremath{\Phi}}
\newunicodechar{Ψ}{\ensuremath{\Psi}}
\newunicodechar{Ω}{\ensuremath{\Omega}}
\newunicodechar{↦}{\ensuremath{\mapsto}}
\newunicodechar{∈}{\ensuremath{\in}}
\newunicodechar{∉}{\ensuremath{\notin}}
\newunicodechar{∧}{\ensuremath{\wedge}}
\newunicodechar{⇔}{\ensuremath{\Leftrightarrow}}
\newunicodechar{⟺}{\ensuremath{\Longleftrightarrow}}
\newunicodechar{⇒}{\ensuremath{\Rightarrow}}
\newunicodechar{⟹}{\ensuremath{\Longrightarrow}}
\newunicodechar{∘}{\ensuremath{\circ}}
\newunicodechar{∪}{\ensuremath{\cup}}
\newunicodechar{∩}{\ensuremath{\cap}}
\newunicodechar{≡}{\ensuremath{\equiv}}
\newunicodechar{⊂}{\ensuremath{\subset}}
\newunicodechar{⊥}{\ensuremath{\perp}}
\newunicodechar{∃}{\ensuremath{\exists}}
\newunicodechar{∀}{\ensuremath{\forall}}
\newunicodechar{∛}{\ensuremath{\sqrt[3]{\,}}}
\newunicodechar{∜}{\ensuremath{\sqrt[4]{\,}}}
\newunicodechar{⃗}{\ensuremath{{}^{\to}}}
\newunicodechar{ⁿ}{\ifmmode^{n}\else\textsuperscript{\ensuremath{n}}\fi}
\newunicodechar{ˣ}{\ifmmode^{x}\else\textsuperscript{\ensuremath{x}}\fi}
\newunicodechar{ᵇ}{\ifmmode^{b}\else\textsuperscript{\ensuremath{b}}\fi}
\newunicodechar{ʳ}{\ifmmode^{r}\else\textsuperscript{\ensuremath{r}}\fi}
\newunicodechar{ᵘ}{\ifmmode^{u}\else\textsuperscript{\ensuremath{u}}\fi}
\newunicodechar{ʸ}{\ifmmode^{y}\else\textsuperscript{\ensuremath{y}}\fi}
\newunicodechar{ᵃ}{\ifmmode^{a}\else\textsuperscript{\ensuremath{a}}\fi}
\newunicodechar{ᵉ}{\ifmmode^{e}\else\textsuperscript{\ensuremath{e}}\fi}
\newunicodechar{ᵏ}{\ifmmode^{k}\else\textsuperscript{\ensuremath{k}}\fi}
\newunicodechar{ᵖ}{\ifmmode^{p}\else\textsuperscript{\ensuremath{p}}\fi}
\newunicodechar{ᴺ}{\ifmmode^{N}\else\textsuperscript{\ensuremath{N}}\fi}
\newunicodechar{ᶜ}{\ifmmode^{c}\else\textsuperscript{\ensuremath{c}}\fi}
\newunicodechar{ʰ}{\ifmmode^{h}\else\textsuperscript{\ensuremath{h}}\fi}
\newunicodechar{ᵠ}{\ifmmode^{q}\else\textsuperscript{\ensuremath{q}}\fi}
\newunicodechar{ₙ}{\ifmmode_{n}\else\textsubscript{\ensuremath{n}}\fi}
\newunicodechar{ₐ}{\ifmmode_{a}\else\textsubscript{\ensuremath{a}}\fi}
\newunicodechar{ᵢ}{\ifmmode_{i}\else\textsubscript{\ensuremath{i}}\fi}
\newunicodechar{ᵣ}{\ifmmode_{r}\else\textsubscript{\ensuremath{r}}\fi}
\newunicodechar{ₖ}{\ifmmode_{k}\else\textsubscript{\ensuremath{k}}\fi}
\newunicodechar{ᵦ}{\ifmmode_{\beta}\else\textsubscript{\ensuremath{\beta}}\fi}
\newunicodechar{ₓ}{\ifmmode_{x}\else\textsubscript{\ensuremath{x}}\fi}
\newfontfamily\popdisplay{ArchivoBlack-Regular.ttf}[Path=../../../fonts/]
\newfontfamily\poplogo{SpaceGrotesk-Bold.ttf}[Path=../../../fonts/]
\newfontfamily\popsemi{PlusJakartaSans-SemiBold.ttf}[Path=../../../fonts/]
\newfontfamily\popxbold{PlusJakartaSans-ExtraBold.ttf}[Path=../../../fonts/]
\newfontfamily\popxboldls{PlusJakartaSans-ExtraBold.ttf}[Path=../../../fonts/,LetterSpace=3.0]

\setlist[enumerate]{leftmargin=1.7em,itemsep=5pt,topsep=4pt}
\setlist[itemize]{leftmargin=1.7em,itemsep=4pt,topsep=4pt,label={\color{poporange}\textbullet}}
\setlength{\parindent}{0pt}
\setlength{\parskip}{5pt}
\renewcommand{\arraystretch}{1.25}

% pgfplots : style Pop par défaut
\pgfplotsset{
  every axis/.append style={axis line style={popink,line width=1pt},tick style={popink},
    grid style={popline,line width=0.5pt},label style={font=\small},tick label style={font=\footnotesize}},
  popcurve/.style={poporange,line width=1.8pt,no markers,smooth},
}
% Même style disponible dans un \draw TikZ simple (hors pgfplots)
\tikzset{popcurve/.style={poporange,line width=1.8pt}}
\tikzset{popgrid/.style={popline,very thin}}
\colorlet{popcurve}{poporange} % le nom du style est parfois utilisé comme couleur

% ============================================================
% Pill / squiggle
% ============================================================
\newcommand{\poppill}[3]{%
  \tikz[baseline=-0.55ex]\node[draw=popink,line width=1.2pt,rounded corners=2.6pt,%
    fill=#2,inner xsep=6.5pt,inner ysep=3pt]{%
    \popxboldls\fontsize{7}{7}\selectfont\color{#3}\MakeUppercase{#1}};%
}
\newcommand{\popsquiggle}[1]{%
  \tikz[baseline=-0.4ex]{
    \draw[#1,line width=2.6pt,line cap=round]
      plot [smooth] coordinates {(0,0.09) (0.34,0.01) (0.68,0.21) (1.02,0.01) (1.36,0.10)};
  }%
}
% Mot-clé surligné (marqueur orange sous le texte)
\newcommand{\cle}[1]{{\popxbold\color{popdark}#1}}

% ============================================================
% En-tête / pied
% ============================================================
\pagestyle{fancy}
\fancyhf{}
\renewcommand{\headrulewidth}{0pt}
\renewcommand{\footrulewidth}{0pt}
\lhead{\raisebox{-0.15ex}{\poplogo\fontsize{13}{13}\selectfont\color{popink}OnBuch\color{poporange}+}}
\rhead{\poppill{\DOCMATIERE\ \textbullet\ \DOCNIVEAU\ \textbullet\ Leçon \DOCLECON}{white}{popink}}
\lfoot{\popxbold\fontsize{7.6}{7.6}\selectfont\color{popmuted}\MakeUppercase{Cours signé OnBuch+}\ \textbullet\ {\popsemi\DOCTITRE}}
\rfoot{\tikz[baseline=-0.6ex]\node[rounded corners=8.5pt,fill=popink,minimum height=17pt,minimum width=17pt,inner xsep=5pt,inner sep=0]{\popxbold\fontsize{8}{8}\selectfont\color{popcream}\thepage\,/\,\pageref*{LastPage}};}

% ============================================================
% Titres
% ============================================================
\newcommand{\chaptitre}[2]{%
  \begin{tcolorbox}[enhanced,colback=poporange,colframe=popink,boxrule=2.6pt,arc=11pt,
      drop shadow=popink,left=15pt,right=15pt,top=12pt,bottom=11pt,before skip=3pt,after skip=18pt]
    \poppill{#2}{popink}{popcream}\par\vspace{9pt}
    {\raggedright\hyphenpenalty=10000\exhyphenpenalty=10000\popdisplay\fontsize{22}{24}\selectfont\color{popcream}\MakeUppercase{#1}\par}
  \end{tcolorbox}
}
% Section numérotée : pastille orange avec le numéro + titre Archivo
\newcounter{popsec}
\newcommand{\coursec}[1]{%
  \stepcounter{popsec}\setcounter{popsub}{0}%
  \par\vspace{16pt}\noindent
  \begin{minipage}{\linewidth}
  \tikz[baseline=-0.7ex]\node[draw=popink,line width=1.6pt,fill=poporange,rounded corners=4pt,minimum size=22pt,inner sep=0]{\popdisplay\fontsize{12}{12}\selectfont\color{popcream}\thepopsec};%
  \hspace{8pt}{\popdisplay\fontsize{15}{17}\selectfont\color{popink}\MakeUppercase{#1}}\par
  \vspace{4pt}\noindent{\color{poporange}\rule{36pt}{3.6pt}}
  \end{minipage}\par\nopagebreak\vspace{8pt}
}
\newcounter{popsub}
\newcommand{\courssub}[1]{%
  \stepcounter{popsub}%
  \par\vspace{9pt}\noindent{\popxbold\fontsize{11.5}{13}\selectfont\color{popink}{\color{poporange}\thepopsec.\thepopsub}\hspace{6pt}#1}\par\nopagebreak\vspace{4pt}
}

% ============================================================
% Boîtes pédagogiques — même recette : fond blanc, bordure épaisse colorée,
% ombre dure encre, badge plein. Titre optionnel [#1].
% ============================================================
\tcbset{popbase/.style={breakable,enhanced,colback=white,boxrule=2.3pt,arc=9pt,
  shadow={3pt}{-3pt}{0pt}{popink},left=14pt,right=14pt,top=11pt,bottom=11pt,before skip=11pt,after skip=15pt,
  fontupper=\popsemi\fontsize{10.6}{14.5}\selectfont\color{popink},
  fontlower=\popsemi\fontsize{10.6}{14.5}\selectfont\color{popink}}}
% Coche dessinée (le glyphe ✓ n'existe pas dans les polices du document)
\newcommand{\popcheck}[1]{\tikz[baseline=-0.5ex]\draw[#1,line width=1.6pt,line cap=round,line join=round](-0.13,0.0)--(-0.04,-0.09)--(0.13,0.1);}
\providecommand{\coche}{\popcheck{popgreen}}
\providecommand{\resultat}[1]{\par\smallskip\noindent\fcolorbox{popgreen}{popgreenL}{\parbox{\dimexpr\linewidth-2\fboxsep-2\fboxrule\relax}{\textbf{Résultat.} #1}}\par\smallskip}
\newcommand{\popboxhead}[3]{\poppill{#1}{#2}{white}\ifx\relax#3\relax\else\hspace{6pt}{\popxbold\fontsize{10.6}{12}\selectfont\color{popink}#3}\fi\par\vspace{7pt}}

\newtcolorbox{aretenir}[1][]{popbase,colframe=popgreen,before upper={\popboxhead{À retenir}{popgreen}{#1}}}
\newtcolorbox{exemplebox}[1][]{popbase,colframe=poporange,before upper={\popboxhead{Exemple}{poporange}{#1}}}
\newtcolorbox{definition}[1][]{popbase,colframe=popblue,before upper={\popboxhead{Définition}{popblue}{#1}}}
\newtcolorbox{propriete}[1][]{popbase,colframe=poppurple,before upper={\popboxhead{Propriété}{poppurple}{#1}}}
\newtcolorbox{methode}[1][]{popbase,colframe=popgold,before upper={\popboxhead{Méthode}{popgold}{#1}}}
\newtcolorbox{attention}[1][]{popbase,colframe=poppink,before upper={\popboxhead{Attention}{poppink}{#1}}}
\newtcolorbox{experience}[1][]{popbase,colframe=popblue,colback=popblueL,before upper={\popboxhead{Expérience}{popblue}{#1}}}
% « Le savais-tu ? » : carte violette pleine (contexte, culture, Cameroun)
\newtcolorbox{savaistu}[1][]{popbase,colback=poppurple,colframe=popink,
  fontupper=\popsemi\fontsize{10.4}{14.2}\selectfont\color{white},
  before upper={\poppill{Le savais-tu\,?}{white}{poppurple}\ifx\relax#1\relax\else\hspace{6pt}{\popxbold\color{white}#1}\fi\par\vspace{7pt}}}
% Exercice résolu : énoncé en haut, \tcblower puis la solution sur fond crème
\newcounter{popexo}\newcounter{popexr}
\newtcolorbox{exoresolu}[1][]{popbase,colframe=poporange,colbacklower=popcream,
  segmentation style={popink,line width=1.2pt,dashed},
  before upper={\stepcounter{popexr}\popboxhead{Exercice résolu \thepopexr}{poporange}{#1}},
  before lower={\poppill{Solution}{popink}{popcream}\par\vspace{6pt}}}
% Exercice (non résolu en place) — la correction est dans la partie Corrigés
\newtcolorbox{exercice}[1][]{popbase,colframe=popink,
  before upper={\stepcounter{popexo}\popboxhead{Exercice \thepopexo}{popink}{#1}}}
% Corrigé
\newtcolorbox{corrige}[1][]{popbase,colframe=popgreen,colback=popgreenL,
  before upper={\popboxhead{Corrigé}{popgreen}{#1}}}
% Objectifs / prérequis / bilan
\newtcolorbox{objectifs}{popbase,colframe=popink,colback=poporangeL,
  before upper={\popboxhead{Objectifs de la leçon}{popink}{}}}
\newtcolorbox{prerequis}{popbase,colframe=popink,colback=popblueL,
  before upper={\popboxhead{Prérequis}{popblue}{}}}
\newtcolorbox{bilan}{popbase,colframe=popink,colback=white,boxrule=2.8pt,
  before upper={\popboxhead{Fiche bilan}{poporange}{L'essentiel à connaître pour l'examen}}}
% Figure encadrée avec légende
\newtcolorbox{popfigure}[1][]{enhanced,breakable=false,colback=white,colframe=popline,boxrule=1.4pt,arc=8pt,
  left=10pt,right=10pt,top=10pt,bottom=8pt,before skip=10pt,after skip=12pt,halign=center,
  lower separated=false,halign lower=center,
  fontlower=\popsemi\fontsize{9}{11.5}\selectfont\color{popmuted},
  #1}
\newcommand{\legende}[1]{\par\vspace{6pt}{\popsemi\fontsize{9}{11.5}\selectfont\color{popmuted}#1}}

% ============================================================
% Signature OnBuch+
% ============================================================
\newcommand{\poplogoinline}[1]{{\poplogo\fontsize{#1}{#1}\selectfont\color{popink}OnBuch\color{poporange}+}}
\newcommand{\signatureonbuch}{%
  \begin{tcolorbox}[enhanced,colback=popink,colframe=popink,boxrule=2.2pt,arc=10pt,
      drop shadow=poporange,left=14pt,right=14pt,top=10pt,bottom=10pt,before skip=0pt,after skip=0pt]
    \tikz[baseline=-0.7ex]\node[circle,fill=popgreen,draw=popcream,line width=1.3pt,minimum size=22pt,inner sep=0]{\popcheck{white}};%
    \hspace{9pt}\begin{minipage}[c]{0.62\linewidth}
      {\popxbold\fontsize{12}{13}\selectfont\color{popcream}Signé\ \textbullet\ {\poplogo OnBuch\color{poporange}+}}\par\vspace{2pt}
      {\raggedright\popsemi\fontsize{8}{10.5}\selectfont\color{popline}\MakeUppercase{Équipe pédagogique OnBuch+}\\\MakeUppercase{Conforme au programme officiel MINESEC}\par}
    \end{minipage}\hfill
    {\poplogo\fontsize{22}{22}\selectfont\color{popcream}OnBuch\color{poporange}+}
  \end{tcolorbox}%
}

% ============================================================
% Page de garde
% #1 titre · #2 matière · #3 niveau · #4 module · #5 n° leçon · #6 durée ·
% #7 description / objectifs (texte ou itemize)
% ============================================================
\newcommand{\pagegarde}[7]{%
  \thispagestyle{empty}
  \newgeometry{a4paper,margin=1.7cm,top=1.6cm,bottom=1.4cm}%
  \begin{tikzpicture}[remember picture,overlay]
    \fill[poporange!18] ([xshift=-3.2cm,yshift=-3.6cm]current page.north east) circle (4.6cm);
    \fill[poppurple!14] ([xshift=2.4cm,yshift=7.6cm]current page.south west) circle (3.8cm);
  \end{tikzpicture}%
  {\poplogo\fontsize{30}{30}\selectfont\color{popink}OnBuch\color{poporange}+}\hfill
  \raisebox{0.6ex}{\poppill{Cours signé \textbullet\ Premium}{popink}{popcream}}\par
  \vspace{1pt}\popsquiggle{poppurple}\par
  \vspace{8pt}
  \begin{tcolorbox}[enhanced,colback=poporange,colframe=popink,boxrule=2.8pt,arc=13pt,
      drop shadow=popink,left=18pt,right=18pt,top=12pt,bottom=13pt,after skip=10pt]
    \poppill{#2 \textbullet\ #3}{popink}{popcream}\hspace{5pt}\poppill{#4}{white}{popink}\par\vspace{8pt}
    {\popdisplay\fontsize{14}{14}\selectfont\color{popink}LEÇON #5}\par\vspace{4pt}
    {\raggedright\hyphenpenalty=10000\exhyphenpenalty=10000\popdisplay\fontsize{25}{27}\selectfont\color{popcream}\MakeUppercase{#1}\par}
    \vspace{8pt}
    {\popxbold\fontsize{9.5}{9.5}\selectfont\color{popink}\MakeUppercase{Durée conseillée : #6}}
  \end{tcolorbox}
  \begin{tcolorbox}[enhanced,colback=white,colframe=popink,boxrule=2.2pt,arc=10pt,
      drop shadow=popink,left=16pt,right=16pt,top=10pt,bottom=10pt,after skip=8pt]
    \poppill{Dans cette leçon}{poppurple}{white}\par\vspace{5pt}
    \popsemi\fontsize{9.3}{12.6}\selectfont\color{popink}#7
  \end{tcolorbox}
  \noindent\begin{minipage}[t]{0.487\linewidth}
    \begin{tcolorbox}[enhanced,colback=popgreen,colframe=popink,boxrule=2.2pt,arc=10pt,
        drop shadow=popink,left=11pt,right=11pt,top=10pt,bottom=10pt,height=3.1cm,valign=center]
      \tcbox[colback=white,colframe=popink,boxrule=1.3pt,arc=5pt,boxsep=0pt,nobeforeafter,
        left=3pt,right=3pt,top=3pt,bottom=3pt,valign=center]{\qrcode[height=1.9cm]{https://play.google.com/store/apps/details?id=com.luvvix.onbuch&hl=fr}}%
      \hspace{8pt}\begin{minipage}[c]{0.5\linewidth}
        {\popdisplay\fontsize{11}{12.5}\selectfont\color{white}\MakeUppercase{Télécharge OnBuch}}\par\vspace{4pt}
        {\raggedright\popsemi\fontsize{8.4}{11}\selectfont\color{white}Gratuit \textbullet\ Google Play\\Tuteur IA, annales, concours, résultats\par}
      \end{minipage}
    \end{tcolorbox}
  \end{minipage}\hfill
  \begin{minipage}[t]{0.487\linewidth}
    \begin{tcolorbox}[enhanced,colback=poppurple,colframe=popink,boxrule=2.2pt,arc=10pt,
        drop shadow=popink,left=11pt,right=11pt,top=10pt,bottom=10pt,height=3.1cm,valign=center]
      \tikz[baseline=-0.7ex]\node[circle,fill=white,draw=popink,line width=1.3pt,minimum size=1.6cm,inner sep=0]{\popdisplay\fontsize{24}{24}\selectfont\color{poppurple}+};%
      \hspace{8pt}\begin{minipage}[c]{0.6\linewidth}
        {\popdisplay\fontsize{11}{12.5}\selectfont\color{white}\MakeUppercase{Passe à OnBuch+}}\par\vspace{4pt}
        {\raggedright\popsemi\fontsize{8.4}{11}\selectfont\color{white}Tous les cours signés, Tuteur IA illimité, fiches PDF — onglet OnBuch+ de l'app\par}
      \end{minipage}
    \end{tcolorbox}
  \end{minipage}
  \vspace{8pt}
  \popcontenu
  \vfill
  \signatureonbuch
  \restoregeometry
  \newpage
}

% Rangée « Ce cours contient » (page de garde)
\newcommand{\poptile}[3]{% #1 couleur #2 gros texte #3 légende
  \begin{tcolorbox}[enhanced,colback=#1,colframe=popink,boxrule=2pt,arc=9pt,shadow={2.5pt}{-2.5pt}{0pt}{popink},
      left=6pt,right=6pt,top=9pt,bottom=9pt,halign=center,valign=center,height=1.75cm,width=\linewidth,nobeforeafter]
    {\popdisplay\fontsize{12.5}{13}\selectfont\color{white}\MakeUppercase{#2}}\par\vspace{3pt}
    {\centering\popsemi\fontsize{7.8}{9.5}\selectfont\color{white}#3\par}
  \end{tcolorbox}}
\newcommand{\popcontenu}{%
  \par\noindent{\popxboldls\fontsize{7.5}{7.5}\selectfont\color{popmuted}CE COURS SIGNÉ CONTIENT}\par\vspace{7pt}
  \noindent\begin{minipage}[t]{0.235\linewidth}\poptile{popblue}{Cours}{complet, illustré, pas à pas}\end{minipage}\hfill
  \begin{minipage}[t]{0.235\linewidth}\poptile{poporange}{Méthodes}{et exercices résolus}\end{minipage}\hfill
  \begin{minipage}[t]{0.235\linewidth}\poptile{poppink}{Exercices}{type examen + corrigés}\end{minipage}\hfill
  \begin{minipage}[t]{0.235\linewidth}\poptile{popgreen}{Fiche bilan}{l'essentiel à retenir}\end{minipage}\par}

% Sommaire
\newtcolorbox{sommaire}{enhanced,colback=white,colframe=popink,boxrule=2.2pt,arc=10pt,
  drop shadow=popink,left=18pt,right=18pt,top=13pt,bottom=13pt,before skip=6pt,after skip=20pt,
  before upper={\poppill{Sommaire}{poppurple}{white}\par\vspace{9pt}\popsemi\fontsize{10.6}{14.5}\selectfont\color{popink}}}

% Signature de fin de cours — poussée tout en bas de la dernière page
\newcommand{\findecours}{%
  \par\vfill\nopagebreak
  \begin{center}\popsquiggle{poporange}\end{center}\vspace{4pt}
  \signatureonbuch
}
\AtBeginDocument{\def\llbracket{\mathopen{[\mkern-3.5mu[}}\def\rrbracket{\mathclose{]\mkern-3.5mu]}}} % intervalles d'entiers (glyphe ⟦ absent de la police)
% Glyphes absents de Fira Math : redéfinis à partir de glyphes disponibles
\AtBeginDocument{%
  \def\setminus{\mathbin{\backslash}}\let\smallsetminus\setminus
  \def\vdots{\mathord{\vbox{\baselineskip3.2pt\lineskiplimit0pt\kern4pt\hbox{.}\hbox{.}\hbox{.}}}}%
  \def\ddots{\mathinner{\mkern1mu\raise6.5pt\hbox{.}\mkern2mu\raise3.6pt\hbox{.}\mkern2mu\raise0.7pt\hbox{.}\mkern1mu}}%
  \def\triangle{\mathord{\scalebox{0.9}{$\symup{\Delta}$}}}%
  \def\nabla{\mathord{\raisebox{\depth}{\scalebox{1}[-1]{$\symup{\Delta}$}}}}%
  \def\checkmark{\popcheck{popdark}}%
  \def\star{\mathbin{\ast}}\def\bigstar{\mathord{\ast}}%
  \def\diamond{\mathbin{\scalebox{0.6}{$\Diamond$}}}%
  \def\bigcup{\mathop{\mathchoice{\raisebox{-0.3em}{\scalebox{1.7}{$\cup$}}}{\scalebox{1.25}{$\cup$}}{\cup}{\cup}}\displaylimits}%
  \def\bigcap{\mathop{\mathchoice{\raisebox{-0.3em}{\scalebox{1.7}{$\cap$}}}{\scalebox{1.25}{$\cap$}}{\cap}{\cap}}\displaylimits}%
  \def\lesseqqgtr{\mathrel{\lessgtr}}\def\bigoplus{\mathop{\oplus}\displaylimits}%
}
\providecommand{\ssi}{si et seulement si} % abréviation fréquente
\AtBeginDocument{\providecommand{\wideparen}{\overparen}} % arc de cercle

%% FIGFIX-BEGIN v5 — correctifs globaux des figures (docs/onbuch-generation/tools/figfix)
\usepackage{adjustbox}\usepackage{etoolbox}\tcbuselibrary{raster}
% toute figure TikZ/pgfplots plus large que la ligne est réduite à la largeur disponible
\BeforeBeginEnvironment{tikzpicture}{\begin{adjustbox}{max width=\linewidth}}
\AfterEndEnvironment{tikzpicture}{\end{adjustbox}}
% légendes pgfplots lisibles par-dessus les courbes
\pgfplotsset{every axis/.append style={legend style={fill=white,fill opacity=0.92,text opacity=1,draw=black!15,inner sep=3pt}}}
% étiquettes d'arêtes (syntaxe edge["..."]) avec fond blanc
\tikzset{every edge quotes/.append style={fill=white,inner sep=1.2pt}}
%% FIGFIX-END
\begin{document}

\pagegarde{VI : produire un rapport}{Français}{1re Tech}{Français – classe de Première technique}{N22}{10 séances (fiche de progression harmonisée)}{Cette leçon apprend à l'élève de Première technique à produire un rapport complet : maîtrise de la structure externe et interne, rédaction de l'introduction, du corps et de la conclusion, emploi du texte explicatif et distinction entre lexique commun et lexique spécialisé. Elle s'appuie sur des situations concrètes camerounaises et sur la lecture de textes fonctionnels et d'affiches de sécurité.
\vspace{4pt}
\begin{multicols}{2}\small
\begin{itemize}
\item À la fin de cette leçon, je sais définir le rapport et identifier ses caractéristiques (objectivité, exactitude, clarté)
\item À la fin de cette leçon, je sais identifier la structure externe d'un rapport (en-tête, références, objet, signature)
\item À la fin de cette leçon, je sais organiser la structure interne d'un rapport (introduction, corps, conclusion)
\item À la fin de cette leçon, je sais rédiger l'introduction et la conclusion d'un rapport
\item À la fin de cette leçon, je sais rédiger le corps d'un rapport en utilisant les procédés du texte explicatif
\item À la fin de cette leçon, je sais distinguer le lexique commun du lexique spécialisé et les employer à bon escient
\end{itemize}\end{multicols}}

\begin{sommaire}
\begin{multicols}{2}
\begin{enumerate}
\item Découverte du rapport : définition et caractéristiques
\item La structure externe du rapport
\item La structure interne : introduction, corps, conclusion
\item Le texte explicatif au service du rapport
\item Lexique commun et lexique spécialisé
\item Lecture de textes fonctionnels et affiches de sécurité
\item Production et amélioration d'un rapport complet
\item Activité d'intégration
\item Méthodes et exercices résolus
\item Exercices
\item Corrigés des exercices
\item Fiche bilan
\end{enumerate}
\end{multicols}
\end{sommaire}

\begin{objectifs}
\begin{itemize}
\item À la fin de cette leçon, je sais définir le rapport et identifier ses caractéristiques (objectivité, exactitude, clarté)
\item À la fin de cette leçon, je sais identifier la structure externe d'un rapport (en-tête, références, objet, signature)
\item À la fin de cette leçon, je sais organiser la structure interne d'un rapport (introduction, corps, conclusion)
\item À la fin de cette leçon, je sais rédiger l'introduction et la conclusion d'un rapport
\item À la fin de cette leçon, je sais rédiger le corps d'un rapport en utilisant les procédés du texte explicatif
\item À la fin de cette leçon, je sais distinguer le lexique commun du lexique spécialisé et les employer à bon escient
\item À la fin de cette leçon, je sais lire et interpréter des textes fonctionnels et des affiches d'indications sécuritaires
\item À la fin de cette leçon, je sais produire et améliorer un rapport complet à partir d'une situation concrète
\item À la fin de cette leçon, je sais justifier ma démarche rédactionnelle et corriger un rapport défectueux
\end{itemize}
\end{objectifs}

\begin{prerequis}
\begin{itemize}
\item Les types de textes (narratif, descriptif, informatif) vus en classe de Seconde
\item Le compte rendu : rapporter fidèlement des faits ou des paroles
\item Les connecteurs logiques et l'organisation du paragraphe
\item La correspondance administrative (lettre officielle) vue en début d'année
\item Le discours direct et le discours rapporté
\end{itemize}
\end{prerequis}

\newpage

\coursec{Découverte du rapport : définition et caractéristiques}

\courssub{Situation d'accroche et problématique}

Imaginez la scène : nous sommes au lycée technique de Nkongsamba. Il est 14 h 30, l'odeur de sciure et d'huile de coupe flotte dans l'atelier de menuiserie. Soudain, une étincelle jaillit d'une scie circulaire mal entretenue, enflamme un tas de copeaux secs. En quelques minutes, la panique s'installe. L'enseignant d'atelier maîtrise le départ de feu grâce à l'extincteur \textbf{CO\textsubscript{2}} situé près de la sortie, mais les dégâts sont réels : une machine hors service, des panneaux de contreplaqué détruits, un élève légèrement brûlé au bras évacué à l'infirmerie.

Le proviseur, M. Fouda, appelle le délégué de la classe de Première technique, Aminatou. « Aminatou, j'ai besoin d'un rapport sur mon bureau avant 17 h pour le transmettre au Délégué régional de l'Éducation. Pas d'émotion, que des faits. »

\begin{attention}[Erreur fréquente : confondre rapport et lettre de plainte]
Beaucoup d'élèves, face à cet événement, voudraient écrire : « Monsieur le Proviseur, je viens vous signaler que l'atelier est dangereux, les machines sont vieilles, c'est scandaleux ! » \textbf{Non.} C'est une lettre de plainte ou une réclamation. Le \cle{rapport} expose des \textbf{faits bruts, vérifiables, chronologiques}. Il ne juge pas, ne réclame pas, n'exprime pas de colère. Il \emph{informe} pour permettre une décision.
\end{attention}

\noindent \textbf{Problématique :} Comment présenter les faits de manière claire, objective et structurée ? Quelles informations mettre dans l'introduction, le corps et la conclusion ? Quel vocabulaire technique employer pour décrire les dégâts ? Cette leçon outille l'élève pour produire un rapport professionnel, compétence indispensable dans la vie scolaire, administrative et professionnelle au Cameroun.

\courssub{Définition et nature du rapport}

\begin{definition}[Le rapport (texte fonctionnel)]
Un \cle{rapport} est un texte fonctionnel à visée informative qui rend compte, de manière \textbf{objective}, \textbf{exacte} et \textbf{structurée}, de faits, d'observations, d'une mission ou d'une situation, à l'intention d'un \textbf{destinataire déterminé} (hiérarchie, client, administration) pour lui permettre de prendre une décision ou de conserver une trace administrative.
\end{definition}

\noindent Contrairement au récit littéraire (qui cherche l'émotion, l'esthétique, le « je » subjectif) ou à l'article de presse (qui vise le grand public, l'actualité chaude), le rapport est un \textbf{outil de gestion et de communication professionnelle}. Il répond à une commande (explicite ou implicite) et suit des normes codifiées.

\begin{aretenir}[Les quatre piliers du rapport]
\begin{enumerate}
    \item \textbf{Objectivité :} Pas de « je pense que », « c'est scandaleux », « heureusement ». Seuls les faits observés ou mesurés.
    \item \textbf{Exactitude :} Chiffres précis (pas « environ 10 mètres » mais « 12,50 m »), noms exacts des machines, références des normes (ex. : NF EN 3-7 pour extincteurs).
    \item \textbf{Clarté :} Phrases courtes, voix passive ou impersonnelle privilégiée, connecteurs logiques (donc, par conséquent, ensuite).
    \item \textbf{Concision :} Aucune phrase superflue. Chaque mot porte une information utile au destinataire.
\end{enumerate}
\end{aretenir}

\courssub{Distinction : rapport, compte rendu, procès-verbal}

Ces trois genres sont souvent confondus car ils partagent l'objectivité. Leur distinction repose sur le \textbf{destinataire}, l'\textbf{objet} et la \textbf{valeur juridique}.

\begin{popfigure}
\centering
\begin{tikzpicture}[
    node distance=2mm and 4mm,
    header/.style={fill=popblue, text=white, font=\bfseries\small, minimum height=8mm, align=center},
    cell/.style={font=\small, minimum height=7mm, align=center, draw=popline},
    titlecell/.style={fill=popblueL, font=\bfseries\small, minimum height=7mm, align=center, draw=popline},
    table/.style={draw=popblue, thick, rounded corners=2pt}
]
% En-têtes
\node[header, minimum width=3.2cm] (c1) {Critère de distinction};
\node[header, minimum width=3.5cm, right=of c1] (c2) {\textbf{Rapport}};
\node[header, minimum width=3.5cm, right=of c2] (c3) {\textbf{Compte rendu}};
\node[header, minimum width=3.5cm, right=of c3] (c4) {\textbf{Procès-verbal (PV)}};
% Lignes
% Destinataire
\node[titlecell, minimum width=3.2cm, below=of c1] (r1c1) {Destinataire};
\node[cell, minimum width=3.5cm, below=of c2] (r1c2) {Hiérarchie supérieure, commanditaire, client (ex. : Délégué régional, Proviseur)};
\node[cell, minimum width=3.5cm, below=of c3] (r1c3) {Participants, membres d'un groupe, service concerné (ex. : collègues de stage, association)};
\node[cell, minimum width=3.5cm, below=of c4] (r1c4) {Autorité judiciaire, administrative, tribunal, assemblée générale (valeur légale)};
% Objet
\node[titlecell, minimum width=3.2cm, below=of r1c1] (r2c1) {Objet principal};
\node[cell, minimum width=3.5cm, below=of r1c2] (r2c2) {Mission, incident, visite, expertise, activité technique ; \emph{analyse + propositions}};
\node[cell, minimum width=3.5cm, below=of r1c3] (r2c3) {Réunion, séance de travail, conférence, soutenance ; \emph{fidélité aux échanges}};
\node[cell, minimum width=3.5cm, below=of r1c4] (r2c4) {Constat d'infraction, élection, assemblée, réception de travaux ; \emph{constat juridique}};
% Structure
\node[titlecell, minimum width=3.2cm, below=of r2c1] (r3c1) {Structure type};
\node[cell, minimum width=3.5cm, below=of r2c2] (r3c2) {En-tête, Intro (contexte), Corps (faits/analyse), Conclusion (bilan/recommandations), Annexes};
\node[cell, minimum width=3.5cm, below=of r2c3] (r3c3) {En-tête, Ordre du jour, Déroulement chronologique (débats/votes), Décisions, Signatures};
\node[cell, minimum width=3.5cm, below=of r2c4] (r3c4) {En-tête légal, Constats techniques signés, Qualifications juridiques, Signatures des officiers/constatateurs};
% Valeur juridique
\node[titlecell, minimum width=3.2cm, below=of r3c1] (r4c1) {Valeur juridique};
\node[cell, minimum width=3.5cm, below=of r3c2] (r4c2) {Preuve administrative, base de décision, non opposable en justice seul};
\node[cell, minimum width=3.5cm, below=of r3c3] (r4c3) {Preuve morale des échanges, engagement des participants};
\node[cell, minimum width=3.5cm, below=of r3c4] (r4c4) {\textbf{Force probante forte} (jusqu'à preuve contraire), acte authentique si officier public};
% Exemple contexte Tech
\node[titlecell, minimum width=3.2cm, below=of r4c1] (r5c1) {Exemple en 1\textsuperscript{re} Tech};
\node[cell, minimum width=3.5cm, below=of r4c2] (r5c2) {Rapport de stage à la SONARA ; Rapport d'incident atelier};
\node[cell, minimum width=3.5cm, below=of r4c3] (r5c3) {Compte rendu de la réunion du bureau du CVL ; CR de TP};
\node[cell, minimum width=3.5cm, below=of r4c4] (r5c4) {PV de réception des nouveaux tours à commande numérique ; PV d'élection délégués};
% Cadre global
\draw[table] (c1.north west) rectangle (r5c4.south east);
% Lignes horizontales
\foreach \n in {c1,r1c1,r2c1,r3c1,r4c1} {
    \draw[popline] (\n.south west) -- (\n.south west -| r5c4.south east);
}
% Lignes verticales
\draw[popline] (c1.north east) -- (r5c1.south east);
\draw[popline] (c2.north east) -- (r5c2.south east);
\draw[popline] (c3.north east) -- (r5c3.south east);
\end{tikzpicture}
\legende{Tableau comparatif : Rapport, Compte rendu et Procès-verbal. Le rapport se distingue par sa visée décisionnelle pour une hiérarchie et sa structure analytique (analyse + recommandations).}
\end{popfigure}

\begin{methode}[Comment choisir le bon genre ?]
\begin{enumerate}
    \item \textbf{Identifiez le destinataire :} Hiérarchie/commanditaire $\rightarrow$ Rapport. Pair/collectif $\rightarrow$ Compte rendu. Justice/Administration formelle $\rightarrow$ PV.
    \item \textbf{Identifiez la finalité :} Aider à décider (analyse + propositions) $\rightarrow$ Rapport. Garder trace des échanges $\rightarrow$ Compte rendu. Constater un fait juridique $\rightarrow$ PV.
    \item \textbf{Vérifiez la signature :} Rapport = signataire unique (auteur). PV = signatures multiples (constatateurs, parties).
\end{enumerate}
\end{methode}

\courssub{Les types de rapports en milieu technique et professionnel}

Le programme de Première technique cible particulièrement les rapports liés à la formation professionnelle. En voici la typologie principale :

\begin{center}
\begin{tabularx}{\linewidth}{|l|X|X|}\hline
\rowcolor{popblue}\textbf{Type de rapport} & \textbf{Contexte déclencheur (Exemple Cameroun)} & \textbf{Spécificité rédactionnelle} \\\hline
\textbf{Rapport d'incident / d'accident} & Incendie atelier menuiserie (Nkongsamba), chute dans escalier, panne critique machine CNC. & \emph{Urgence, chronologie stricte (heure par heure), causes probables, mesures conservatoires immédiates, dommages chiffrés.} \\\hline
\textbf{Rapport de stage} & Stage de 4 semaines en 1\textsuperscript{re} Tech à la SODECOTON (Maroua) ou SONARA (Limbe). & \emph{Description structure entreprise, tâches exécutées, compétences acquises, difficultés techniques, suggestions pédagogiques.} \\\hline
\textbf{Rapport de visite d'entreprise} & Visite pédagogique d'une journée à la brasserie UCAM (Douala) ou centrale ENEO. & \emph{Focus sur l'observation : process industriels, sécurité, maintenance, organisation. Moins d'analyse personnelle que le stage.} \\\hline
\textbf{Rapport d'activité (périodique)} & Bilan mensuel du club robotique, rapport trimestriel du délégué de classe au proviseur. & \emph{Récurrence, indicateurs quantitatifs (effectifs, budget, taux de réalisation), comparatif période précédente.} \\\hline
\textbf{Rapport d'expertise / technique} & Demande d'avis sur la conformité d'une installation électrique (norme NF C 15-100). & \emph{Très normé : références normes, mesures instrumentées, verdict technique (conforme/non conforme), prescriptions.} \\\hline
\end{tabularx}
\end{center}

\begin{savaistu}[Le rapport dans la fonction publique camerounaise]
Au Cameroun, le \cle{rapport} est l'outil privilégié de \textbf{remontée de l'information} vers la hiérarchie. Un chef de service (ex. : Chef de travaux au Lycée technique) rédige un « rapport d'activité » mensuel pour le Proviseur ; le Proviseur synthétise dans un « rapport de synthèse » pour le Délégué Régional ; ce dernier envoie un « rapport statistique » au MINESUP/MINESEC. La carrière d'un fonctionnaire technique dépend souvent de la qualité de ses rapports : clarté, respect de la chaîne hiérarchique, vocabulaire administratif approprié (\"Vu\", \"Attendu que\", \"Je soussigné\", \"Fait à..., le...\").
\end{savaistu}

\courssub{L'influence du destinataire sur le ton et le niveau de langue}

Le rapport n'est pas un texte « neutre » au sens où il s'adresserait à tout le monde. Il est \textbf{sur mesure} pour son destinataire.

\begin{exemplebox}[Adaptation au destinataire : L'incendie de l'atelier menuiserie]
\textbf{Scénario :} Même fait (incendie), trois destinataires différents.
\begin{itemize}
    \item \textbf{Au Proviseur (hiérarchie directe, technique) :} Ton technique, vocabulaire métier (\"scie circulaire 380V triphasé\", \"extincteur CO\textsubscript{2} 5 kg\", \"contreplaqué Okoumé 18 mm\"), analyse des causes techniques (\"absence de capteur de surchauffe\"), propositions concrètes (\"remplacement disjoncteur différentiel 30 mA\").
    \item \textbf{Au Délégué Régional (administration, gestion budgétaire) :} Ton administratif, synthèse financière (\"coût estimé des dégâts : 1 250 000 FCFA\"), impact pédagogique (\"arrêt module usinage 3 semaines\"), demande d'arbitrage (\"demande de dotation exceptionnelle\").
    \item \textbf{À l'Assureur (expertise, indemnisation) :} Ton expert, descriptif précis des biens détruits (numéros de série, factures d'achat, valeur à neuf), circonstances exactes, photos jointes en annexe, réserve sur causes (\"enquête en cours\").
\end{itemize}
\end{exemplebox}

\begin{attention}[Piège : le tutoiement ou le langage familier]
Même si le destinataire est votre tuteur de stage que vous tutoyez à la pause café, le rapport s'adresse à \textbf{la fonction}, pas à la personne. On utilise le \textbf{vouvoiement de fonction} ou l'impersonnel. \textbf{Interdit :} « T'as vu la machine ? Elle a pété un câble. » \textbf{Correct :} « L'examen de la machine a révélé une rupture de l'isolant du câble d'alimentation. »
\end{attention}

\coursec{La structure externe du rapport}

\courssub{Les éléments obligatoires de l'en-tête}

La structure externe (la « mise en page ») conditionne la recevabilité administrative du document. Elle suit un formalisme strict.

\begin{definition}[Structure externe type (norme administrative camerounaise)]
\begin{enumerate}
    \item \textbf{Référence :} Numéro d'ordre / service / année (ex. : N\textsuperscript{o} 2024/0145/LTN-KNG).
    \item \textbf{Objet :} Une phrase nominale précise (ex. : « Départ de feu atelier menuiserie — Machine SC-400 — Blessé léger »).
    \item \textbf{Expéditeur (De) :} Nom, prénom, qualité, service (ex. : Aminatou M., Déléguée de classe 1\textsuperscript{re} TMA).
    \item \textbf{Destinataire (À) :} Titre, nom, service (ex. : M. le Proviseur, Lycée Technique de Nkongsamba).
    \item \textbf{Date et Heure :} Jour, mois, année ; heure de rédaction (ex. : 15 octobre 2024 — 16 h 30).
    \item \textbf{Distribution (optionnel) :} Copie à... (ex. : C/c : Délégué Régional, Chef de travaux).
    \item \textbf{Pièces jointes :} Liste numérotée en bas de page (ex. : 1. Photos (3) ; 2. Devis réparation ; 3. Fiche infirmerie).
\end{enumerate}
\end{definition}

\begin{popfigure}
\centering
\begin{tikzpicture}[
    scale=0.9,
    every node/.style={transform shape},
    box/.style={draw=popdark, thick, rounded corners=3pt, fill=popcream, inner sep=6pt, text width=13cm, align=left, font=\small},
    arr/.style={->, thick, poporange, shorten >=2pt, shorten <=2pt},
    lbl/.style={font=\tiny\bfseries, text=popdark, fill=poporangeL, rounded corners=1pt, inner sep=2pt},
]
% Structure schematic
\node[box, align=center] (header){
\textbf{RAPPORT D'INCIDENT N\textsuperscript{o} 2024/0145/LTN-KNG} \\[1mm]
\textbf{Objet :} Départ de feu atelier menuiserie — Machine SC-400 — Blessé léger \\[1mm]
\textbf{De :} Aminatou M., Déléguée de classe 1\textsuperscript{re} TMA \\
\textbf{À :} M. le Proviseur, Lycée Technique de Nkongsamba \\
\textbf{Date :} 15 octobre 2024 \hfill \textbf{Heure :} 16 h 30 \\
\textbf{C/c :} Délégué Régional MINESEC Littoral ; Chef de travaux \\
\textbf{P.J. :} 1. Photos (3) \quad 2. Devis réparation \quad 3. Fiche intervention infirmerie
};
% Annotations
\node[lbl, align=center, text width=9cm, above=0.8cm of header.north] {\textbf{En-tête normalisé}\\(Référence, Objet, Expéditeur,\\Destinataire, Date, Heure, C/c, P.J.)\\(Garantit \textbf{traçabilité} et \textbf{archivage})};
\node[lbl, align=center, text width=9cm, below=0.8cm of header.south] {\textbf{Objet = Résumé factuel}\\(Mots-clés : Nature, Lieu, Machine,\\Conséquence humaine)\\(Pas de phrase verbe, pas d'avis)};
\end{tikzpicture}
\legende{Schéma de la structure externe (en-tête) d'un rapport administratif type. L'objet doit permettre le classement et la recherche rapide.}
\end{popfigure}

\coursec{La structure interne : introduction, corps, conclusion}

\courssub{Architecture en trois temps}

La structure interne suit une logique \textbf{ententonnoir} : du contexte global (introduction) aux faits détaillés (corps) vers la décision future (conclusion).

\begin{popfigure}
\centering
\begin{tikzpicture}[
    node distance=8mm and 4mm,
    main/.style={draw=popblue, thick, fill=popblueL, rounded corners=4pt, minimum width=3.5cm, minimum height=2.2cm, align=center, font=\small},
    sub/.style={draw=poporange, fill=poporangeL, rounded corners=2pt, minimum width=3.2cm, minimum height=1.2cm, align=center, font=\tiny},
    arr/.style={->, thick, popdark},
]
\node[main, align=center] (intro){\textbf{1. INTRODUCTION}\\(Contexte \& Mandat)\\\textit{Qui ? Quoi ? Où ? Quand ? Pourquoi ?}};
\node[main, right=of intro, align=center] (corps){\textbf{2. CORPS DU RAPPORT}\\(Faits, Analyse, Données)\\\textit{Chronologie OU Thématique}};
\node[main, right=of corps, align=center] (concl){\textbf{3. CONCLUSION}\\(Bilan \& Recommandations)\\\textit{Synthèse + Actions concrètes}};
\draw[arr] (intro.east) -- (corps.west) node[midway, above, font=\tiny, text=popdark] {Contexte $\rightarrow$ Preuves};
\draw[arr] (corps.east) -- (concl.west) node[midway, above, font=\tiny, text=popdark] {Analyse $\rightarrow$ Décision};
% Sub-nodes Intro
\node[sub, below=of intro, xshift=-1.8cm] (i1) {Identité auteur / commande};
\node[sub, below=of intro, xshift=1.8cm] (i2) {Cadre spatio-temporel précis};
\draw[arr, popmuted] (intro.south) -- (i1.north);
\draw[arr, popmuted] (intro.south) -- (i2.north);
% Sub-nodes Corps
\node[sub, below=of corps, xshift=-2.2cm, align=center] (c1){Chronologie des faits\\ (Heures, actions, acteurs)};
\node[sub, below=of corps, xshift=0cm, align=center] (c2){Bilan chiffré\\ (Humain, Matériel, Financier)};
\node[sub, below=of corps, xshift=2.2cm, align=center] (c3){Analyse causes techniques\\ (Méthode 5M, Arbre des causes)};
\draw[arr, popmuted] (corps.south) -- (c1.north);
\draw[arr, popmuted] (corps.south) -- (c2.north);
\draw[arr, popmuted] (corps.south) -- (c3.north);
% Sub-nodes Conclusion
\node[sub, below=of concl, xshift=-1.8cm] (f1) {Bilan global (Maîtrisé / Critique)};
\node[sub, below=of concl, xshift=1.8cm, align=center] (f2){Recommandations\\ (Action, Responsable, Délai, Coût)};
\draw[arr, popmuted] (concl.south) -- (f1.north);
\draw[arr, popmuted] (concl.south) -- (f2.north);
\end{tikzpicture}
\legende{Architecture interne type du rapport technique. Le corps peut être organisé chronologiquement (incident) ou thématiquement (stage, visite).}
\end{popfigure}

\courssub{Rédaction de l'introduction : le « contrat de lecture »}

L'introduction ne raconte pas l'histoire, elle \textbf{pose le cadre}. Elle répond aux 5 questions fondamentales (QQOQCP) sans développer.

\begin{methode}[Rédiger une introduction efficace (3 phrases types)]
\begin{enumerate}
    \item \textbf{La phrase de mandat :} « \textit{Suivant vos instructions verbales du 15 octobre 2024 / Dans le cadre de mon stage du 3 au 28 juin 2024 / En ma qualité de délégué de classe, je vous présente le rapport relatif à...} »
    \item \textbf{La phrase de fait déclencheur :} « \textit{Le 15 octobre 2024 à 14 h 28, un départ de feu s'est déclaré sur la scie circulaire n\textsuperscript{o} SC-400 de l'atelier menuiserie (Bâtiment B, RDC).} »
    \item \textbf{La phrase d'annonce du plan :} « \textit{Ce rapport expose la chronologie des faits, dresse le bilan des dommages, analyse les causes probables et formule des recommandations immédiates.} »
\end{enumerate}
\end{methode}

\courssub{Rédaction du corps du rapport : faits, analyse, données}

Le corps est la « preuve ». Deux organisations possibles selon le type de rapport :

\begin{propriete}[Choix de l'organisation du corps]
\begin{itemize}
    \item \textbf{Organisation CHRONOLOGIQUE} (Rapports d'incident, d'accident, d'intervention) : Minute par minute, phase par phase. Verbes d'action au passé composé. \cle{Mots-clés} : \textit{À 14 h 28... ; Ensuite... ; Par la suite... ; Finalement...}
    \item \textbf{Organisation THÉMATIQUE} (Rapports de stage, de visite, d'activité, d'expertise) : Par pôles (Ex. : 1. Présentation structure ; 2. Activités menées ; 3. Compétences ; 4. Difficultés ; 5. Suggestions). \cle{Mots-clés} : \textit{En ce qui concerne... ; S'agissant de... ; Par ailleurs... ; En matière de...}
\end{itemize}
\end{propriete}

\begin{aretenir}[Règles d'or du corps de rapport]
\begin{itemize}
    \item \textbf{Un paragraphe = une idée / un fait / un thème.}
    \item \textbf{Chiffrer systématiquement :} « 1 blessé léger », « 420 000 FCFA », « 1,2 m de câble », « 92 dB(A) », « 15 L de copeaux ».
    \item \textbf{Séparer FAITS (observés) et ANALYSE (interprétée).} Utiliser des sous-titres clairs : « 2.1 Chronologie », « 2.2 Bilan », « 2.3 Analyse des causes ».
    \item \textbf{Citer les sources :} « selon le registre de maintenance », « d'après le devis ETS MOTOR REPAIR », « conformément à la norme NF EN 3-7 ».
\end{itemize}
\end{aretenir}

\courssub{Rédaction de la conclusion : bilan et recommandations opérationnelles}

La conclusion n'est pas une « fin de dissertation ». C'est un \textbf{outil d'aide à la décision}.

\begin{methode}[Structurer la conclusion en 2 blocs]
\begin{enumerate}
    \item \textbf{Bilan synthétique (1 phrase) :} « L'incident est maîtrisé sans interruption durable de la formation ; la cause racine est un défaut de maintenance préventive couplé à une hygiène de poste insuffisante. »
    \item \textbf{Recommandations numérotées, actionnables, datées, chiffrées (SMART) :}
    \begin{itemize}
        \item \textbf{Action :} Verbe à l'infinitif ou impératif (Mettre, Nettoyer, Réviser, Demander, Installer).
        \item \textbf{Responsable :} Qui fait ? (Chef de travaux, Entreprise X, Élèves).
        \item \textbf{Délai :} « ce jour avant 12 h », « avant le 30/11/2024 », « dès la rentrée ».
        \item \textbf{Coût / Ressource :} « budget maintenance », « dotation exceptionnelle 1 250 000 FCFA ».
    \end{itemize}
\end{enumerate}
\end{methode}

\begin{exemplebox}[Recommandations SMART (Extrait rapport Nkongsamba)]
\begin{enumerate}
    \item \textbf{Mise hors service} machine SC-400 (consignation LOTO) \textbf{jusqu'à réparation validée} — \textit{Responsable : Chef d'atelier}.
    \item \textbf{Nettoyage intégral} atelier (aspiration copeaux) \textbf{ce jour 16/10 avant 12 h} — \textit{Responsable : Élèves 1\textsuperscript{re} TMA / Agent entretien}.
    \item \textbf{Révision planning} maintenance préventive toutes machines (cible : mensuelle) \textbf{avant le 30/11/2024} — \textit{Responsable : Chef de travaux}.
    \item \textbf{Demande dotation budgétaire} pour 5 extincteurs CO\textsubscript{2} supplémentaires (norme 1/150 m\textsuperscript{2}) \textbf{intégrée budget 2025} — \textit{Responsable : Proviseur}.
\end{enumerate}
\end{exemplebox}

\coursec{Le texte explicatif au service du rapport}

\courssub{Caractéristiques linguistiques du texte explicatif}

Le rapport relève du \textbf{texte explicatif} (sous-type du texte informatif). Son but : \emph{faire comprendre un phénomène, un processus, une situation} par l'analyse causale et la structuration logique.

\begin{definition}[Texte explicatif (Stylistique/Rhétorique)]
Le texte explicatif répond à la question « \textbf{Pourquoi ?} » (causes) ou « \textbf{Comment ?} » (mécanismes, procédures). Il mobilise :
\begin{itemize}
    \item \textbf{La définition, la classification, l'exemplification} pour clarifier les concepts techniques.
    \item \textbf{L'analyse causale} (cause $\rightarrow$ effet) pour expliquer les dysfonctionnements.
    \item \textbf{La description technique} (composants, paramètres, normes) pour situer l'objet.
\end{itemize}
\end{definition}

\courssub{Outils linguistiques obligatoires dans le rapport}

\begin{center}
\begin{tabularx}{\linewidth}{|l|X|X|}\hline
\rowcolor{popblue}\textbf{Outil linguistique} & \textbf{Fonction dans le rapport} & \textbf{Exemples (Contexte Nkongsamba / SONARA)} \\\hline
\textbf{Voix passive / Tournures impersonnelles} & Gommer l'énonciateur, focaliser sur le fait/objet. & \textit{Le joint \textbf{a été remplacé} ; Il \textbf{convient de} nettoyer ; \textbf{Aucune perte n'est à déplorer}.} \\\hline
\textbf{Temps verbaux : PC / Imparfait / Présent} & PC = actions ponctuelles ; Imparfait = contexte/état ; Présent = vérités techniques / recommandations. & \textit{L'équipe \textbf{a été sollicitée} (PC) ; La température \textbf{excédait} (Imparfait) la norme ; La sonde \textbf{permet} (Présent) la surveillance.} \\\hline
\textbf{Connecteurs logiques (cause, conséquence, opposition, énumération)} & Structurer le raisonnement factuel. & \textit{\textbf{Due à} l'accumulation (cause) ; \textbf{Par conséquent}, le moteur a fondu (conséquence) ; \textbf{Toutefois}, l'extinction a été rapide (opposition) ; \textbf{Ensuite}, \textbf{puis} (chronologie).} \\\hline
\textbf{Nominalisation / Verbes d'état} & Densifier l'information, objectiver. & \textit{\textbf{L'usure} du joint (nom) au lieu de « le joint s'use » ; \textbf{L'inflammation} des copeaux ; \textbf{La conformité} de l'installation.} \\\hline
\textbf{Lexique spécialisé précis (voir section 5)} & Désigner sans ambiguïté. & \textit{Vanne V-104, joint torique JT-224 Viton, ligne MP 15 bar, VLEP, LOTO, disjoncteur différentiel 30 mA.} \\\hline
\end{tabularx}
\end{center}

\begin{attention}[Pièges stylistiques à bannir]
\begin{itemize}
    \item \textbf{Futur / Conditionnel d'incertitude :} « \textit{On va changer le joint} » $\rightarrow$ « \textit{Le joint sera remplacé / Il convient de remplacer le joint} ».
    \item \textbf{Adjectifs subjectifs :} « \textit{grave, important, minime, choquant, super} » $\rightarrow$ Remplacer par des \textbf{mesures} : « \textit{brûlure 1\textsuperscript{er} degré, surface 3 cm $\times$ 2 cm} », « \textit{coût 420 000 FCFA} ».
    \item \textbf{Connecteurs argumentatifs émotionnels :} « \textit{malheureusement, heureusement, scandaleusement, évidemment} ».
\end{itemize}
\end{attention}

\coursec{Lexique commun et lexique spécialisé}

\courssub{Distinction et enjeu professionnel}

\begin{definition}[Lexique commun vs Lexique spécialisé (Sémantique/Lexicologie)]
\begin{itemize}
    \item \textbf{Lexique commun (langue générale) :} Mots compris de tous, sens stable, contexte courant (ex. : \textit{machine, feu, réparation, chef, pièce, nettoyer}).
    \item \textbf{Lexique spécialisé (langue de spécialité / LSP) :} Mots propres à un domaine professionnel (menuiserie, maintenance, chimie, BTP, agroalimentaire), sens technique précis, souvent normalisés (ex. : \textit{scie circulaire à format, extincteur CO\textsubscript{2}, joint torique Viton, disjoncteur différentiel 30 mA, LOTO, VLEP, égrenage, soufflerie}).
\end{itemize}
\end{definition}

\noindent Le passage du commun au spécialisé marque la \textbf{professionnalisation} de l'élève. Au Probatoire, l'usage pertinent du lexique technique est un critère d'évaluation majeur.

\begin{methode}[Construire son répertoire lexical spécialisé (Fiche méthode)]
\begin{enumerate}
    \item \textbf{Collecter :} Lors des TP, stages, visites, noter les termes exacts utilisés par les professionnels (enseignants, tuteurs, notices machines).
    \item \textbf{Classifier :} Par système/fonction (ex. : \textit{Électrique : disjoncteur, contacteur, variateur ; Mécanique : roulement, palier, vilebrequin ; Sécurité : EPI, LOTO, VLEP, MSDS}).
    \item \textbf{Contextualiser :} Écrire une phrase de définition technique pour chaque terme (ex. : « Le \textbf{LOTO} (Lock Out Tag Out) est une procédure de consignation d'énergie pour maintenance. »).
    \item \textbf{Utiliser :} Réinvestir systématiquement dans les rapports de TP, de stage, d'incident.
\end{enumerate}
\end{methode}

\begin{exemplebox}[Glissement commun $\rightarrow$ spécialisé (Exemples 1\textsuperscript{re} Tech)]
\begin{center}
\begin{tabular}{|l|l|l|}\hline
\rowcolor{popblue}\textbf{Domaine} & \textbf{Lexique commun (À éviter dans le corps technique)} & \textbf{Lexique spécialisé (Exigé)} \\\hline
Menuiserie / Bois & scie, table, copeaux, panneau, colle & \textbf{Scie circulaire à format}, \textbf{table de travail}, \textbf{copeaux / sciure}, \textbf{panneau mélaminé / contreplaqué Okoumé}, \textbf{colle vinylique / polyuréthane} \\\hline
Maintenance / Mécanique & machine cassée, pièce usée, huile, vis, bruit & \textbf{Panne critique}, \textbf{usure anormale}, \textbf{lubrifiant ISO VG 220}, \textbf{vis CHC M10$\times$30}, \textbf{vibration / niveau sonore 92 dB(A)} \\\hline
Électricité & courant, fil, disjoncteur, prise, panne & \textbf{Tension 380V triphasé}, \textbf{câble HO7RNF 4G6}, \textbf{disjoncteur différentiel 30 mA}, \textbf{prise CEE 32A}, \textbf{défaut d'isolement} \\\hline
Sécurité / HSE & danger, protection, masque, gants, règle & \textbf{Risque / Dangers (ISO 12100)}, \textbf{EPI (Équipement Protection Individuelle)}, \textbf{masque FFP2 / FFP3}, \textbf{gants anti-coupure niveau 5}, \textbf{Consigne de sécurité / Mode opératoire} \\\hline
Agroalimentaire (SODECOTON/UCAM) & coton, bière, machine, tuyau, nettoyage & \textbf{Coton-graine / Fibre / Graine}, \textbf{Moût / Bière / Taux d'alcool}, \textbf{Égreneuse Continental Eagle}, \textbf{Tuyauterie inox 316L}, \textbf{NEP (Nettoyage En Place)} \\\hline
\end{tabular}
\end{center}
\end{exemplebox}

\coursec{Lecture de textes fonctionnels et affiches de sécurité}

\courssub{Les textes fonctionnels : définition et lecture efficace}

\begin{definition}[Textes fonctionnels]
Textes à \textbf{visée utilitaire} (pas esthétique), produits dans un cadre professionnel/administratif pour \textbf{agir} ou \textbf{faire agir}. Le rapport en est l'archétype. Autres exemples : modes opératoires, consignes de sécurité, fiches techniques, comptes rendus, PV, cahiers des charges, mails professionnels.
\end{definition}

\begin{methode}[Lecture efficace d'un texte fonctionnel (Stratégie « 3 lectures »)]
\begin{enumerate}
    \item \textbf{Lecture 1 : Identification (Métadonnées).} Qui écrit ? À qui ? Pourquoi ? Quand ? Quel type de document ? (En-tête, objet, référence).
    \item \textbf{Lecture 2 : Compréhension de la structure.} Repérer les parties (Intro, Corps, Conclusion), les titres, les listes, les données chiffrées, les annexes.
    \item \textbf{Lecture 3 : Extraction de l'information cible.} Selon la question posée : « Quel est le coût ? », « Quelle est la cause racine ? », « Quelle action est demandée à qui et pour quand ? ». Surligner / annoter.
\end{enumerate}
\end{methode}

\courssub{Les affiches d'indications sécuritaires : lecture iconique}

Conformément au programme (Réception des textes : Textes iconiques), l'élève doit savoir décoder les \textbf{signes de sécurité normalisés (ISO 7010)}. Ce sont des « textes » visuels à lire comme des phrases.

\begin{popfigure}
\centering
\begin{tikzpicture}[
    scale=1.1,
    sign/.style={draw=popdark, thick, minimum size=2.2cm, rounded corners=2pt, inner sep=2pt},
    inter/.style={sign, fill=white, circle, path picture={\draw[popred, line width=4pt] (path picture bounding box.south west) -- (path picture bounding box.north east); \draw[popred, line width=4pt] (path picture bounding box.north west) -- (path picture bounding box.south east);}},
    danger/.style={sign, shape=regular polygon, regular polygon sides=3, fill=popgoldL, rotate=-90, minimum size=2.5cm, inner sep=0, path picture={\node[popdark, font=\Large\bfseries] at (path picture bounding box.center) {!};}},
    oblig/.style={sign, circle, fill=popblue, minimum size=2.2cm, path picture={\node[white, font=\Large\bfseries] at (path picture bounding box.center) {+};}},
    secours/.style={sign, rectangle, fill=popgreen, minimum size=2.2cm, path picture={\node[white, font=\Large\bfseries] at (path picture bounding box.center) {+};}},
    label/.style={below, font=\small, align=center, text width=2.5cm, node distance=3mm}
]
% Interdiction
\node[inter] (i1) {};
\node[label, align=center] at (i1.south){\textbf{INTERDICTION}\\Rond rouge + barre\\Ex: Interdit de fumer, Accès interdit};
% Danger
\node[danger, right=1.5cm of i1] (d1) {};
\node[label, align=center] at (d1.south){\textbf{DANGER / AVERTISSEMENT}\\Triangle jaune + contour noir\\Ex: Haute tension, Produit toxique, Chute};
% Obligation
\node[oblig, right=1.5cm of d1] (o1) {};
\node[label, align=center] at (o1.south){\textbf{OBLIGATION}\\Rond bleu + pictogramme blanc\\Ex: Port casque, EPI, LOTO, Lavage mains};
% Secours / Sauvetage
\node[secours, right=1.5cm of o1] (s1) {};
\node[label, align=center] at (s1.south){\textbf{SECOURS / SAUVETAGE}\\Carré/Rect. vert + pictogramme blanc\\Ex: Sortie, Extincteur, Poste secours, Douche oculaire};
% Légende codes couleurs
\node[below=1.5cm of d1, font=\small, align=center, text width=10cm] {
\textbf{Code couleur ISO 7010 (Universel) :} \textcolor{popred}{\textbf{Rouge}} = Interdiction / Arrêt / Matériel incendie \quad \textcolor{popgold}{\textbf{Jaune}} = Danger / Attention \quad \textcolor{popblue}{\textbf{Bleu}} = Obligation / Contrainte \quad \textcolor{popgreen}{\textbf{Vert}} = Secours / Issue / Sécurité
};
\end{tikzpicture}
\legende{Les 4 familles de panneaux de sécurité (ISO 7010). L'élève doit les identifier, en expliquer le sens et en respecter la consigne (ex. : rapport d'incident mentionnant « extincteur CO\textsubscript{2} \textcolor{popred}{\textbf{(rouge)}} » ou « zone LOTO \textcolor{popblue}{\textbf{(bleu)}} »).}
\end{popfigure}

\begin{exemplebox}[Analyse d'une affiche composite (Atelier menuiserie)]
\textit{Affiche observée à l'entrée de l'atelier :}
\begin{itemize}
    \item \textcolor{popblue}{\textbf{[Rond bleu]}} Pictogramme lunettes + casque anti-bruit + gants $\rightarrow$ \textbf{Obligation :} Port EPI obligatoire (yeux, ouïe, mains).
    \item \textcolor{popred}{\textbf{[Rond rouge + barre]}} Pictogramme téléphone portable $\rightarrow$ \textbf{Interdiction :} Téléphone portable interdit (distraction, risque étincelle).
    \item \textcolor{popgold}{\textbf{[Triangle jaune]}} Pictogramme lame circulaire $\rightarrow$ \textbf{Danger :} Risque coupure / projection (scies).
    \item \textcolor{popgreen}{\textbf{[Carré vert]}} Pictogramme extincteur + flèche $\rightarrow$ \textbf{Secours :} Emplacement extincteur (CO\textsubscript{2} pour électrique).
\end{itemize}
\textbf{Transposition en phrase de rapport :} « L'affichage de sécurité à l'entrée de l'atelier (conforme ISO 7010) prescrit le port d'EPI (lunettes, casque, gants), interdit l'usage du téléphone portable, signale le risque de coupure sur machines tournantes et localise l'extincteur CO\textsubscript{2}. »
\end{exemplebox}

\coursec{Production et amélioration d'un rapport complet}

\courssub{Méthodologie globale : de la commande à l'envoi}

\begin{methode}[Produire un rapport complet : 5 étapes]
\begin{enumerate}
    \item \textbf{ANALYSER LA COMMANDE :} Qui demande ? Quoi ? Pour quand ? Quel format (papier, mail, logiciel GMAO) ? Quel niveau de détail ?
    \item \textbf{COLLECTER LES FAITS (Terrain) :} Observer, mesurer, photographier, interroger (témoins), consulter registres (maintenance, stock, infirmerie). \textbf{Noter brut (carnet de terrain).}
    \item \textbf{STRUCTURER ET RÉDIGER (Bureau) :}
    \begin{itemize}
        \item Remplir l'en-tête (Réf, Objet, De, À, Date, P.J.).
        \item Rédiger l'introduction (Mandat + Fait + Plan).
        \item Rédiger le corps (Chrono OU Thématique) : Faits $\rightarrow$ Bilan chiffré $\rightarrow$ Analyse causes (5M : Main d'œuvre, Matériel, Milieu, Méthode, Management).
        \item Rédiger la conclusion (Bilan + Recommandations SMART).
    \end{itemize}
    \item \textbf{RELIRE ET AMÉLIORER (Contrôle qualité) :} Appliquer la check-list « Objectivité / Exactitude / Clarté / Concision / Lexique spécialisé / Temps verbaux / Mise en page / Signature / Visa / P.J. ».
    \item \textbf{TRANSMETTRE ET ARCHIVER :} Remise en main propre / mail avec accusé de réception / dépôt GED. Conserver une copie (archive personnelle).
\end{enumerate}
\end{methode}

\courssub{Exercice d'entraînement : Production guidée (Sujet type Probatoire)}

\begin{exercice}[Sujet : Rapport d'incident atelier — 1h30]
\textbf{Situation :} Vous êtes élève en 1\textsuperscript{re} TMA (Technique Menuisier Agenceur) au Lycée Technique de Douala-Bonabéri. Le mardi 12 novembre 2024, vers 10 h 15, lors d'un TP de dégauchissage sur la machine n\textsuperscript{o} DG-200 (n\textsuperscript{o} série 5544-D), l'arbre porte-outils s'est bloqué brutalement. Un éclat de métal a été projeté, brisant la lunette de protection de votre camarade K. Paul (blessure superficielle paupière gauche, évacué infirmerie). La machine a été arrêtée par le bouton d'urgence. L'enseignant M. Ngono a constaté la casse d'un couteau HSS (réf. CT-HSS-50x12x1,5) et l'absence de la protection arrière du carter (démontée pour maintenance non remontée).

\textbf{Consignes :}
\begin{enumerate}
    \item Rédigez le \textbf{rapport d'incident} adressé au Proviseur (M. Essomba), référence LTN-DLA/2024/0312.
    \item Respectez la structure externe et interne complète.
    \item Utilisez le lexique spécialisé menuiserie/maintenance/sécurité.
    \item Appliquez les règles d'objectivité (temps, voix, connecteurs, chiffres).
    \item Formulez 3 recommandations SMART.
\end{enumerate}
\end{exercice}

\begin{corrige}[Exercice n°1 — Proposition de rapport type Probatoire]
\textbf{RAPPORT D'INCIDENT N\textsuperscript{o} 2024/0312/LTN-DLA} \\[1mm]
\textbf{Objet :} Blocage arbre porte-outils dégauchisseuse DG-200 — Projection éclat — Blessé léger \\[1mm]
\textbf{De :} [Prénom Nom], Élève délégué 1\textsuperscript{re} TMA \\
\textbf{À :} M. le Proviseur, Lycée Technique de Douala-Bonabéri \\
\textbf{Date :} 12 novembre 2024 \hfill \textbf{Heure :} 15 h 45 \\
\textbf{C/c :} Chef de travaux Menuiserie ; Infirmerie \\
\textbf{P.J. :} 1. Photos machine (2) \quad 2. Fiche intervention infirmerie \quad 3. Fiche vie machine DG-200 \\[2mm]

\textbf{1. INTRODUCTION} \\
Suivant la consigne permanente de signalement des incidents d'atelier, je vous présente le rapport relatif au dysfonctionnement survenu sur la dégauchisseuse n\textsuperscript{o} DG-200 le 12 novembre 2024. Cet exposé détaille la chronologie, le bilan, l'analyse des causes et les mesures correctives. \\[2mm]

\textbf{2. CORPS DU RAPPORT} \\
\textbf{2.1 Chronologie des faits} \\
\begin{itemize}
    \item 10 h 12 : Début TP dégauchissage panneau hêtre massif (ép. 40 mm), groupe de 3 élèves poste DG-200. EPI portés (lunettes, casque, gants).
    \item 10 h 15 : Blocage brutal arbre porte-outils (vitesse 5000 tr/min). Projection éclat métallique (couteau HSS cassé) vers l'avant.
    \item 10 h 15 : Impact lunette protection élève K. Paul (paupière gauche). Arrêt d'urgence actionné par élève poste adjacent.
    \item 10 h 16 : Intervention M. Ngono (enseignant). Mise hors tension (sectionneur général). Appel infirmerie.
    \item 10 h 20 : Prise en charge K. Paul (lavage oculaire, collyre, pansement). Avis médecin non requis. Reprise cours 11 h 00.
\end{itemize}

\textbf{2.2 Bilan des dommages} \\
\begin{itemize}
    \item \textbf{Humain :} 1 blessé léger (K. Paul, 1\textsuperscript{re} TMA) : éraflure conjonctivale gauche, sans atteinte cornéenne (certificat infirmerie).
    \item \textbf{Matériel :} Machine DG-200 : 1 couteau HSS cassé (réf. CT-HSS-50x12x1,5, valeur 18 500 FCFA). Arbre porte-outils : contrôle jeu radial à prévoir. Carter protection arrière : absent (démonté maintenance préventive 05/11/2024, non remonté).
    \item \textbf{Production :} Arrêt machine 45 min. Perte temps TP : 1 groupe (3 élèves) sur 4.
\end{itemize}

\textbf{2.3 Analyse des causes (Méthode 5M)} \\
\begin{itemize}
    \item \textbf{Matériel :} Couteau HSS fragilisé (usure > 80 h affûtage, micro-fissures non détectées contrôle visuel). Absence protection arrière carter (barrière physique manquante).
    \item \textbf{Méthode :} Procédure remontage carter post-maintenance non formalisée / non vérifiée (check-list absente).
    \item \textbf{Main d'œuvre / Management :} Maintenance préventive du 05/11 non clôturée (non-repos carter). Surveillance conformité EPI efficace (lunettes ont protégé l'œil).
    \item \textbf{Milieu :} Éclairage correct. Sol sec, non glissant.
\end{itemize}

\textbf{3. CONCLUSION ET RECOMMANDATIONS} \\
L'incident, maîtrisé grâce au port des EPI et à la réactivité de l'arrêt d'urgence, révèle une défaillance du processus de maintenance (non-remontage protection) et du suivi outillage (couteau en fin de vie). \\[1mm]
\textbf{Recommandations immédiates (SMART) :} \\
\begin{enumerate}
    \item \textbf{Remonter et verrouiller} le carter protection arrière machine DG-200 \textbf{ce jour 12/11 avant 16 h} — \textit{Responsable : M. Ngono / Agent maintenance}.
    \item \textbf{Remplacer les 4 couteaux} jeu complet DG-200 (réf. CT-HSS-50x12x1,5) \textbf{avant reprise TP 14/11} — \textit{Coût : 74 000 FCFA (budget atelier)} — \textit{Responsable : Chef de travaux}.
    \item \textbf{Instaurer check-list « Remise en service post-maintenance »} (points : protections, outillages, alignement, essai à vide) \textbf{opérationnelle au 25/11/2024} — \textit{Responsable : Chef de travaux / Qualité}.
\end{enumerate}
\textbf{Signature :} \hfill \textit{[Prénom Nom]} \\[1mm]
\textbf{Visa Enseignant :} \hfill \textit{M. Ngono} \\[1mm]
\textbf{Visa Chef de travaux :} \hfill \textit{[Nom]}
\end{corrige}

\courssub{Exercice d'application : Repérage et correction (Style)}

\begin{exercice}[Entraînement : Repérage et correction]
\textbf{Consigne :} Le paragraphe ci-dessous est extrait d'un « rapport » rédigé par un élève de 1\textsuperscript{re} Tech après une visite à la SODECOTON Maroua. Il contient 5 erreurs graves contre l'objectivité et le style du rapport. Relevez-les et réécrivez le paragraphe correctement.

\textit{« Moi j'ai trouvé l'usine super grande et impressionnante. Les machines font un bruit d'enfer, c'est dingue ! Heureusement que j'avais mes bouchons d'oreilles. Le chef d'atelier nous a dit que le coton vient du Nord, c'est logique. Par contre, je trouve qu'il y a trop de poussière, c'est pas très safe pour les poumons. Ils devraient vraiment mettre des aspirateurs partout. En tout cas, moi j'ai kiffé la visite, c'était trop bien pour mon orientation. »}
\end{exercice}

\begin{corrige}[Exercice n°1]
\textbf{Erreurs repérées :}
\begin{enumerate}
    \item \textbf{Subjectivité affective :} « super grande », « impressionnante », « bruit d'enfer », « c'est dingue », « j'ai kiffé », « trop bien ».
    \item \textbf{Présence du « Je » énonciatif :} « Moi j'ai trouvé », « j'avais », « je trouve », « moi j'ai kiffé ».
    \item \textbf{Langage familier / oral :} « c'est pas très safe », « kiffé », « trop ».
    \item \textbf{Absence de précision technique :} « machines » (lesquelles ?), « poussière » (quelle nature ? taux ?), « aspirateurs » (type ?).
    \item \textbf{Recommandation non professionnelle :} « Ils devraient vraiment mettre... » (donner des ordres au lieu de formuler une préconisation technique).
\end{enumerate}

\textbf{Proposition de réécriture (style rapport de visite) :}
\textit{« La visite de l'unité d'égrenage de la SODECOTON Maroua, effectuée le 20 novembre 2024, a porté sur une surface de production de 12 000 m\textsuperscript{2}. Le niveau sonore mesuré au poste d'égreneuse n\textsuperscript{3} (machine Continental Eagle, modèle 2018) atteint 92 dB(A), imposant le port de protections auditives individuelles (fournies). L'approvisionnement en coton-graine provient exclusivement de la zone Nord (campagne 2023-2024 : 240 000 tonnes). Une concentration de poussières de coton (particules < 10 µm) a été observée visuellement dans la zone de soufflerie, dépassant visuellement la VLEP (Valeur Limite d'Exposition Professionnelle) de 10 mg/m\textsuperscript{3} pour poussières totales. \textbf{Recommandation :} Étude d'installation d'un système de captage à la source avec filtration à manches sur les lignes de soufflerie, conformément à la norme ISO 16890. »}
\end{corrige}

\courssub{Synthèse : Le rapport, compétence transversale}

Produire un rapport en Première technique, ce n'est pas seulement réussir un exercice scolaire. C'est acquérir une \textbf{compétence transférable} :
\begin{itemize}
    \item \textbf{Scolaire :} Rapport de TP sciences, rapport de stage (obligatoire pour le Probatoire), compte rendu de visite.
    \item \textbf{Professionnel :} Rapport d'intervention maintenance, rapport de chantier BTP, rapport de contrôle qualité industrie agroalimentaire (SODECOTON, SOCAPALM, SONARA).
    \item \textbf{Administratif / Citoyen :} Déclaration de sinistre assurance, constat amiable, signalement voirie, témoignage écrit.
\end{itemize}

\begin{aretenir}[Check-list « Mon rapport est-il prêt ? »]
Avant de rendre votre copie ou d'envoyer votre mail :
\begin{itemize}
    \item[$\square$] L'en-tête est-il complet (Expéditeur, Destinataire, Objet, Date, Référence, C/c, P.J.) ?
    \item[$\square$] L'introduction pose-t-elle le contexte (Qui ? Quoi ? Où ? Quand ? Pourquoi ? Plan) ?
    \item[$\square$] Le corps suit-il une logique claire (Chronologie OU Thématique) avec des faits \textbf{chiffrés} ?
    \item[$\square$] L'analyse des causes est-elle technique, non accusatoire, étayée (5M, Arbre des causes) ?
    \item[$\square$] La conclusion synthétise-t-elle et propose-t-elle des \textbf{actions concrètes, datées, chiffrées, responsables} (SMART) ?
    \item[$\square$] Le vocabulaire est-il \textbf{spécialisé} (lexique technique) et \textbf{neutre} (pas d'adjectifs subjectifs) ?
    \item[$\square$] Les temps verbaux sont-ils cohérents (PC/Imparfait pour faits, Présent pour vérités/recos) ?
    \item[$\square$] La voix passive / impersonnelle domine-t-elle ? (Pas de « Je » affectif).
    \item[$\square$] La mise en page est-elle aérée (paragraphes courts, listes à puces, gras sur mots-clés) ?
    \item[$\square$] Signature + Visa hiérarchique + Pièces jointes mentionnées ?
\end{itemize}
\end{aretenir}

\coursec{La structure externe du rapport}

Dans la section précédente, nous avons découvert ce qu'est un rapport : un document écrit, objectif et structuré, par lequel une personne rend compte à une autorité (chef de service, principal, délégué, ministre) d'une mission, d'un événement ou d'une situation qu'elle a observée ou étudiée. Nous avons vu que le rapport se distingue de la lettre ordinaire par son caractère administratif et par son organisation rigoureuse. Mais comment reconnaît-on, au premier coup d'œil, qu'un document est un rapport ? C'est précisément ce que nous allons apprendre maintenant : avant même de lire le contenu, un rapport se reconnaît à sa \cle{présentation matérielle}, c'est-à-dire à sa structure externe. C'est elle que nous étudions dans cette section.

\courssub{Observation guidée : que voit-on sur un rapport ?}

Commençons par une démarche d'investigation. Le professeur distribue à la classe deux documents authentiques (ou reproduits au tableau) : d'une part le rapport d'un chef d'atelier sur une panne de machine, d'autre part le rapport d'un surveillant général sur un incident survenu pendant les cours. Lisons-les rapidement, sans entrer dans le fond, et répondons aux questions suivantes.

\begin{experience}[Observer pour découvrir]
\textbf{Matériel :} deux rapports administratifs réels ou reconstitués.\par
\textbf{Protocole :} observer chaque document pendant trois minutes, de haut en bas, et noter tout ce qui apparaît \emph{avant} le texte proprement dit et \emph{après} lui.\par
\textbf{Observations attendues :}
\begin{itemize}
\item En haut à gauche : le nom et la fonction de celui qui écrit (ex. : « Moussa Ibrahim, Chef d'atelier de maintenance ») ;
\item En haut à droite : le nom et la fonction de celui qui reçoit (ex. : « Monsieur le Proviseur »), ainsi qu'un lieu et une date ;
\item Une ligne commençant par « Objet : » ;
\item Un texte divisé en paragraphes, parfois numérotés, parfois précédés de petits titres ;
\item En bas : une signature, et parfois des mentions comme « Pièce jointe » ou « Ampliation ».
\end{itemize}
\textbf{Interprétation :} ces éléments reviennent dans les deux documents, toujours à la même place. Ils ne dépendent donc pas du sujet traité : ils constituent le \emph{cadre fixe}, la « vitrine » de tout rapport.
\end{experience}

Cette observation nous conduit à la définition suivante.

\begin{definition}[La structure externe du rapport]
La \cle{structure externe} d'un rapport est l'ensemble des \textbf{éléments matériels d'identification et de présentation} qui encadrent le texte : l'\cle{en-tête} (auteur et destinataire), les \cle{références}, l'\cle{objet}, la mention du \cle{lieu} et de la \cle{date}, le \cle{corps du document} (titre, paragraphes, intertitres), la \cle{signature} et les \cle{mentions finales} (pièces jointes, ampliations). Ces éléments occupent des places conventionnelles sur la page et permettent d'identifier immédiatement qui écrit, à qui, à quel sujet, quand et où.
\end{definition}

\begin{aretenir}[Les six zones visuelles de la page]
Un rapport bien présenté comporte six zones repérables à l'œil nu : (1) l'en-tête (auteur), (2) le destinataire, le lieu et la date, (3) les références et l'objet, (4) le corps du document, (5) la signature, (6) les mentions finales (pièces jointes, ampliations). Aucune ne doit manquer.
\end{aretenir}

\courssub{L'en-tête : identifier l'auteur et le destinataire}

L'\cle{en-tête} occupe la partie supérieure de la page. Il répond à deux questions : \emph{qui écrit ?} et \emph{à qui ?}

\begin{itemize}
\item \textbf{En haut à gauche} figurent l'identité de l'auteur : son \textbf{nom}, sa \textbf{fonction} et son \textbf{service} (ou son établissement). On peut y ajouter l'intitulé officiel de la structure (ministère, lycée, entreprise), parfois sous forme d'en-tête imprimé.
\item \textbf{En haut à droite} (ou sous le bloc de gauche) figure le \textbf{destinataire}, désigné par sa fonction : « Monsieur le Principal », « Monsieur le Délégué régional des Enseignements secondaires du Littoral ».
\end{itemize}

\begin{exemplebox}[En-tête d'un rapport]
\begin{center}
\begin{tabularx}{\linewidth}{@{}X@{}X@{}}
\textbf{MINISTÈRE DES ENSEIGNEMENTS SECONDAIRES} & \\
Lycée technique de Deïdo & \\
Le Chef de travaux, & \textbf{À Monsieur le Délégué régional} \\
M. ONANA Jean-Claude & \textbf{des Enseignements secondaires} \\
 & \textbf{du Littoral} \\
\end{tabularx}
\end{center}
On lit immédiatement que M. Onana, chef de travaux dans un lycée technique, s'adresse à son autorité hiérarchique régionale.
\end{exemplebox}

\begin{attention}[Erreur fréquente]
Ne confonds pas la lettre personnelle et le rapport : dans un rapport, on n'écrit jamais « Cher Monsieur » ni de formule de politesse développée. Le destinataire est désigné par sa \textbf{fonction}, non par une formule affectueuse. De même, l'auteur se présente par sa fonction, pas seulement par son nom.
\end{attention}

\courssub{Les références et l'objet}

Sous l'en-tête, on trouve souvent une ligne de \cle{références} : elle indique le numéro d'enregistrement du document dans le service (ex. : « N° 045/LTD/CT/26 ») et peut rappeler les documents antérieurs auxquels le rapport se rattache (« Réf. : votre lettre n° 12 du 3 mars »). Les références permettent de classer et de retrouver le document.

Vient ensuite la ligne la plus importante de la structure externe : l'\cle{objet}. L'objet annonce, en quelques mots, le sujet du rapport. Il doit être \textbf{bref} et \textbf{précis}.

\begin{methode}[Formuler l'objet d'un rapport]
\begin{enumerate}
\item Repère le \textbf{sujet exact} du rapport : de quoi rend-on compte ?
\item Choisis un \textbf{nom} qui désigne ce sujet (incendie, panne, visite, stage, incident\ldots) et fais-le précéder du mot « Rapport sur\ldots » ou « Compte rendu de\ldots ».
\item Ajoute les \textbf{précisions indispensables} : le lieu, la date, le service concerné, sous forme de compléments du nom.
\item Rédige une \textbf{phrase nominale} (sans verbe conjugué), d'une seule ligne si possible.
\end{enumerate}
\textbf{Modèle :} « Objet : Rapport sur l'incendie de l'atelier du 12 mars. »
\end{methode}

\begin{attention}[L'objet : deux pièges à éviter]
\begin{itemize}
\item \textbf{Oublier l'objet} : sans objet, le destinataire ne sait pas de quoi traite le document avant de le lire ; le rapport perd sa valeur administrative.
\item \textbf{Formuler l'objet comme une phrase verbale longue} : « Objet : Je vous écris pour vous raconter que l'atelier a pris feu le 12 mars et que nous avons essayé d'éteindre » est incorrect. L'objet est une \emph{phrase nominale courte} : « Rapport sur l'incendie de l'atelier du 12 mars ».
\end{itemize}
\end{attention}

\courssub{La mention du lieu et de la date}

Tout rapport est \textbf{daté} et \textbf{localisé}. La mention se place en haut à droite, au-dessus ou au-dessous du destinataire, selon la formule consacrée : « Douala, le 15 mars 2026 ». La date est essentielle : un rapport sert souvent de preuve dans un dossier administratif ; sans date, il ne vaut rien. On écrit le mois en toutes lettres (« 15 mars 2026 » et non « 15/03/26 » dans un rapport formel).

\courssub{Le corps du document : titre, paragraphes, intertitres}

Entre l'objet et la signature se trouve le \cle{corps du document}, c'est-à-dire le texte du rapport lui-même. Sa présentation externe obéit à des règles simples :

\begin{itemize}
\item Un \textbf{titre} éventuel, en gras et centré, peut annoncer le thème général (surtout dans les rapports longs) ;
\item Le texte est découpé en \textbf{paragraphes} marqués par des alinéas ; chaque paragraphe traite une seule idée ;
\item Les parties peuvent être \textbf{numérotées} (I, II, III ou 1, 2, 3) et précédées d'\textbf{intertitres} (petits titres de parties, ex. : « 1. Circonstances de l'incident », « 2. Dégâts constatés », « 3. Mesures proposées ») ;
\item La mise en page doit être \textbf{aérée} : marges régulières, espaces entre les paragraphes, pas de ratures.
\end{itemize}

Le contenu de ce corps (introduction, développement, conclusion) fera l'objet de la section suivante ; retenons ici seulement sa \emph{présentation matérielle}.

\courssub{La signature et les mentions finales}

En bas à droite du document, l'auteur \textbf{signe}, précédé ou suivi de son nom et de sa fonction en toutes lettres : « Le Chef d'atelier, [signature], MOUSSA Ibrahim ». La signature engage la responsabilité de l'auteur : c'est elle qui donne au rapport sa valeur officielle.

Sous la signature, à gauche, peuvent figurer deux mentions finales :

\begin{itemize}
\item \textbf{Pièces jointes (P. J.)} : la liste des documents annexés au rapport (photos, factures, procès-verbaux). Ex. : « P. J. : 3 photos de l'atelier ; 1 devis de réparation. »
\item \textbf{Ampliations} : la liste des autres personnes ou services qui reçoivent une copie du rapport.
\end{itemize}

\begin{savaistu}[À quoi servent les « ampliations » ?]
Le mot \emph{ampliation} vient du latin \emph{ampliatio} (« élargissement ») : il désigne les destinataires auxquels on « élargit » l'envoi. Dans l'administration camerounaise de l'enseignement technique, on écrit par exemple : « Ampliations : MINESEC ; DRES ; archives. » Chacun sait ainsi qui d'autre est informé, ce qui évite les doublons et garantit la transparence du circuit administratif. Au Probatoire et au Baccalauréat technique, la présence correcte de cette mention peut valoir des points de présentation.
\end{savaistu}

\courssub{Schéma récapitulatif : la page d'un rapport}

\begin{popfigure}
\begin{center}
\begin{tikzpicture}
\draw[thick,popdark] (0,0) rectangle (10,13);
% Zone 1: En-tête Auteur
\draw[fill=popblueL,draw=popblue] (0.5,11.2) rectangle (4.8,12.6);
\node[align=center,font=\small] at (2.65,11.9) {1. En-tête : auteur\\(nom, fonction, service)};
% Zone 2: Destinataire, Lieu, Date
\draw[fill=popblueL,draw=popblue] (5.2,11.2) rectangle (9.5,12.6);
\node[align=center,font=\small] at (7.35,11.9) {2. Destinataire\\Lieu et date};
% Zone 3: Références et Objet
\draw[fill=poporangeL,draw=poporange] (0.5,10.2) rectangle (9.5,10.9);
\node[align=center,font=\small] at (5,10.55) {3. Références \quad --- \quad \textbf{Objet :} phrase nominale brève};
% Zone 4: Corps du document
\draw[fill=popgreenL,draw=popgreen] (0.5,4.6) rectangle (9.5,9.9);
\node[align=center,font=\small] at (5,9.2) {4. Corps du document};
\node[align=left,font=\small] at (5,7.5) {Titre éventuel (centré)\\Paragraphes avec alinéas\\Parties numérotées (1, 2, 3)\\Intertitres};
% Zone 5: Signature
\draw[fill=poppurpleL,draw=poppurple] (5.2,2.6) rectangle (9.5,4.3);
\node[align=center,font=\small] at (7.35,3.45) {5. Signature\\(nom et fonction)};
% Zone 6: Mentions finales
\draw[fill=popgoldL,draw=popgold] (0.5,0.3) rectangle (9.5,2.3);
\node[align=left,font=\small] at (5,1.3) {6. Mentions finales :\\Pièces jointes (P. J.) --- Ampliations};
\end{tikzpicture}
\end{center}
\legende{Schéma d'une page de rapport : les six zones de la structure externe. 1. En-tête (auteur) ; 2. Destinataire, lieu et date ; 3. Références et objet ; 4. Corps du document ; 5. Signature ; 6. Pièces jointes et ampliations.}
\end{popfigure}

\courssub{Exercices d'application}

\begin{exoresolu}[Reconstituer un rapport mélangé]
Les éléments externes suivants ont été mélangés. Remets-les dans l'ordre correct de la page (de haut en bas) :\par
a) « Objet : Rapport sur la panne de la fraiseuse de l'atelier mécanique. »\par
b) « Le Chef d'atelier, [signature], ABENA Samuel »\par
c) « À Monsieur le Proviseur du Lycée technique d'Ebolowa »\par
d) « P. J. : 1 devis ; Ampliations : DRES du Sud »\par
e) « Ebolowa, le 4 février 2026 »\par
f) « ABENA Samuel, Chef d'atelier mécanique »\par
\tcblower
\textbf{Solution :} l'ordre correct est : \textbf{f} (en-tête : auteur, en haut à gauche) ; \textbf{c} (destinataire, en haut à droite) ; \textbf{e} (lieu et date) ; \textbf{a} (objet, phrase nominale) ; puis viendrait le corps du document ; \textbf{b} (signature, en bas à droite) ; \textbf{d} (mentions finales : pièces jointes et ampliations, en bas à gauche).\par
\textbf{Vérification :} on lit d'abord qui écrit (f), à qui (c), où et quand (e), de quoi il s'agit (a), puis on trouve la signature (b) et les annexes (d). La logique de la page est respectée.
\end{exoresolu}

\begin{exercice}[Formuler un objet correct]
Voici trois objets rédigés par des élèves. Un seul est correct. Identifie-le et corrige les deux autres.\par
a) « Objet : Je vous informe que la machine à coudre de l'atelier de couture est tombée en panne hier et que les élèves n'ont pas pu travailler. »\par
b) « Objet : Rapport sur la panne de la machine à coudre de l'atelier de couture du 21 avril. »\par
c) « Objet : Panne. »
\end{exercice}

\begin{corrige}[Exercice : formuler un objet correct]
L'objet correct est le \textbf{b} : c'est une phrase nominale brève et précise (sujet + lieu + date).\par
Correction de a : supprimer la phrase verbale et rédiger « Rapport sur la panne de la machine à coudre de l'atelier de couture du 21 avril ».\par
Correction de c : « Panne » est trop vague : préciser de quelle machine, de quel atelier et de quel jour il s'agit.
\end{corrige}

\begin{exercice}[Rédiger la structure externe complète]
Tu es délégué de classe de Première technique. Le 9 mai 2026, un court-circuit a endommagé deux prises électriques de ta salle de classe au Lycée technique de Bafoussam. Rédige uniquement la \textbf{structure externe} du rapport que tu adresses au Proviseur : en-tête, destinataire, lieu et date, objet, signature, mentions finales (une photo jointe, copie au Censeur).
\end{exercice}

\begin{corrige}[Exercice : structure externe complète]
Éléments attendus (la formulation peut varier) :\par
En haut à gauche : « Le Délégué de la classe de Première technique, [Nom Prénom] ».\par
En haut à droite : « À Monsieur le Proviseur du Lycée technique de Bafoussam » ; « Bafoussam, le 9 mai 2026 ».\par
Objet : « Rapport sur le court-circuit survenu en salle de Première technique le 9 mai 2026. »\par
En bas à droite : « Le Délégué de classe, [signature], [Nom Prénom] ».\par
En bas à gauche : « P. J. : 1 photo ; Ampliations : Monsieur le Censeur. »
\end{corrige}

\begin{aretenir}[Bilan de la section]
La structure externe est le cadre matériel du rapport : en-tête (auteur : nom, fonction, service ; destinataire : fonction), références, objet en phrase nominale brève, lieu et date, corps présenté en paragraphes aérés et numérotés, signature, pièces jointes et ampliations. Maîtriser ce cadre, c'est garantir que le rapport sera pris au sérieux avant même d'être lu. La section suivante nous conduira à l'intérieur du texte : l'introduction, le corps et la conclusion du rapport.
\end{aretenir}

\coursec{La structure interne : introduction, corps, conclusion}

Dans la section précédente, nous avons étudié la \cle{structure externe} du rapport : l'en-tête, l'intitulé du service, l'objet, les références, la date, la signature. Ces éléments donnent au rapport son \emph{habillage administratif}. Mais une belle enveloppe ne suffit pas : il faut encore que le contenu soit organisé de manière claire et logique. C'est précisément l'objet de la \cle{structure interne}, que nous allons maintenant explorer.

Imagine un commerçant du marché Mokolo qui raconte un incendie à un assureur. S'il commence par la fin («~tout a brûlé~»), revient au début («~ce matin-là, j'étais chez mon fournisseur~»), puis mélange les faits avec ses projets d'avenir, son interlocuteur sera perdu. Un bon récit suit un ordre : d'où l'on vient, ce qui s'est passé, ce qu'on en conclut. Le rapport obéit exactement à cette logique, en trois temps.

\begin{definition}[La structure interne du rapport]
La \cle{structure interne} d'un rapport est l'organisation logique de son contenu en trois parties obligatoires :
\begin{itemize}
\item l'\cle{introduction}, qui situe le contexte et annonce le sujet ;
\item le \cle{corps} (ou développement), qui présente les faits de manière ordonnée ;
\item la \cle{conclusion}, qui dresse le bilan et formule éventuellement des suggestions.
\end{itemize}
Ces trois parties forment un tout cohérent : l'introduction pose la question, le corps apporte les éléments de réponse, la conclusion tire les enseignements.
\end{definition}

\begin{popfigure}
\begin{center}
\begin{tikzpicture}[font=\small]
% Entonnoir en trois étages
\fill[popblueL] (-4,4) -- (4,4) -- (2.8,2.2) -- (-2.8,2.2) -- cycle;
\fill[poporangeL] (-2.8,2.2) -- (2.8,2.2) -- (1.6,0.4) -- (-1.6,0.4) -- cycle;
\fill[popgreenL] (-1.6,0.4) -- (1.6,0.4) -- (0.6,-0.8) -- (-0.6,-0.8) -- cycle;
\node[align=center] at (0,3.1) {\textbf{INTRODUCTION}\\ contexte : qui ? quoi ? quand ? où ? pourquoi ?};
\node[align=center] at (0,1.3) {\textbf{CORPS}\\ faits ordonnés, constats, chiffres};
\node[align=center,text width=3.4cm] at (0,-0.1) {\textbf{CONCLUSION}\\ bilan + suggestions};
% Flèche centrale
\draw[-{Stealth[length=3mm]},popdark,very thick] (5,4) -- (5,-0.8) node[midway,right,align=left]{progression\\ logique};
\end{tikzpicture}
\end{center}
\legende{Schéma en entonnoir de la structure interne : du contexte général (introduction) vers les faits précis (corps), puis vers le bilan resserré (conclusion).}
\end{popfigure}

\courssub{L'introduction du rapport}

\textbf{Le rôle de l'introduction}\par
L'introduction est la porte d'entrée du rapport. En quelques lignes (cinq à dix lignes maximum), elle doit permettre au lecteur — souvent un supérieur hiérarchique très occupé — de comprendre immédiatement \emph{de quoi il s'agit} et \emph{pourquoi ce rapport a été rédigé}. Une introduction réussie contient trois éléments :

\begin{enumerate}
\item \textbf{Le contexte} : qui rédige, pour qui, quand, où. On rappelle la mission confiée ou l'événement survenu.
\item \textbf{Les circonstances} : l'ordre de mission, la demande de la hiérarchie, l'incident qui motive le rapport.
\item \textbf{L'annonce du plan} (facultative mais appréciée) : une phrase qui indique comment le rapport est organisé.
\end{enumerate}

\begin{methode}[Rédiger l'introduction d'un rapport]
Pour construire une introduction complète, réponds dans l'ordre aux cinq questions suivantes :
\begin{enumerate}
\item \textbf{Qui ?} — identifie l'auteur du rapport et sa fonction (ex. : «~le délégué de classe de Première F3~», «~l'infirmier de service~»).
\item \textbf{Quoi ?} — précise l'objet de la mission ou la nature de l'événement (visite d'entreprise, accident, stage, réunion).
\item \textbf{Quand ?} — indique la date ou la période exacte (ex. : «~le mardi 14 janvier 2026, de 8 h à 12 h~»).
\item \textbf{Où ?} — nomme le lieu précis (ex. : «~à la centrale hydroélectrique de Songloulou~»).
\item \textbf{Pourquoi ?} — explique la circonstance : ordre de mission, demande du chef d'établissement, incident à signaler.
\end{enumerate}
Termine, si nécessaire, par l'annonce du plan : «~Le présent rapport rend compte de cette visite en trois points : l'accueil, la visite des installations et les enseignements tirés.~»
\end{methode}

\begin{exemplebox}[Une introduction modèle]
\emph{Conformément à l'ordre de mission n° 12/DIR/2026 du 10 janvier 2026, les élèves de la classe de Première F3 du Lycée technique de Douala, accompagnés de deux enseignants, ont effectué le mardi 14 janvier 2026 une visite d'étude à la centrale hydroélectrique de Songloulou. L'objectif de cette sortie était d'observer le fonctionnement d'une unité de production d'électricité. Le présent rapport rend compte du déroulement de la visite, des observations faites et des enseignements tirés.}
\end{exemplebox}

Observe : en cinq lignes, le lecteur sait \textbf{qui} (les élèves de Première F3), \textbf{quoi} (une visite d'étude), \textbf{quand} (le 14 janvier 2026), \textbf{où} (Songloulou) et \textbf{pourquoi} (ordre de mission, objectif pédagogique). C'est le modèle à imiter.

\courssub{Le corps du rapport : présenter les faits de manière ordonnée}

\textbf{Deux ordres possibles}\par
Le corps est la partie la plus longue du rapport. Il présente les faits, les constats et les données recueillis. Mais des faits jetés en désordre ne valent rien : il faut choisir un \cle{ordre de présentation}.

\begin{aretenir}[Les deux ordres du corps de rapport]
\begin{itemize}
\item L'\cle{ordre chronologique} : on raconte les événements dans l'ordre où ils se sont déroulés (arrivée, visite de la salle des machines, visite du barrage, départ). Cet ordre convient aux rapports de visite, de voyage, d'événement.
\item L'\cle{ordre thématique} : on regroupe les faits par thèmes (état des bâtiments, effectif du personnel, matériel disponible). Cet ordre convient aux rapports d'inspection, d'inventaire, de diagnostic.
\end{itemize}
\end{aretenir}

\textbf{Le contenu du corps : faits, chiffres, témoignages}\par
Le corps d'un rapport ne contient que des \emph{faits vérifiables}, jamais des rumeurs ni des impressions personnelles déguisées. Trois types d'éléments nourrissent le corps :

\begin{itemize}
\item les \textbf{constats} : ce que l'auteur a vu ou entendu lui-même («~la salle des machines abrite quatre turbines~») ;
\item les \textbf{données chiffrées} : nombres, quantités, montants, durées, qui donnent au rapport sa précision et sa crédibilité ;
\item les \textbf{témoignages} : propos recueillis auprès des personnes concernées, cités avec leur identité et leur fonction («~selon M. Ngo Bassa, chef de poste, ...~»).
\end{itemize}

\begin{exemplebox}[Un corps de rapport intégrant des données chiffrées]
\emph{À notre arrivée à 8 h 30, le bilan provisoire de l'incendie du marché s'établissait comme suit : 47 boutiques entièrement détruites, 12 autres partiellement endommagées. Le nombre de sinistrés s'élève à 59 commerçants, dont 8 blessés légers pris en charge à l'hôpital de district. Selon l'évaluation conjointe du chef de marché et du service communal, le montant des dégâts est estimé à environ 85 000 000 FCFA, dont 60 000 000 FCFA de marchandises et 25 000 000 FCFA d'infrastructures. Les pompiers, alertés à 3 h 15, sont intervenus à 3 h 50 et ont maîtrisé le sinistre à 6 h 20.}
\end{exemplebox}

Remarque la précision : heures exactes, nombres de sinistrés, montants en FCFA. Ce sont ces chiffres qui distinguent un rapport sérieux d'un simple récit.

\textbf{L'organisation en parties et sous-parties avec intertitres}\par
Dès que le corps dépasse une page, il faut le découper en \cle{parties} numérotées (I, II, III ou 1, 2, 3), elles-mêmes subdivisées en \cle{sous-parties} si nécessaire. Chaque partie porte un \cle{intertitre} : un titre court et informatif qui annonce le contenu. Le lecteur pressé doit pouvoir comprendre l'essentiel du rapport en lisant seulement les intertitres.

\begin{popfigure}
\begin{center}
\begin{tikzpicture}[
grow=right,
level 1/.style={sibling distance=3.4cm, level distance=4.2cm},
level 2/.style={sibling distance=1.1cm, level distance=3.6cm},
font=\small,
every node/.style={draw, rounded corners, align=center, inner sep=2pt}
]
\node[fill=popblueL, text width=3cm]{\textbf{Rapport de visite}\\ centrale de Songloulou}
child { node[fill=popgreenL, text width=2.6cm]{III. Enseignements et suggestions}
  child { node[fill=popgreenL!40, text width=2.6cm]{b) Suggestions}}
  child { node[fill=popgreenL!40, text width=2.6cm]{a) Bilan pédagogique}}
}
child { node[fill=poporangeL, text width=2.6cm]{II. Déroulement de la visite}
  child { node[fill=poporangeL!40, text width=2.6cm]{c) Séance de questions}}
  child { node[fill=poporangeL!40, text width=2.6cm]{b) Le barrage}}
  child { node[fill=poporangeL!40, text width=2.6cm]{a) Salle des machines}}
}
child { node[fill=poppinkL, text width=2.6cm]{I. Accueil et présentation}
  child { node[fill=poppinkL!40, text width=2.6cm]{b) Historique de la centrale}}
  child { node[fill=poppinkL!40, text width=2.6cm]{a) Accueil par le chef de poste}}
};
\end{tikzpicture}
\end{center}
\legende{Arbre de plan d'un rapport de visite de la centrale hydroélectrique de Songloulou : trois parties numérotées, chacune subdivisée en sous-parties repérées par des intertitres.}
\end{popfigure}

\courssub{La conclusion du rapport}

\textbf{Le rôle de la conclusion}\par
La conclusion clôt le rapport. Elle est courte (trois à huit lignes) et contient deux éléments, plus une formule finale :

\begin{enumerate}
\item \textbf{Le bilan des faits} : on résume en une ou deux phrases l'essentiel de ce qui a été constaté, sans répéter tout le corps.
\item \textbf{Les suggestions ou recommandations} (éventuelles) : l'auteur propose des solutions, des mesures à prendre. C'est souvent la partie la plus attendue par le destinataire.
\item \textbf{La formule de courtoisie administrative} : une phrase de politesse adaptée au destinataire.
\end{enumerate}

\begin{methode}[Rédiger la conclusion d'un rapport]
\begin{enumerate}
\item Rédige d'abord une phrase de bilan qui reprend l'idée principale du corps : «~Cette visite a permis aux élèves de comprendre concrètement les étapes de la production d'électricité.~»
\item Ajoute, si le contexte s'y prête, une ou deux suggestions formulées avec des verbes comme \emph{suggérer, recommander, proposer, souhaiter} : «~Nous suggérons que ce type de visite soit renouvelé chaque année.~»
\item Termine par la formule de courtoisie : «~Nous vous prions d'agréer, Monsieur le Proviseur, l'expression de notre respectueuse considération.~»
\item Vérifie qu'aucun fait nouveau n'apparaît dans la conclusion.
\end{enumerate}
\end{methode}

\begin{attention}[Piège fréquent : les faits nouveaux dans la conclusion]
La conclusion ne doit contenir \textbf{aucun fait nouveau}. Tout ce qui s'y trouve doit avoir été annoncé et développé dans le corps. Si tu découvres en rédigeant ta conclusion une information importante (un chiffre oublié, un témoignage), ne la glisse pas à la fin : remonte-la dans le corps du rapport, à sa place logique. Autre erreur à éviter : transformer la conclusion en répétition mot à mot du corps. La conclusion \emph{résume}, elle ne \emph{recopie} pas.
\end{attention}

\begin{exemplebox}[Une conclusion modèle]
\emph{En définitive, cette visite a été très enrichissante : les élèves ont pu observer directement le fonctionnement des turbines et comprendre le rôle stratégique de la centrale de Songloulou dans l'approvisionnement en électricité du pays. Nous suggérons au chef d'établissement de pérenniser cette sortie pédagogique et d'étendre l'expérience aux classes de Terminale. Nous vous prions d'agréer, Monsieur le Proviseur, l'expression de notre respectueuse considération.}
\end{exemplebox}

\courssub{Les articulations logiques entre les parties}

Un rapport n'est pas une juxtaposition de paragraphes : les idées doivent être \emph{enchaînées}. Ce sont les \cle{connecteurs} qui assurent cette liaison.

\begin{aretenir}[Les connecteurs du rapport]
\begin{itemize}
\item \textbf{Connecteurs chronologiques} (ordre des faits) : \emph{d'abord, ensuite, puis, après, enfin, au terme de, dès notre arrivée}.
\item \textbf{Connecteurs logiques} (cause et conséquence) : \emph{par conséquent, c'est pourquoi, ainsi, de ce fait, en effet, cependant, néanmoins}.
\item \textbf{Connecteurs d'annonce} (passage d'une partie à l'autre) : \emph{en premier lieu, en second lieu, quant à, en ce qui concerne}.
\end{itemize}
\end{aretenir}

Exemple d'enchaînement : «~\emph{D'abord}, le guide nous a présenté l'historique de la centrale. \emph{Ensuite}, nous avons visité la salle des machines. \emph{Enfin}, nous nous sommes rendus au barrage. Les installations étant vétustes, \emph{par conséquent} des travaux de modernisation s'imposent. \emph{C'est pourquoi} nous recommandons...~»

\courssub{Analyse d'un rapport modèle et atelier de production}

\begin{exoresolu}[Coloriage et justification du découpage d'un rapport]
Énoncé. Voici un rapport abrégé. Repère ses trois parties (introduction, corps, conclusion) et justifie le découpage.

\emph{(1) Sur instruction du chef d'établissement en date du 3 février 2026, j'ai procédé le 5 février 2026 à l'inspection de l'atelier d'électrotechnique du lycée technique de Bafoussam, à la suite de la panne signalée par le professeur de travaux pratiques. (2) J'ai d'abord constaté que trois postes de travail sur douze sont hors service. Ensuite, l'inventaire du matériel a révélé l'absence de cinq multimètres. Enfin, le tableau électrique général présente des câbles dénudés dangereux. Le préjudice est estimé à 1 200 000 FCFA. (3) En somme, l'atelier nécessite une remise en état urgente. Je suggère l'achat de cinq multimètres et la réfection du tableau électrique. Je vous prie d'agréer, Monsieur le Proviseur, l'expression de ma haute considération.}
\tcblower
Solution détaillée.
\begin{itemize}
\item \textbf{Introduction = phrase (1)}. Justification : elle répond aux cinq questions — \emph{qui} (l'auteur, sur instruction du chef d'établissement), \emph{quoi} (une inspection), \emph{quand} (le 5 février 2026), \emph{où} (atelier d'électrotechnique du lycée technique de Bafoussam), \emph{pourquoi} (panne signalée).
\item \textbf{Corps = phrase (2)}. Justification : il présente les faits en \emph{ordre thématique} (postes hors service, matériel manquant, danger électrique), avec des \emph{données chiffrées} (3 postes sur 12, 5 multimètres, 1 200 000 FCFA) et des connecteurs chronologiques (\emph{d'abord, ensuite, enfin}).
\item \textbf{Conclusion = phrase (3)}. Justification : elle contient un \emph{bilan} («~remise en état urgente~»), des \emph{suggestions} (achat de multimètres, réfection du tableau) et une \emph{formule de courtoisie}. Aucun fait nouveau n'y apparaît : tout renvoie au corps.
\end{itemize}
\end{exoresolu}

\begin{exercice}[Atelier : de notes désordonnées à l'introduction et à la conclusion]
Voici les notes prises en désordre par un élève après une visite d'entreprise :

\emph{«~départ du lycée 7 h 30 / entreprise SOCAPALM de Dibombari / objectif : voir la transformation du palmier à huile / mardi 21 avril 2026 / classe de Première G1, 48 élèves, 3 enseignants / ordre de mission n° 08/2026 du proviseur / retour 16 h / les élèves veulent recommander la visite aux autres classes / visite très instructive / accueil par le directeur adjoint.~»}

Travail à faire : (1) rédige l'introduction du rapport en répondant aux questions qui ? quoi ? quand ? où ? pourquoi ? ; (2) rédige la conclusion (bilan + suggestion + formule de courtoisie), sans ajouter aucun fait nouveau.
\end{exercice}

\begin{corrige}[Exercice 1]
\textbf{Introduction possible} : «~Conformément à l'ordre de mission n° 08/2026 du proviseur, les 48 élèves de la classe de Première G1, accompagnés de trois enseignants, ont effectué le mardi 21 avril 2026 une visite d'étude à l'entreprise SOCAPALM de Dibombari. Le départ du lycée a eu lieu à 7 h 30 et le retour à 16 h. L'objectif de cette sortie était d'observer les étapes de la transformation du palmier à huile. Le présent rapport rend compte du déroulement de cette visite.~»

\textbf{Conclusion possible} : «~En somme, cette visite a été très instructive : les élèves ont découvert concrètement le fonctionnement d'une unité industrielle de transformation. Nous recommandons que cette expérience soit étendue aux autres classes de l'établissement. Nous vous prions d'agréer, Monsieur le Proviseur, l'expression de notre respectueuse considération.~»

Remarque : les éléments «~accueil par le directeur adjoint~» et les détails de la visite appartiennent au \emph{corps} du rapport, non à l'introduction ni à la conclusion.
\end{corrige}

\begin{savaistu}[Le rapport administratif au Cameroun]
Dans la fonction publique camerounaise, le rapport suit un formalisme strict défini par les instructions générales du Premier Ministre. Tout agent public doit savoir rédiger un \emph{rapport d'activité}, un \emph{rapport de mission} ou un \emph{rapport d'incident}. Le respect de la structure interne (introduction, corps, conclusion) et l'usage du lexique spécialisé (\"vu\", \"considérant\", \"arrête\") sont évalués aux concours professionnels et au Probatoire technique.
\end{savaistu}

\begin{aretenir}[L'essentiel de la section]
\begin{itemize}
\item La structure interne du rapport comporte trois parties : \cle{introduction} (contexte : qui, quoi, quand, où, pourquoi), \cle{corps} (faits ordonnés chronologiquement ou thématiquement, avec chiffres et témoignages), \cle{conclusion} (bilan + suggestions + courtoisie).
\item Le corps se découpe en parties numérotées avec \cle{intertitres}.
\item Les \cle{connecteurs} chronologiques (\emph{d'abord, ensuite, enfin}) et logiques (\emph{par conséquent, c'est pourquoi}) enchaînent les idées.
\item Jamais de fait nouveau dans la conclusion.
\end{itemize}
\end{aretenir}

\coursec{Le texte explicatif au service du rapport}

Dans la section précédente, nous avons appris à organiser un rapport selon sa structure interne : une introduction qui annonce le contexte et l'objectif, un corps qui présente les faits et les analyses, une conclusion qui propose des suites. Mais une bonne architecture ne suffit pas : encore faut-il que chaque paragraphe du corps du rapport soit \emph{compréhensible}. Un lecteur — le proviseur, le chef d'atelier, l'inspecteur — doit pouvoir saisir \emph{pourquoi} un incident s'est produit et \emph{comment} une machine fonctionne ou a cessé de fonctionner. C'est précisément le rôle du \cle{texte explicatif}, que nous étudions maintenant comme l'outil d'écriture fondamental du corps du rapport.

\courssub{Qu'est-ce qu'expliquer ?}

Commençons par une situation simple. Le groupe électrogène du lycée s'arrête en plein cours. Deux réactions sont possibles :

\begin{itemize}
\item \emph{Constater} : « Le groupe électrogène s'est arrêté à 10 h 15. »
\item \emph{Expliquer} : « Le groupe électrogène s'est arrêté à 10 h 15 \textbf{parce que} le niveau de gasoil était insuffisant, \textbf{ce qui} a provoqué la désamorçage de la pompe d'alimentation. »
\end{itemize}

La première phrase informe ; la seconde rend le phénomène \emph{intelligible}. Expliquer, c'est donc répondre aux questions \emph{pourquoi ?} et \emph{comment ?}, et non seulement à la question \emph{quoi ?}.

\begin{definition}[Le texte explicatif]
Un \cle{texte explicatif} est un texte qui rend compte des \textbf{causes}, du \textbf{fonctionnement} ou du \textbf{sens} d'un fait, d'un phénomène ou d'un objet. Il vise à rendre intelligible ce qui, sans lui, resterait obscur pour le lecteur. Il repose sur un enchaînement logique : chaque affirmation est justifiée, illustrée ou précisée par la suivante.
\end{definition}

Dans un rapport technique ou administratif, le texte explicatif intervient surtout dans le \textbf{corps} : c'est là qu'on explique les causes d'un accident, le fonctionnement d'une installation, les raisons d'un retard ou d'une dépense. Sans explication, le rapport se réduit à une liste de faits que le lecteur ne peut ni comprendre ni exploiter pour décider.

\courssub{Les procédés de l'explication}

Pour expliquer, le rédacteur dispose d'outils précis. Observons d'abord un extrait de rapport rédigé par un élève de Première technique après une panne :

\begin{exemplebox}[Observation : un paragraphe qui explique]
« Le compresseur de l'atelier ne démarre plus. Un \textbf{compresseur}, c'est-à-dire une machine qui comprime l'air pour alimenter les outils pneumatiques, ne peut fonctionner que si sa courroie est tendue. Or la courroie était détendue, \textbf{parce que} son galet tendeur était usé. Cette usure s'explique \textbf{par} l'absence de graissage depuis six mois. \textbf{Par exemple}, le manuel d'entretien impose un graissage mensuel. \textbf{Autrement dit}, la panne résulte d'un défaut d'entretien, et non d'un vice de fabrication. »
\end{exemplebox}

Ce paragraphe mobilise, sans le nommer, les cinq grands procédés de l'explication. Dégageons-les rigoureusement.

\begin{propriete}[Les cinq procédés majeurs de l'explication]
\begin{enumerate}
\item \textbf{Définir} : donner le sens précis d'un terme (« un compresseur est une machine qui... »).
\item \textbf{Exemplifier} : illustrer une affirmation générale par un cas concret (« par exemple, le manuel impose... »).
\item \textbf{Comparer} : rapprocher ou opposer deux réalités pour éclairer l'une par l'autre (« à la différence d'un moteur essence, un moteur diesel... »).
\item \textbf{Expliquer par les causes et les conséquences} : établir la chaîne causale (« la courroie était détendue \emph{parce que} le galet était usé ; \emph{par conséquent} le compresseur ne démarre plus »).
\item \textbf{Reformuler} : dire autrement pour clarifier (« autrement dit, c'est-à-dire »).
\end{enumerate}
Ces procédés se combinent : un bon paragraphe explicatif en utilise généralement plusieurs.
\end{propriete}

\begin{attention}[Erreur fréquente : empiler les faits sans liens logiques]
Dans le corps du rapport, beaucoup d'élèves juxtaposent les constatations : « La courroie était détendue. Le galet était usé. Le graissage n'était pas fait. Le compresseur ne démarre plus. » Le lecteur doit alors deviner seul les relations entre les faits : le rapport perd sa valeur explicative. \textbf{Règle d'or} : entre deux faits liés, il doit toujours y avoir un connecteur ou une reprise qui explicite le lien (cause, conséquence, précision).
\end{attention}

\courssub{Les connecteurs de l'explication}

Les connecteurs sont les articulations du raisonnement. Ils signalent au lecteur la nature du lien logique entre deux énoncés. Il faut les classer par \textbf{fonction} pour les choisir correctement.

\begin{popfigure}
\begin{center}
\begin{tikzpicture}[
  box/.style={draw=popblue, fill=popblueL, rounded corners=2pt, align=left, text width=4.6cm, inner sep=5pt, font=\small},
  head/.style={draw=popdark, fill=popdark, text=white, rounded corners=2pt, align=center, text width=4.6cm, inner sep=4pt, font=\small\bfseries},
  node distance=0.35cm and 0.45cm]
\node[head] (h1) {La cause};
\node[box, below=of h1, align=center] (b1){parce que\\ car\\ puisque\\ en effet\\ c'est pourquoi (cause annoncée)\\ grâce à, à cause de};
\node[head, right=of h1] (h2) {La conséquence};
\node[box, below=of h2, align=center] (b2){donc\\ alors\\ ainsi\\ par conséquent\\ c'est pourquoi\\ de sorte que\\ si bien que};
\node[head, below=1.1cm of b1] (h3) {La reformulation};
\node[box, below=of h3, align=center] (b3){c'est-à-dire\\ autrement dit\\ en d'autres termes\\ soit (suivi d'une précision)};
\node[head, right=of h3] (h4) {L'illustration};
\node[box, below=of h4, align=center] (b4){par exemple\\ notamment\\ ainsi (exemple)\\ comme (comparaison)\\ tel que};
\end{tikzpicture}
\end{center}
\legende{Les principaux connecteurs de l'explication classés par fonction logique. Attention : \emph{en effet} introduit une justification (cause), tandis que \emph{ainsi} peut introduire une conséquence ou un exemple selon le contexte.}
\end{popfigure}

Deux pièges méritent d'être signalés. D'abord, \emph{en effet} n'introduit jamais une conséquence : il justifie ce qui vient d'être dit (« Le moteur a calé ; \textbf{en effet}, le filtre à gasoil était colmaté »). Ensuite, \emph{car} et \emph{parce que} ne sont pas interchangeables partout : \emph{parce que} répond à la question « pourquoi ? » et peut ouvrir une phrase, tandis que \emph{car} reste à l'intérieur de la phrase, dans un style soutenu très apprécié du rapport.

\courssub{Les temps verbaux du texte explicatif}

Le choix des temps n'est pas neutre : il distingue ce qui est \emph{toujours vrai} de ce qui s'est passé \emph{une fois}.

\begin{aretenir}[Temps verbaux et fonctions]
\begin{itemize}
\item Le \textbf{présent de vérité générale} exprime les lois, les fonctionnements permanents : « Un moteur diesel \emph{comprime} l'air avant l'injection du carburant. »
\item Le \textbf{passé composé} rapporte des faits accomplis et datés : « Le 12 mars, la courroie \emph{a cédé}. »
\item L'\textbf{imparfait} décrit le décor, les états durables ou les actions en cours au moment du fait : « Le galet \emph{était} usé ; la machine \emph{vibrait} depuis plusieurs jours. »
\end{itemize}
Dans un rapport, le corps alterne donc : présent pour le fonctionnement normal, imparfait et passé composé pour les faits rapportés.
\end{aretenir}

\courssub{La phrase explicative : clarté et enchaînement}

Une phrase explicative efficace respecte trois exigences : un fait par phrase, un lien logique explicite entre les phrases, et une reprise claire des éléments déjà mentionnés. La \textbf{reprise nominale} (répéter le nom ou le remplacer par un synonyme : « le groupe électrogène... la machine... ») et les \textbf{pronoms anaphoriques} (\emph{il, elle, ce, celle-ci, cela}) évitent les répétitions lourdes tout en maintenant la clarté.

\begin{methode}[Améliorer la clarté d'un paragraphe explicatif]
\begin{enumerate}
\item \textbf{Isoler les faits} : une phrase = un fait ou une idée. Supprimer les phrases à trois ou quatre verbes conjugués.
\item \textbf{Ordonner} : aller du général au particulier, ou suivre la chaîne chronologique et causale.
\item \textbf{Relier} : insérer entre les phrases les connecteurs adaptés (cause, conséquence, reformulation, exemple).
\item \textbf{Reprendre} : remplacer les répétitions par des pronoms anaphoriques ou des reprises nominales, en vérifiant qu'on ne peut pas hésiter sur ce qu'ils désignent.
\item \textbf{Relire à voix haute} : si l'on doit relire une phrase deux fois pour la comprendre, la réécrire.
\end{enumerate}
\end{methode}

\courssub{Application : d'un paragraphe confus à un paragraphe explicatif}

\begin{exoresolu}[Transformer un paragraphe de rapport d'accident en atelier]
\textbf{Énoncé.} Voici le paragraphe rédigé à chaud par un élève après un accident survenu à la scie circulaire de l'atelier de bois : « L'élève s'est blessé à la main. La lame n'avait pas de carter de protection. Le carter était démonté. Le professeur avait demandé de le remonter la semaine dernière. Personne ne l'a fait. La scie a été utilisée quand même. Il y a eu un rebond de la pièce de bois. » Réécrivez ce paragraphe en un texte explicatif clair, en utilisant les connecteurs et les temps appropriés.
\tcblower
\textbf{Solution détaillée.} On identifie d'abord la chaîne causale : carter démonté $\rightarrow$ consigne non exécutée $\rightarrow$ scie utilisée sans protection $\rightarrow$ rebond $\rightarrow$ blessure. On rédige ensuite :

« L'élève s'est blessé à la main gauche \textbf{parce que} la scie circulaire \textbf{était} utilisée sans carter de protection. Ce carter, c'est-à-dire le capot métallique qui recouvre la lame, \textbf{avait été} démonté pour une réparation et n'\textbf{avait pas été} remis en place, \textbf{bien que} le professeur d'atelier \textbf{eût demandé} son remontage la semaine précédente. \textbf{Par conséquent}, lorsque la pièce de bois \textbf{a rebondi}, la main de l'élève \textbf{a rencontré} la lame découverte. \textbf{Autrement dit}, l'accident résulte d'une consigne de sécurité non appliquée, et non d'une faute de manipulation. »

On vérifie : un fait par phrase, connecteurs explicites (\emph{parce que, par conséquent, autrement dit}), définition du terme technique (\emph{carter}), imparfait pour les états, passé composé et plus-que-parfait pour les faits.
\end{exoresolu}

\begin{exemplebox}[Paragraphe explicatif : panne du groupe électrogène (lycée de Garoua)]
« Le groupe électrogène du lycée technique de Garoua s'est arrêté brutalement le mardi 14 novembre à 9 h 40, en pleine évaluation. Cet arrêt s'explique par la surchauffe du moteur : en effet, le radiateur était obstrué par la poussière accumulée pendant la saison sèche, période où le vent d'harmattan dépose une couche de sable sur toutes les grilles d'aération. Or un moteur mal refroidi voit sa température monter jusqu'au déclenchement de la sécurité thermique, c'est-à-dire un dispositif qui coupe automatiquement l'alimentation pour éviter la casse. Par conséquent, le groupe s'est mis en sécurité et l'établissement est resté sans courant pendant deux heures. Autrement dit, la panne ne vient pas d'une défaillance du matériel, mais d'un entretien insuffisant : par exemple, le nettoyage hebdomadaire du radiateur, prévu par le carnet d'entretien, n'avait pas été effectué depuis un mois. »
\end{exemplebox}

\courssub{Lien avec le Probatoire et le Baccalauréat technique}

À l'examen, la compétence explicative est évaluée de plusieurs façons : dans le \textbf{commentaire} et l'\textbf{explication de texte}, il faut dégager les procédés par lesquels un auteur rend son propos intelligible (définitions, exemples, connecteurs) ; dans la \textbf{dissertation} et l'\textbf{expression écrite}, c'est à vous d'expliquer clairement vos idées, avec un fait par phrase et des liens logiques visibles. Le rapport, lui, peut être demandé comme sujet de production d'écrits : un corps de rapport sans explication claire est systématiquement pénalisé.

\begin{savaistu}[Pourquoi les notices et rapports techniques privilégient-ils le présent ?]
Dans l'industrie, les notices de machines et les rapports d'audit sont rédigés au présent de vérité générale : « la pompe aspire le gasoil », « le disjoncteur coupe le circuit ». Ce présent, dit \emph{intemporel}, permet à n'importe quel technicien, à n'importe quelle date, de comprendre le fonctionnement décrit. C'est la même logique qui fait écrire au présent les consignes des affiches de sécurité que nous étudierons dans la section 6 : « Porter obligatoirement les lunettes de protection. »
\end{savaistu}

\begin{exercice}[À vous de jouer]
Le réfrigérateur du laboratoire de biologie ne refroidit plus. Voici les faits dans le désordre : la porte est restée ouverte toute la nuit ; le givre a envahi l'évaporateur ; le moteur a forcé ; le thermostat s'est déréglé ; les vaccins stockés ont été perdus. Rédigez un paragraphe explicatif de huit à dix lignes pour le corps du rapport adressé au proviseur, en utilisant au moins quatre connecteurs différents, une définition et un exemple, et en respectant les temps du texte explicatif.
\end{exercice}

\begin{corrige}[Exercice]
Éléments de correction attendus : la chaîne causale doit être ordonnée (porte ouverte $\rightarrow$ givre $\rightarrow$ moteur qui force $\rightarrow$ thermostat déréglé $\rightarrow$ perte des vaccins) ; au moins quatre connecteurs de fonctions différentes (\emph{parce que, par conséquent, c'est-à-dire, par exemple}) ; une définition d'un terme technique (\emph{l'évaporateur, c'est-à-dire la pièce où le froid se produit}) ; imparfait pour les états (« la porte était restée ouverte »), passé composé pour les faits (« le givre a envahi », « les vaccins ont été perdus »), présent de vérité générale pour le fonctionnement (« un moteur qui force surchauffe »). Toute juxtaposition de faits sans lien logique doit être sanctionnée dans l'auto-évaluation.
\end{corrige}

\coursec{Lexique commun et lexique spécialisé}

Dans la section précédente, nous avons vu que le rapport repose sur le \cle{texte explicatif} : il faut expliquer clairement les faits, les procédures et les résultats. Mais expliquer, c'est d'abord \emph{choisir ses mots}. Or, dans un rapport d'atelier ou de stage, deux types de mots se côtoient : les mots que tout le monde comprend et les mots propres au métier. Savoir distinguer le \cle{lexique commun} du \cle{lexique spécialisé}, et savoir passer de l'un à l'autre selon le destinataire, est une compétence essentielle du technicien qui rédige. C'est l'objet de cette section.

\courssub{Le lexique commun : les mots de tous}

\begin{definition}[Lexique commun]
Le \cle{lexique commun} est l'ensemble des mots compris et utilisés par tous les locuteurs d'une langue, quel que soit leur métier ou leur niveau de formation. Ces mots désignent les réalités et les actions de la vie quotidienne : \emph{table, écrire, machine, eau, réparer, mesurer, casser}.
\end{definition}

L'intuition est simple : si vous dites à votre voisin « la machine est en panne », il comprend immédiatement, même s'il n'est pas mécanicien. Les mots \emph{machine} et \emph{panne} appartiennent au lexique commun. Ils ne demandent aucune connaissance particulière.

\begin{exemplebox}[Phrases en lexique commun]
\begin{itemize}
\item « L'appareil ne fonctionne plus. »
\item « Nous avons nettoyé la pièce et remplacé la partie cassée. »
\item « Le chef d'atelier a vérifié le travail. »
\end{itemize}
Ces phrases sont correctes mais imprécises : quelle pièce ? quelle partie ? quel appareil ? Pour un rapport technique, ce flou est un défaut.
\end{exemplebox}

\courssub{Le lexique spécialisé : les mots du métier}

\begin{definition}[Lexique spécialisé]
Le \cle{lexique spécialisé} est l'ensemble des termes propres à un domaine technique, scientifique ou professionnel. Ces mots ne sont pleinement compris que par les spécialistes du domaine : \emph{vilebrequin} (mécanique), \emph{court-circuit} (électrotechnique), \emph{sinusite} (médecine), \emph{grand livre} (comptabilité).
\end{definition}

Le lexique spécialisé a deux qualités majeures pour le rédacteur d'un rapport :
\begin{itemize}
\item la \cle{précision} : un seul mot désigne exactement une pièce, une opération ou un phénomène, sans ambiguïté ;
\item l'\cle{économie} : dire « vilebrequin » évite la longue périphrase « axe moteur qui transforme le mouvement de va-et-vient des pistons en mouvement de rotation ».
\end{itemize}

\begin{exemplebox}[Champs lexicaux spécialisés par filière technique]
\begin{center}
\begin{tabularx}{\linewidth}{|l|X|}\hline
\rowcolor{popblueL} \textbf{Filière} & \textbf{Exemples de termes spécialisés} \\ \hline
Mécanique & vilebrequin, piston, embrayage, culasse, jeu de soupapes \\ \hline
Électrotechnique & court-circuit, disjoncteur, tension, intensité, relais, bobine \\ \hline
Bâtiment & chaînage, linteau, fondation, enduit, ferraillage, pignon \\ \hline
Secrétariat & classement chronologique, bordereau, courrier arrivée, parapheur \\ \hline
Comptabilité & bilan, journal, grand livre, amortissement, créance, TVA \\ \hline
\end{tabularx}
\end{center}
\end{exemplebox}

\begin{attention}[Le jargon non expliqué]
L'erreur la plus fréquente dans un rapport de stage est d'utiliser un \cle{jargon} technique sans le définir devant un destinataire non spécialiste. Écrire au directeur d'école « le relais thermique a déclenché à cause d'un déséquilibre de phases » sans aucune explication, c'est produire un rapport inutile : le lecteur ne comprend ni la cause ni la gravité de la panne. Règle d'or : \textbf{tout terme spécialisé adressé à un non-spécialiste doit être défini ou remplacé par son équivalent en lexique commun}.
\end{attention}

\courssub{Observation : repérer les deux lexiques dans des textes réels}

Avant de définir une méthode, observons. Voici deux extraits courts, tels qu'on peut les lire dans une notice d'appareil et dans un rapport d'atelier.

\begin{exemplebox}[Extrait 1 — notice d'un moteur électrique]
« Avant toute intervention, couper l'alimentation au \textbf{disjoncteur} et vérifier l'absence de \textbf{tension} aux \textbf{bornes} du moteur. »
\end{exemplebox}

\begin{exemplebox}[Extrait 2 — rapport d'atelier (mécanique)]
« Le moteur a été ouvert. La \textbf{culasse} présentait des dépôts de calamine. Les \textbf{segments} étaient usés et ont été remplacés. Après remontage, le moteur a été essayé : il tourne bien. »
\end{exemplebox}

\begin{experience}[Démarche d'observation : trier les mots des deux extraits]
\textbf{Protocole.} Soulignez tous les noms et verbes des deux extraits, puis classez-les selon un critère simple : « une personne sans formation technique comprend-elle ce mot sans aide ? »\par
\textbf{Observations.}
\begin{itemize}
\item Lexique commun : \emph{couper, vérifier, absence, moteur, ouvrir, remplacer, remonter, essayer, tourner}.
\item Lexique spécialisé : \emph{disjoncteur, tension, bornes, culasse, calamine, segments}.
\end{itemize}
\textbf{Interprétation.} Les deux lexiques cohabitent dans le même texte : le lexique commun porte les actions générales (ouvrir, remplacer, essayer), le lexique spécialisé nomme les pièces et les phénomènes avec précision. Un bon rapport technique mélange les deux, mais dose le lexique spécialisé selon le lecteur.
\end{experience}

\courssub{La zone d'intersection : des mots qui changent de statut}

Certains mots, nés dans un domaine spécialisé, finissent par passer dans l'usage courant. \emph{Virus} vient de la médecine, \emph{batterie} de l'électrotechnique, \emph{télécharger} de l'informatique : aujourd'hui, tout le monde les emploie. Ces mots appartiennent à la fois au lexique spécialisé (ils gardent un sens technique précis) et au lexique commun (ils sont compris de tous).

\begin{popfigure}
\begin{center}
\begin{tikzpicture}
\draw[fill=popblue!25,draw=popblue,thick] (-1.3,0) circle (2.4);
\draw[fill=poporange!25,draw=poporange,thick] (1.3,0) circle (2.4);
\node[popblue,font=\bfseries] at (-2.9,1.7) {Lexique commun};
\node[poporange,font=\bfseries] at (2.9,1.7) {Lexique spécialisé};
\node[align=center,font=\small] at (-2.3,0.4) {table\\ écrire\\ machine\\ réparer};
\node[align=center,font=\small] at (2.3,0.4) {vilebrequin\\ disjoncteur\\ chaînage\\ amortissement};
\node[align=center,font=\small\itshape] at (0,0.2) {virus\\ batterie\\ télécharger\\ ordinateur};
\node[align=center,font=\footnotesize,popmuted] at (0,-1.1) {termes passés\\ dans l'usage courant};
\end{tikzpicture}
\end{center}
\legende{Diagramme de Venn : le lexique commun et le lexique spécialisé se recoupent. La zone centrale contient les termes techniques devenus courants.}
\end{popfigure}

\begin{savaistu}[Des mots techniques devenus courants]
Le vocabulaire n'est pas figé : il circule entre les métiers et la vie quotidienne. Le mot \emph{ordinateur} a été inventé en 1955 à la demande d'IBM France pour nommer une machine réservée aux spécialistes ; soixante ans plus tard, il est dans toutes les bouches. De même, \emph{télécharger}, \emph{virus}, \emph{batterie} ou \emph{stress} étaient des termes de spécialistes avant de devenir des mots de tous les jours. Au Cameroun, des mots comme \emph{forage} (hydraulique) ou \emph{groupe électrogène} sont passés dans le langage courant des quartiers. Conséquence pratique : dans un rapport, ces mots n'ont pas besoin d'être définis, même pour un lecteur non spécialiste.
\end{savaistu}

\courssub{Adapter le lexique au destinataire du rapport}

Le choix du lexique dépend d'une question unique : \textbf{qui va lire mon rapport ?}

\begin{exemplebox}[Le même fait, deux rédactions]
\begin{itemize}
\item \textbf{Pour un technicien (lexique spécialisé) :} « Le moteur s'est arrêté à la suite d'un court-circuit dans l'enroulement du stator. »
\item \textbf{Pour un non-spécialiste (lexique commun + définition) :} « Le moteur s'est arrêté à la suite d'un court-circuit, c'est-à-dire un contact accidentel entre deux fils conducteurs, qui a brûlé une partie interne du moteur. »
\end{itemize}
Le fait rapporté est identique ; seul le niveau de langue change. La deuxième version \emph{vulgarise} : elle garde le terme exact mais l'explique immédiatement en lexique commun.
\end{exemplebox}

\begin{methode}[Adapter le niveau de lexique au destinataire]
\begin{enumerate}
\item \textbf{Identifier le destinataire} : technicien de la même filière, chef d'entreprise, client, autorité administrative ?
\item \textbf{Repérer les termes spécialisés} du brouillon : souligner tous les mots qu'un lecteur extérieur au métier ne comprendrait pas.
\item \textbf{Décider pour chaque terme} : le garder tel quel (destinataire spécialiste), le définir entre parenthèses ou par « c'est-à-dire » (destinataire mixte), ou le remplacer par un équivalent en lexique commun (destinataire profane).
\item \textbf{Ajouter un glossaire} en fin de rapport si les termes spécialisés sont nombreux : liste alphabétique des termes suivis de leur définition courte.
\item \textbf{Relire en se mettant à la place du lecteur} : toute phrase qui exige une connaissance du métier non fournie dans le texte doit être réécrite.
\end{enumerate}
\end{methode}

\begin{aretenir}[Définir un terme spécialisé dans un rapport]
Trois procédés simples permettent de définir un terme technique sans alourdir le texte :
\begin{itemize}
\item la \cle{parenthèse explicative} : « le disjoncteur (appareil qui coupe automatiquement le courant en cas de défaut) » ;
\item la \cle{reformulation} introduite par \emph{c'est-à-dire}, \emph{autrement dit} : « un court-circuit, c'est-à-dire un contact accidentel entre deux conducteurs » ;
\item le \cle{glossaire} placé à la fin du rapport pour les documents longs.
\end{itemize}
\end{aretenir}

\courssub{Exercice résolu : classer et réécrire}

\begin{exoresolu}[Atelier de classement lexical]
\textbf{Énoncé.} Voici vingt mots relevés dans un rapport d'atelier d'électrotechnique : \emph{moteur, disjoncteur, ouvrir, tension, nettoyer, relais, panne, bobine, mesurer, isolant, casser, borne, huile, court-circuit, vérifier, stator, remplacer, câble, serrer, ampèremètre}.
\begin{enumerate}
\item Classez ces mots en deux colonnes : lexique commun / lexique spécialisé.
\item Réécrivez la phrase suivante pour un directeur non technicien : « Le relais thermique a déclenché en raison d'une surintensité provoquée par le blocage du rotor. »
\end{enumerate}
\tcblower
\textbf{Solution détaillée.}
\begin{enumerate}
\item Classement :
\begin{center}
\begin{tabularx}{\linewidth}{|X|X|}\hline
\rowcolor{popblueL} \textbf{Lexique commun} & \textbf{Lexique spécialisé} \\ \hline
moteur, ouvrir, nettoyer, panne, mesurer, casser, huile, vérifier, remplacer, câble, serrer & disjoncteur, tension, relais, bobine, isolant, borne, court-circuit, stator, ampèremètre \\ \hline
\end{tabularx}
\end{center}
Remarque : \emph{câble} et \emph{moteur} sont acceptés en lexique commun car ils sont passés dans l'usage courant, même s'ils ont un sens technique précis (zone d'intersection du diagramme).
\item Réécriture pour un non-spécialiste : « Un appareil de sécurité (le relais thermique) a automatiquement coupé le moteur parce que le courant est devenu trop fort : la partie tournante du moteur était bloquée. »
\end{enumerate}
La réécriture conserve le terme technique entre parenthèses (pour la traçabilité) mais explique chaque notion en lexique commun.
\end{exoresolu}

\begin{exercice}[Réécrire pour deux destinataires]
Voici un passage de rapport de chantier (filière bâtiment) : « Le ferraillage des linteaux a été contrôlé avant le coulage du béton. Le chaînage vertical est conforme aux plans. »
\begin{enumerate}
\item Soulignez les termes du lexique spécialisé.
\item Réécrivez ce passage pour le chef de chantier (version technique).
\item Réécrivez-le pour le propriétaire de la maison, qui n'est pas du métier (version vulgarisée, avec définitions).
\end{enumerate}
\end{exercice}

\begin{corrige}[Exercice : réécrire pour deux destinataires]
\begin{enumerate}
\item Termes spécialisés : \emph{ferraillage, linteaux, coulage, béton (au sens technique), chaînage vertical}.
\item Version technique (pour le chef de chantier) : le passage peut rester tel quel, car le destinataire maîtrise le lexique du bâtiment. On peut ajouter une précision : « Le ferraillage des linteaux (section et espacement des barres) a été contrôlé avant le coulage du béton ; le chaînage vertical est conforme aux plans d'exécution. »
\item Version vulgarisée (pour le propriétaire) : « Avant de couler le béton, nous avons vérifié l'armature en fer placée à l'intérieur des linteaux, c'est-à-dire les poutres situées au-dessus des portes et des fenêtres. Les renforts verticaux qui solidarisent les murs (le chaînage) sont conformes aux plans. »
\end{enumerate}
\end{corrige}

\begin{aretenir}[L'essentiel de la section]
\begin{itemize}
\item Le \cle{lexique commun} est compris de tous ; le \cle{lexique spécialisé} est propre à un métier et garantit la précision.
\item Certains termes techniques (\emph{virus, batterie, télécharger}) passent avec le temps dans le lexique commun.
\item Dans un rapport, le niveau de lexique s'adapte au \cle{destinataire} : précision pour le technicien, vulgarisation pour le non-spécialiste.
\item Tout terme spécialisé adressé à un profane doit être défini : parenthèse explicative, reformulation ou glossaire.
\end{itemize}
Cette maîtrise du vocabulaire vous servira directement dans la section suivante, consacrée à la lecture des textes fonctionnels et des affiches de sécurité, où le choix des mots conditionne la compréhension immédiate du message.
\end{aretenir}

\coursec{Lecture de textes fonctionnels et affiches de sécurité}

Dans la section précédente, nous avons distingué le \cle{lexique commun} du \cle{lexique spécialisé} : nous avons vu qu'un rapport technique emploie un vocabulaire précis, adapté au domaine professionnel concerné. Mais avant même de rédiger un rapport, il faut savoir \emph{lire} les documents qui circulent dans un atelier, un chantier ou une entreprise. Or ces documents ne sont pas des textes littéraires : ce sont des \cle{textes fonctionnels}, parmi lesquels les affiches de sécurité occupent une place essentielle. C'est à leur lecture méthodique que cette section est consacrée, car bien souvent, le fait à rapporter dans un rapport est précisément le respect ou la violation d'une consigne affichée.

\courssub{Qu'est-ce qu'un texte fonctionnel ?}

\textbf{Intuition.} Le matin, en entrant à l'atelier de chaudronnerie, tu croises une affiche « Port des lunettes obligatoire », une notice d'utilisation de la perceuse à colonne, un formulaire de sortie de matériel, un panneau « Défense de fumer ». Aucun de ces textes ne cherche à te faire rêver ou à te raconter une histoire : tous visent une \emph{action} ou une \emph{information immédiate}. Ce sont des textes fonctionnels.

\begin{definition}[Le texte fonctionnel]
Un \cle{texte fonctionnel} (ou texte utilitaire) est un texte de la vie quotidienne et professionnelle qui vise une \textbf{action} ou une \textbf{information immédiate}. Il ne cherche ni à divertir ni à émouvoir : il sert à \emph{faire} ou à \emph{faire faire}. En font partie : la \textbf{notice} d'utilisation, la \textbf{consigne} de sécurité, l'\textbf{affiche}, le \textbf{formulaire}, le \textbf{mode d'emploi}, le \textbf{règlement intérieur}, l'\textbf{horaire} de transport.
\end{definition}

\textbf{Caractéristiques des textes fonctionnels.} On reconnaît un texte fonctionnel à plusieurs traits réunis :

\begin{itemize}
\item la \textbf{brièveté} : le message doit être lu et compris en quelques secondes ;
\item l'emploi de l'\textbf{impératif} (\og Éteignez les machines après usage \fg{}) ou de l'\textbf{infinitif} (\og Ne pas fumer \fg{}, \og Porter un casque \fg{}), modes de l'ordre et de la consigne ;
\item une \textbf{typographie} soignée : caractères gras, majuscules, taille importante, encadrés ;
\item le recours à des \textbf{pictogrammes} et à des \textbf{couleurs codées} qui rendent le message compréhensible même par un lecteur pressé ou peu francophone ;
\item un \textbf{vocabulaire précis}, souvent spécialisé (nous retrouvons ici le lexique technique étudié à la section précédente).
\end{itemize}

\begin{attention}[Ne pas confondre texte fonctionnel et texte littéraire]
À l'examen, ne décris jamais une affiche ou une notice comme on analyse un poème : pas de recherche de sentiments de l'auteur, pas de figures de style ornementales. La question à poser est : \emph{que doit faire ou savoir le lecteur après avoir lu ce texte ?} Un texte fonctionnel réussi est un texte qui déclenche la bonne action.
\end{attention}

\begin{exemplebox}[Panorama des textes fonctionnels en milieu technique]
\begin{center}
\begin{tabularx}{\linewidth}{|l|X|X|}
\hline
\rowcolor{popblueL} \textbf{Type} & \textbf{Fonction principale} & \textbf{Exemple en atelier} \\
\hline
Notice / Mode d'emploi & Guider l'utilisation d'un matériel & Notice de la perceuse à colonne \\
\hline
Consigne / Affiche de sécurité & Imposer un comportement ou signaler un danger & \og Port des gants obligatoire \fg{} (cercle bleu) \\
\hline
Formulaire & Recueillir des informations standardisées & Fiche de sortie de matériel \\
\hline
Règlement intérieur & Fixer les règles collectives & Règlement de l'atelier de soudure \\
\hline
Étiquetage / Signalétique & Identifier, localiser, orienter & \og Arrivée eau froide \fg{}, \og Issue de secours \fg{} \\
\hline
\end{tabularx}
\end{center}
\end{exemplebox}

\courssub{Les affiches d'indications sécuritaires : un texte iconique}

\textbf{Observation.} À l'entrée de l'atelier de mécanique du lycée technique, trois affiches sont visibles : un éclair jaune dans un triangle (\og Danger électrique \fg{}), un cercle bleu avec un personnage coiffé (\og Port du casque obligatoire \fg{}), un rectangle vert avec un personnage qui court vers une porte (\og Issue de secours \fg{}). Remarque : avant même de lire les mots, tu as compris l'essentiel grâce aux \emph{images} et aux \emph{couleurs}. C'est toute la force du \cle{texte iconique}.

\begin{definition}[Le texte iconique]
Un \cle{texte iconique} est un texte dont le message passe en tout ou partie par l'\textbf{image} : dessin, photographie, schéma ou \textbf{pictogramme} (image simplifiée et conventionnelle). Dans une affiche de sécurité, le pictogramme, la couleur et la consigne écrite travaillent ensemble.
\end{definition}

\textbf{Le code couleur des affiches de sécurité.} Les couleurs ne sont pas choisies au hasard : elles forment un code que tout technicien doit connaître par cœur.

\begin{center}
\begin{tabularx}{\linewidth}{|l|X|X|}
\hline
\rowcolor{popblueL} \textbf{Couleur / forme} & \textbf{Signification} & \textbf{Exemples} \\
\hline
Rouge (cercle barré ou carré) & Interdiction ou matériel de lutte contre l'incendie & Défense de fumer ; extincteur \\
\hline
Jaune (triangle) & Danger, avertissement & Danger électrique ; sol glissant \\
\hline
Bleu (cercle) & Obligation & Port du casque ; port des lunettes \\
\hline
Vert (rectangle ou carré) & Secours, issue, sécurité & Issue de secours ; poste de secours \\
\hline
\end{tabularx}
\end{center}

\begin{attention}[Erreur fréquente : négliger le code couleur]
Beaucoup d'élèves décrivent une affiche en ne citant que la phrase écrite (\og Port du casque obligatoire \fg{}) et oublient la couleur et la forme. Or la couleur porte déjà la moitié du message : un cercle \textbf{bleu} signifie une \emph{obligation}, un cercle \textbf{rouge barré} une \emph{interdiction}, un triangle \textbf{jaune} un \emph{danger}, un rectangle \textbf{vert} une \emph{voie de secours}. Dans la lecture comme dans le rapport, cite toujours le pictogramme, la couleur et la consigne.
\end{attention}

\begin{popfigure}
\begin{center}
\begin{tikzpicture}[scale=1]
% Danger électrique : triangle jaune
\begin{scope}[shift={(0,0)}]
\fill[popgold!85] (0,0) -- (2.4,0) -- (1.2,2.1) -- cycle;
\draw[line width=2pt,popdark] (0,0) -- (2.4,0) -- (1.2,2.1) -- cycle;
\draw[line width=1.6pt,popdark,fill=popdark] (1.35,1.55) -- (0.95,0.85) -- (1.2,0.85) -- (1.05,0.35) -- (1.5,1.05) -- (1.25,1.05) -- cycle;
\node[below,align=center] at (1.2,-0.1) {\small \textbf{Danger électrique}\\ \footnotesize (jaune = avertissement)};
\end{scope}
% Port obligatoire : cercle bleu
\begin{scope}[shift={(4.6,1.05)}]
\fill[popblue] (0,0) circle (1.05);
\draw[line width=2pt,popdark] (0,0) circle (1.05);
\fill[white] (0,0.35) circle (0.28);
\draw[line width=2.5pt,white] (0,0.05) -- (0,-0.55);
\draw[line width=2.5pt,white] (-0.4,-0.2) -- (0.4,-0.2);
\draw[line width=2.5pt,white] (0,-0.55) -- (-0.3,-0.85);
\draw[line width=2.5pt,white] (0,-0.55) -- (0.3,-0.85);
\draw[line width=2pt,white] (-0.32,0.42) arc (160:20:0.36);
\node[below,align=center] at (0,-1.2) {\small \textbf{Port du casque obligatoire}\\ \footnotesize (bleu = obligation)};
\end{scope}
% Issue de secours : rectangle vert
\begin{scope}[shift={(8.2,0)}]
\fill[popgreen!80] (0,0) rectangle (2.6,2.1);
\draw[line width=2pt,popdark] (0,0) rectangle (2.6,2.1);
\fill[white] (0.95,1.45) circle (0.2);
\draw[line width=2.2pt,white] (0.95,1.25) -- (0.75,0.7);
\draw[line width=2.2pt,white] (0.95,1.1) -- (1.35,0.95);
\draw[line width=2.2pt,white] (0.75,0.7) -- (0.5,0.25);
\draw[line width=2.2pt,white] (0.75,0.7) -- (1.15,0.3);
\draw[line width=2.2pt,white] (1.9,0.5) rectangle (2.35,1.7);
\draw[line width=2.2pt,white,->] (1.35,1.05) -- (1.85,1.05);
\node[below,align=center] at (1.3,-0.1) {\small \textbf{Issue de secours}\\ \footnotesize (vert = secours)};
\end{scope}
\end{tikzpicture}
\end{center}
\legende{Planche de pictogrammes de sécurité reproduits en schéma : le triangle jaune signale un danger, le cercle bleu impose une obligation, le rectangle vert indique une issue ou un équipement de secours.}
\end{popfigure}

\begin{savaistu}[Une norme internationale]
Les pictogrammes de sécurité sont \textbf{normalisés au niveau international} par la norme \cle{ISO 7010}. Cela signifie qu'un technicien camerounais formé à Yaoundé comprendra immédiatement les affiches de sécurité d'une usine en France, en Côte d'Ivoire ou en Chine, sans connaître la langue du pays. Le pictogramme est une langue universelle au service de la sécurité : c'est l'exemple même du texte fonctionnel réussi.
\end{savaistu}

\courssub{Méthode : décoder une affiche de sécurité}

Lire une affiche de sécurité n'est pas un acte spontané : c'est une démarche qui s'apprend et qui sera utile dans tout atelier, chantier ou bureau.

\begin{methode}[Décoder une affiche de sécurité en quatre étapes]
\begin{enumerate}
\item \textbf{Repérer la couleur et la forme} du panneau : rouge (interdiction ou incendie), jaune (danger), bleu (obligation), vert (secours). C'est le premier niveau du message.
\item \textbf{Identifier le pictogramme} : que représente l'image simplifiée (éclair, flamme, casque, personnage qui court) ? Quel objet ou quelle action évoque-t-elle ?
\item \textbf{Lire la consigne écrite} : relever le mode verbal employé (infinitif : \og Ne pas fumer \fg{} ; impératif : \og Portez vos lunettes \fg{}) et le vocabulaire spécialisé éventuel.
\item \textbf{Formuler le message complet} en une phrase claire : qui est concerné, quelle action est imposée ou interdite, dans quel lieu, pour quel danger.
\end{enumerate}
\end{methode}

\textbf{La relation texte-image : redondance et complémentarité.} Dans une affiche de sécurité, le texte et l'image entretiennent deux types de relations :

\begin{itemize}
\item la \textbf{redondance} : l'image et le texte disent la même chose (le dessin d'une cigarette barrée + \og Défense de fumer \fg{}). La répétition renforce le message et le rend accessible à tous, même aux non-lecteurs ;
\item la \textbf{complémentarité} : l'image et le texte apportent chacun une information que l'autre ne donne pas (le pictogramme montre la flamme, le texte précise : \og En cas d'incendie, empruntez l'escalier B, n'utilisez pas l'ascenseur \fg{}).
\end{itemize}

\begin{exemplebox}[Analyse d'une affiche de chantier]
À l'entrée du chantier du nouveau atelier du lycée technique de Bafoussam, une affiche présente un \textbf{cercle bleu} contenant un personnage coiffé d'un casque, avec la mention \og \textbf{Port du casque obligatoire} \fg{}. Analyse : le cercle bleu signale une \emph{obligation} ; le pictogramme montre l'équipement concerné ; la consigne à l'infinitif (\og port \fg{}) exprime l'ordre de manière impersonnelle ; la relation texte-image est ici \emph{redondante} : chacun, ouvrier ou visiteur, comprend qu'il ne peut entrer sans casque.
\end{exemplebox}

\begin{popfigure}
\begin{center}
\begin{tikzpicture}[scale=1]
% Affiche centrale
\draw[line width=1.5pt,popdark,fill=popcream] (4.2,0.4) rectangle (7.8,4.6);
\fill[popblue] (6,3.1) circle (0.85);
\fill[white] (6,3.35) circle (0.2);
\draw[line width=2pt,white] (6,3.15) -- (6,2.7);
\draw[line width=2pt,white] (5.7,2.95) -- (6.3,2.95);
\draw[line width=2pt,white] (6,2.7) -- (5.75,2.45);
\draw[line width=2pt,white] (6,2.7) -- (6.25,2.45);
\draw[line width=1.8pt,white] (5.72,3.42) arc (160:20:0.32);
\node[align=center,font=\bfseries] at (6,1.6) {PORT DU CASQUE\\ OBLIGATOIRE};
\node[align=center,font=\footnotesize] at (6,0.85) {Chantier du lycée technique};
% Flèches d'analyse
\draw[line width=1.2pt,poporange,->] (4.2,4.3) -- (1.9,4.3);
\node[left,align=right,text=popdark] at (1.8,4.3) {\small \textbf{1. Cadre et titre}\\ \footnotesize lieu et contexte};
\draw[line width=1.2pt,poporange,->] (5.15,3.4) -- (1.9,3.1);
\node[left,align=right,text=popdark] at (1.8,3.1) {\small \textbf{2. Pictogramme}\\ \footnotesize image simplifiée};
\draw[line width=1.2pt,poporange,->] (5.4,1.6) -- (1.9,1.7);
\node[left,align=right,text=popdark] at (1.8,1.7) {\small \textbf{3. Consigne écrite}\\ \footnotesize infinitif ou impératif};
\draw[line width=1.2pt,poporange,->] (7.8,3.9) -- (10.1,3.9);
\node[right,align=left,text=popdark] at (10.2,3.9) {\small \textbf{4. Couleur codée}\\ \footnotesize bleu = obligation};
\draw[line width=1.2pt,poporange,->] (7.8,1.0) -- (10.1,1.0);
\node[right,align=left,text=popdark] at (10.2,1.0) {\small \textbf{5. Relation texte-image}\\ \footnotesize redondance ou complémentarité};
\end{tikzpicture}
\end{center}
\legende{Schéma d'analyse d'une affiche de sécurité : on lit successivement le cadre, le pictogramme, la consigne écrite, la couleur, puis la relation entre texte et image.}
\end{popfigure}

\courssub{Lire des affiches en contexte : atelier, circulation, incendie}

Appliquons la méthode à trois situations concrètes de la vie d'un élève de Première technique.

\textbf{1. Les affiches d'atelier.} Dans un atelier de mécanique, on trouve typiquement : \og Port des lunettes obligatoire \fg{} (cercle bleu, obligation), \og Défense de fumer \fg{} (cercle rouge barré, interdiction), \og Attention, pièces chaudes \fg{} (triangle jaune, danger). Ces affiches emploient le lexique spécialisé du métier (\emph{lunettes de protection}, \emph{écran facial}, \emph{gants de manutention}) : le rapporteur qui citera ces affiches devra reprendre ce vocabulaire exact, sans le paraphraser approximativement.

\textbf{2. Les affiches et panneaux de circulation à Yaoundé.} Sur le boulevard de l'Unité ou à l'entrée du carrefour Nsam, les panneaux \og Sens interdit \fg{} (disque rouge), \og Cédez le passage \fg{} (triangle blanc bordé de rouge), \og Passage pour piétons \fg{} (panneau bleu) obéissent au même code : forme, couleur, pictogramme, brièveté. Le conducteur lit en une fraction de seconde : c'est la preuve que le texte fonctionnel est conçu pour une \emph{lecture immédiate}.

\textbf{3. Les consignes anti-incendie.} Dans les bâtiments scolaires, l'affiche anti-incendie combine souvent plusieurs couleurs : le rouge pour localiser l'extincteur, le vert pour indiquer l'issue de secours, et une consigne écrite détaillée (\og En cas d'incendie : 1. Donnez l'alerte. 2. Évacuez par l'escalier. 3. Rassemblement devant le portail. \fg{}). Ici, la relation texte-image est \emph{complémentaire} : le pictogramme localise, le texte numéroté ordonne les actions.

\begin{exoresolu}[Décodage complet d'une affiche d'atelier]
Énoncé : à l'entrée de l'atelier d'électrotechnique, on lit une affiche composée d'un triangle jaune contenant un éclair noir, avec la mention \og Danger : haute tension. Ne pas ouvrir l'armoire électrique. \fg{} Analyse cette affiche selon la méthode en quatre étapes.
\tcblower
Solution détaillée :
\begin{enumerate}
\item \textbf{Couleur et forme} : triangle jaune, donc \emph{avertissement} : il existe un danger réel dans ce lieu.
\item \textbf{Pictogramme} : un éclair noir, symbole conventionnel de l'électricité ; il désigne la nature précise du danger (le courant électrique).
\item \textbf{Consigne écrite} : deux phrases brèves ; la première nomme le danger (\og haute tension \fg{}, lexique spécialisé du domaine électrique), la seconde emploie l'infinitif négatif \og Ne pas ouvrir \fg{}, mode typique de l'interdiction impersonnelle.
\item \textbf{Message complet} : toute personne présente dans l'atelier est avertie que l'armoire électrique présente un risque d'électrocution et qu'il est interdit de l'ouvrir. La relation texte-image est ici complémentaire : le pictogramme signale l'électricité, le texte précise l'action interdite.
\end{enumerate}
\end{exoresolu}

\courssub{Des affiches de sécurité au rapport : la consigne comme fait à rapporter}

Pourquoi étudier les affiches de sécurité dans une leçon consacrée au rapport ? Parce que, dans la vie professionnelle, un grand nombre de rapports portent précisément sur le \textbf{respect ou la violation des consignes affichées}. Le chef d'atelier, le surveillant, le délégué de classe doivent rapporter des faits : or le fait à rapporter est souvent le suivant — \emph{une consigne affichée n'a pas été respectée}.

Le lien entre l'affiche et le rapport impose trois règles de rédaction :

\begin{itemize}
\item \textbf{citer exactement la consigne} concernée, entre guillemets, sans la déformer (\og Port des lunettes obligatoire \fg{}, et non \og il fallait mettre des lunettes \fg{}) ;
\item \textbf{décrire objectivement le fait observé} : date, heure, lieu, personnes concernées, comportement constaté ;
\item \textbf{rester neutre} : le rapport constate, il n'accuse pas ; on écrit \og trois élèves travaillaient sans lunettes de protection \fg{} et non \og trois élèves indisciplinés se moquaient du règlement \fg{}.
\end{itemize}

\begin{exercice}[Rédaction d'un court rapport]
Dans l'atelier de soudure de ton lycée, tu as observé ce mardi 12 mai, vers 10 h 30, que plusieurs élèves soudaient sans écran facial et qu'un élève fumait près des bouteilles de gaz, alors que deux affiches sont visibles : \og Port de l'écran facial obligatoire \fg{} (cercle bleu) et \og Défense de fumer — Gaz inflammable \fg{} (cercle rouge barré et triangle jaune). Rédige un court rapport (une quinzaine de lignes) adressé au chef d'atelier, en citant les affiches concernées.
\end{exercice}

\begin{corrige}[Exercice : rapport sur le non-respect des consignes]
\textbf{Rapport sur le non-respect des consignes de sécurité à l'atelier de soudure}

Monsieur,

J'ai l'honneur de porter à votre connaissance les faits suivants, observés ce mardi 12 mai, vers 10 h 30, dans l'atelier de soudure.

Lors de la séance de travaux pratiques, j'ai constaté que quatre élèves effectuaient des soudures à l'arc sans écran facial, alors que l'affiche \og Port de l'écran facial obligatoire \fg{} (panneau bleu d'obligation) est apposée à l'entrée de l'atelier. Par ailleurs, un élève fumait à proximité immédiate des bouteilles de gaz, en violation de l'affiche \og Défense de fumer — Gaz inflammable \fg{} (panneau rouge d'interdiction accompagné d'un triangle jaune de danger).

Ces comportements exposent les élèves à des risques graves : brûlures oculaires pour les uns, explosion pour l'ensemble de l'atelier.

Je vous propose de rappeler les consignes au début de chaque séance et de vérifier le port des équipements de protection avant toute mise en route des postes.

Je vous prie d'agréer, Monsieur, l'expression de ma considération respectueuse.

Le délégué de l'atelier.
\end{corrige}

\begin{aretenir}[L'essentiel de la section]
\begin{itemize}
\item Un \cle{texte fonctionnel} vise une action ou une information immédiate : notice, consigne, affiche, formulaire, mode d'emploi.
\item Il se reconnaît à sa brièveté, à l'emploi de l'infinitif ou de l'impératif, à sa typographie et à ses pictogrammes.
\item Le \cle{code couleur} des affiches de sécurité : rouge = interdiction/incendie, jaune = danger, bleu = obligation, vert = secours.
\item On décode une affiche en quatre étapes : couleur, pictogramme, consigne écrite, message complet ; la relation texte-image est redondante ou complémentaire.
\item Dans un rapport, la consigne affichée se cite exactement ; le fait observé se décrit objectivement, avec date, lieu et personnes, sans jugement personnel.
\end{itemize}
\end{aretenir}

Cette compétence de lecture des textes fonctionnels acquise, nous disposons désormais de tous les outils — structure, lexique, observation des faits — pour passer à l'étape finale de l'unité : la \emph{production et l'amélioration d'un rapport complet}, objet de la section suivante.

\coursec{Production et amélioration d'un rapport complet}

Dans la section précédente, nous avons appris à \emph{lire} les textes fonctionnels — notices, consignes, affiches de sécurité — pour en extraire l'information utile. Nous franchissons maintenant l'étape décisive : passer de la \cle{réception} à la \cle{production}. Lire une affiche de sécurité dans un atelier, c'est bien ; mais si un incident survient malgré tout, c'est vous, futur technicien, qui devrez rédiger le rapport. Cette section vous donne la démarche complète, de la situation d'écriture jusqu'à l'amélioration du texte, pour produire un rapport clair, objectif et conforme aux attentes du monde professionnel comme à celles du Probatoire et du Baccalauréat technique.

\courssub{Rappel de la démarche de production d'écrit}

\begin{definition}[La démarche de production]
La \cle{démarche de production d'écrit} est l'ensemble ordonné des étapes par lesquelles un rédacteur transforme une situation de communication en un texte abouti. Pour le rapport, elle comporte cinq étapes : comprendre la \cle{situation d'écriture}, collecter et trier les informations, planifier, rédiger, relire et améliorer.
\end{definition}

L'intuition est simple : aucun bon technicien ne répare une machine sans l'avoir d'abord examinée, sans avoir réuni ses outils et sans vérifier son travail à la fin. Il en va de même pour l'écriture. Rédiger directement, sans préparation, c'est démonter un moteur sans avoir identifié la panne : on perd du temps et on oublie des pièces.

\begin{methode}[Produire un rapport en cinq étapes]
\begin{enumerate}
\item \textbf{Comprendre la situation d'écriture.} Identifier : qui écrit (votre fonction), à qui (le destinataire : chef de travaux, directeur, proviseur), pourquoi (incident, visite, stage), sur quelle consigne.
\item \textbf{Collecter et trier les informations.} Noter les faits : dates, heures, lieux, personnes, matériel, chiffres. Trier : garder l'essentiel, écarter le superflu et tout ce qui relève de l'opinion.
\item \textbf{Planifier.} Construire le plan en trois parties : introduction (contexte et objet), corps (faits organisés chronologiquement ou thématiquement), conclusion (bilan et propositions).
\item \textbf{Rédiger.} Écrire des phrases courtes, à la première personne si vous êtes acteur, au passé composé et à l'imparfait pour les faits, avec un lexique adapté (commun pour le destinataire généraliste, spécialisé pour le destinataire technique).
\item \textbf{Relire et améliorer.} Vérifier le fond (faits exacts, chiffres justes), la forme (orthographe, grammaire) et la structure (en-tête, trois parties, signature).
\end{enumerate}
\end{methode}

\begin{popfigure}
\begin{center}
\begin{tikzpicture}[node distance=0.55cm,
etape/.style={rectangle, rounded corners=3pt, draw=popblue, fill=popblueL, text width=3.4cm, align=center, minimum height=0.85cm, font=\small},
fleche/.style={-{Stealth[length=2.5mm]}, thick, popdark}]
\node[etape, align=center] (s){\textbf{1. Situation}\\ qui ? à qui ? pourquoi ?};
\node[etape, below=of s, align=center] (c){\textbf{2. Collecte et tri}\\ faits, dates, chiffres};
\node[etape, below=of c, align=center] (p){\textbf{3. Plan}\\ introduction, corps, conclusion};
\node[etape, below=of p, align=center] (r){\textbf{4. Rédaction}\\ phrases courtes, objectivité};
\node[etape, below=of r, align=center] (re){\textbf{5. Relecture}\\ fond, forme, structure};
\node[etape, below=of re, draw=poporange, fill=poporangeL, align=center] (a){\textbf{6. Amélioration}\\ version finale};
\draw[fleche] (s) -- (c);
\draw[fleche] (c) -- (p);
\draw[fleche] (p) -- (r);
\draw[fleche] (r) -- (re);
\draw[fleche] (re) -- (a);
\draw[fleche, poporange, dashed] (a.east) .. controls +(1.8,0) and +(1.8,0) .. (r.east) node[midway, right, font=\footnotesize\itshape, popdark]{nouvelle relecture si besoin};
\end{tikzpicture}
\end{center}
\legende{Organigramme de la démarche de production d'un rapport : six étapes enchaînées, avec une boucle de retour de l'amélioration vers la rédaction.}
\end{popfigure}

\courssub{Situation d'écriture n°1 : le rapport d'incident}

\begin{exemplebox}[Situation : court-circuit au laboratoire d'électrotechnique]
Vous êtes élève de Première F3 (Électrotechnique) au lycée technique de Nkongsamba. Le mardi 14 octobre, vers 10 h 30, pendant une séance de travaux pratiques, un court-circuit s'est produit sur un banc d'essai : un étudiant a été légèrement brûlé à la main, le disjoncteur général a coupé l'alimentation du laboratoire. Le chef de travaux vous demande, en tant que délégué de classe présent sur les lieux, un rapport écrit pour le lendemain.
\end{exemplebox}

Appliquons la démarche. \textbf{Situation} : j'écris en tant que témoin et délégué, au chef de travaux, pour rendre compte d'un incident. \textbf{Collecte} : date (14 octobre), heure (10 h 30), lieu (laboratoire d'électrotechnique, banc n°4), fait (court-circuit sur un montage en dérivation), victime (un élève, brûlure légère à la main droite, soigné à l'infirmerie), conséquence matérielle (disjoncteur général déclenché, fusible du banc n°4 hors service), cause probable (fil de connexion dénudé). \textbf{Plan} : introduction (contexte de la séance), corps (déroulement chronologique des faits), conclusion (bilan et propositions : vérifier les câbles, rappeler les consignes de sécurité).

\begin{exoresolu}[Rédiger l'introduction et le début du corps du rapport d'incident]
\textbf{Énoncé.} À partir de la situation précédente, rédigez l'en-tête, l'introduction et le premier paragraphe du corps du rapport (environ 120 mots).
\tcblower
\textbf{Solution détaillée.}

\emph{De} : Essomba Léa, déléguée de la classe de Première F3. \emph{À} : Monsieur le chef de travaux d'électrotechnique. \emph{Date} : 15 octobre. \emph{Objet} : rapport sur l'incident survenu au laboratoire d'électrotechnique le 14 octobre.

\textbf{Introduction.} Le mardi 14 octobre, la classe de Première F3 effectuait une séance de travaux pratiques au laboratoire d'électrotechnique, portant sur le montage de circuits en dérivation. Au cours de cette séance, un incident s'est produit sur le banc d'essai n°4. Le présent rapport relate les faits, en précise les causes probables et propose des mesures de prévention.

\textbf{Corps (début).} Vers 10 h 30, alors que le groupe de travail n°4 mettait le montage sous tension, une étincelle s'est produite au niveau d'un fil de connexion dénudé. Un élève du groupe, touché à la main droite, a ressenti une brûlure légère. Le disjoncteur général du laboratoire s'est immédiatement déclenché, interrompant l'alimentation électrique de l'ensemble des bancs.

\emph{Commentaire} : le texte ne contient que des faits datés, localisés et vérifiables ; le lexique spécialisé (\emph{disjoncteur, montage en dérivation, mise sous tension}) est employé à bon escient car le destinataire est un technicien.
\end{exoresolu}

\courssub{Situation d'écriture n°2 : le rapport de visite d'entreprise}

\begin{exemplebox}[Situation : visite d'une cimenterie]
Dans le cadre de la séquence, votre classe de Première F4 (Génie civil) a visité une cimenterie de la région du Littoral le vendredi 21 novembre. Le professeur de construction vous demande un rapport de visite destiné au proviseur, qui n'est pas un spécialiste du bâtiment.
\end{exemplebox}

Le piège ici est le \cle{destinataire} : le proviseur n'est pas technicien. Le rapport doit donc privilégier le \cle{lexique commun} et expliquer brièvement les termes spécialisés indispensables (par exemple : « le clinker, c'est-à-dire le constituant de base du ciment, obtenu par cuisson du calcaire et de l'argile »). La structure reste identique : introduction (cadre de la visite, date, entreprise, accompagnateurs), corps (accueil, déroulement de la visite atelier par atelier, observations chiffrées : capacité de production, nombre d'employés), conclusion (bilan pédagogique, remerciements, intérêt pour la formation).

\begin{attention}[Ne pas mélanger faits et opinions dans le corps du rapport]
Le corps du rapport relate des \cle{faits} observables et vérifiables : « le four tourne à 1 450 °C », « l'entreprise emploie 320 salariés ». Les appréciations personnelles — « la visite était géniale », « l'usine est trop poussiéreuse », « je pense que le responsable était incompétent » — n'ont pas leur place dans le corps. Si un jugement est nécessaire, il se formule avec mesure dans la \emph{conclusion}, sous forme de bilan ou de proposition : « la visite a permis de concrétiser les notions du cours de technologie ».
\end{attention}

\courssub{La grille d'auto-évaluation}

Avant de remettre votre rapport, évaluez-le vous-même à l'aide de la grille suivante, inspirée des critères de correction de l'examen.

\begin{center}
\begin{tabularx}{\linewidth}{|l|X|c|}
\hline
\rowcolor{popblueL} \textbf{Critère} & \textbf{Ce qui est attendu} & \textbf{Barème} \\
\hline
Structure externe & En-tête complet (émetteur, destinataire, date, objet), signature & 4 \\
\hline
Introduction & Contexte, circonstances, annonce de l'objet du rapport & 3 \\
\hline
Corps du rapport & Faits organisés (ordre chronologique ou thématique), informations exactes et suffisantes & 5 \\
\hline
Conclusion & Bilan, propositions éventuelles, formule adaptée & 3 \\
\hline
Objectivité et lexique & Aucune opinion dans le corps ; lexique commun ou spécialisé adapté au destinataire & 3 \\
\hline
Langue & Orthographe, grammaire, ponctuation, phrases claires & 2 \\
\hline
\rowcolor{popgoldL} \textbf{Total} & & \textbf{20} \\
\hline
\end{tabularx}
\end{center}

\begin{methode}[Améliorer un rapport : la grille de relecture en trois passes]
\begin{enumerate}
\item \textbf{Relecture du fond.} Les faits sont-ils exacts ? Les dates, heures et chiffres sont-ils vérifiés ? Manque-t-il une information essentielle (qui, quoi, où, quand, comment) ?
\item \textbf{Relecture de la forme.} Orthographe (accords, conjugaisons du passé composé et de l'imparfait), ponctuation, majuscules. Reformuler toute phrase qui dépasse deux lignes ou qui exige d'être relue pour être comprise.
\item \textbf{Relecture de la structure.} L'en-tête est-il complet ? Distingue-t-on nettement introduction, corps, conclusion ? Le titre ou l'objet annonce-t-il fidèlement le contenu ?
\end{enumerate}
\end{methode}

\courssub{La correction mutuelle : améliorer un rapport d'élève}

La \cle{relecture croisée} consiste à échanger son brouillon avec un camarade : chacun applique la grille de relecture au texte de l'autre, annote dans la marge et justifie ses remarques. Un regard extérieur repère ce que l'auteur ne voit plus.

\begin{exemplebox}[Avant et après amélioration : comparaison commentée]
\textbf{Avant (extrait du corps d'un rapport d'élève).} « Ensuite on est allés voir le four qui était super impressionnant et il faisait très chaud, le guide nous a expliqué plein de trucs sur le clinker et après on a vu les sacs de ciment, c'était bien. »

\emph{Problèmes relevés} : opinions personnelles (\emph{super impressionnant, c'était bien}), phrase unique trop longue, familiarité (\emph{plein de trucs}), aucune information précise.

\textbf{Après amélioration.} « Le groupe a ensuite visité le hall du four rotatif, où la température atteint 1 450 °C. Le guide a présenté la fabrication du clinker, constituant de base du ciment. La visite s'est poursuivie par l'atelier d'ensachage, qui conditionne 2 400 sacs de 50 kg par heure. »

\emph{Améliorations} : faits chiffrés et vérifiables, lexique adapté avec explication du terme spécialisé, phrases courtes, ordre chronologique marqué (\emph{ensuite, s'est poursuivie}).
\end{exemplebox}

\begin{exercice}[Améliorer un extrait de rapport]
Améliorez l'extrait suivant, tiré du rapport d'incident d'un élève, en appliquant la grille de relecture (fond, forme, structure) : « Le truc a explosé vers 10 h et demi je crois, heureusement personne n'était gravement blessé, c'est inadmissible que le matériel soit aussi vieux, le prof a coupé le courant et on est sortis. »
\end{exercice}

\begin{corrige}[Exercice : améliorer un extrait de rapport]
Extrait amélioré : « Vers 10 h 30, un court-circuit s'est produit sur le banc d'essai n°4, provoquant une étincelle. Un élève a été légèrement brûlé à la main droite et conduit à l'infirmerie. Le professeur de travaux pratiques a immédiatement coupé l'alimentation électrique et fait évacuer la salle. » \emph{Corrections effectuées} : suppression de l'imprécision (\emph{je crois}), du registre familier (\emph{le truc}) et du jugement personnel (\emph{c'est inadmissible}), qui peut éventuellement devenir une proposition en conclusion : « il serait souhaitable de vérifier l'état des câbles des bancs d'essai » ; ajout des faits précis (heure, lieu, victime, conduite tenue).
\end{corrige}

\courssub{Compte rendu collectif et remédiation}

Après la correction mutuelle, la classe procède à un \cle{compte rendu collectif} : le professeur projette ou dicte des extraits de rapports d'élèves \emph{anonymisés}, et la classe identifie ensemble les erreurs récurrentes. Les plus fréquentes sont :

\begin{itemize}
\item l'absence d'en-tête ou un objet trop vague (« Rapport » au lieu de « Rapport sur l'incident du 14 octobre au laboratoire ») ;
\item le mélange des temps (présent de narration au lieu du passé composé et de l'imparfait) ;
\item les opinions personnelles glissées dans le corps du texte ;
\item les chiffres inexacts ou invérifiables ;
\item les phrases trop longues qui perdent le lecteur.
\end{itemize}

La \cle{remédiation} consiste alors à reprendre collectivement un extrait fautif au tableau, à le corriger ensemble, puis à réécrire individuellement la partie défaillante de son propre rapport. C'est la boucle « amélioration » de notre organigramme : on revient à la rédaction, puis on relit une dernière fois.

\begin{savaistu}[Le rapport à l'examen et dans la vie professionnelle]
Au Probatoire et au Baccalauréat technique, la production écrite fonctionnelle — rapport, lettre administrative, compte rendu — est fréquemment évaluée, car elle correspond aux écrits réels du technicien : rapport de stage, rapport de chantier, rapport d'incident, fiche de maintenance. Maîtriser le rapport, c'est donc préparer à la fois votre examen et votre insertion professionnelle : dans une entreprise, un rapport clair et objectif est la preuve écrite de votre sérieux.
\end{savaistu}

\courssub{Trace écrite finale : fiche-modèle de rapport}

\begin{aretenir}[Fiche-modèle de rapport à conserver dans le classeur]
\textbf{Structure externe} : nom et qualité de l'émetteur ; nom et fonction du destinataire ; lieu et date de rédaction ; objet précis du rapport ; signature.

\textbf{Structure interne} :
\begin{enumerate}
\item \emph{Introduction} : contexte et circonstances (qui, quand, où, pourquoi) ; annonce de l'objet.
\item \emph{Corps} : faits uniquement, organisés chronologiquement ou thématiquement ; dates, lieux, chiffres exacts ; aucune opinion.
\item \emph{Conclusion} : bilan des faits, propositions éventuelles, formule de courtoisie adaptée.
\end{enumerate}

\textbf{Langue} : passé composé (faits ponctuels) et imparfait (décor, circonstances) ; phrases courtes ; lexique commun pour un destinataire généraliste, lexique spécialisé (expliqué si nécessaire) pour un destinataire technique.

\textbf{Démarche} : situation $\rightarrow$ collecte et tri $\rightarrow$ plan $\rightarrow$ rédaction $\rightarrow$ relecture (fond, forme, structure) $\rightarrow$ amélioration.
\end{aretenir}

\begin{exercice}[Production complète]
Votre classe de Première technique a visité une brasserie de Douala le 5 décembre. Rédigez le rapport de visite complet (en-tête, introduction, corps, conclusion) destiné au proviseur, en respectant la fiche-modèle. Vous inventerez des faits et des chiffres plausibles et vous les signalerez comme tels à la relecture.
\end{exercice}

\begin{corrige}[Exercice : production complète (éléments de correction)]
Le rapport attendu comporte : un en-tête complet (émetteur : délégué de classe ; destinataire : M. le proviseur ; date : 6 décembre ; objet : rapport de visite de la brasserie du 5 décembre) ; une introduction situant la visite (sortie pédagogique de la classe, encadrée par deux professeurs) ; un corps organisé chronologiquement (accueil par le responsable qualité, visite de la salle de brassage, de la ligne d'embouteillage, du magasin de stockage), avec des faits chiffrés plausibles et un lexique expliqué (par exemple la \emph{fermentation}) ; une conclusion avec bilan pédagogique et remerciements ; enfin la signature. L'évaluation s'effectue avec la grille sur 20 points vue en classe : structure externe (4), introduction (3), corps (5), conclusion (3), objectivité et lexique (3), langue (2).
\end{corrige}

\newpage

\coursec{Activité d'intégration}

\courssub{Objectifs de l'activité}

À l'issue de cette activité d'intégration, l'élève sera capable de :
\begin{itemize}
\item trier et hiérarchiser des informations issues d'un dossier documentaire hétérogène (témoignages, constats, données chiffrées, document iconique) ;
\item élaborer le plan d'un rapport d'incident en respectant la structure externe (en-tête, objet, destinataire, date, signature) et la structure interne (introduction, corps, conclusion) ;
\item rédiger un rapport complet, clair et objectif, en mobilisant les procédés du texte explicatif (définition, explication, illustration, reformulation) ;
\item employer à bon escient le lexique commun et le lexique spécialisé du domaine technique et administratif ;
\item analyser une affiche de sécurité (texte iconique) et l'intégrer comme élément de preuve dans le rapport ;
\item évaluer et améliorer une production écrite à l'aide d'une grille de relecture, puis justifier ses choix de rédaction.
\end{itemize}

\courssub{Situation}

Tu es élève en classe de Première F4 (Génie mécanique) au lycée technique de Nkongsamba, dans la région du Littoral. Tu es le délégué de ta classe. Le mardi 14 octobre 2025, vers 10 h 30, un incident s'est produit dans l'atelier de mécanique générale pendant la séance de travaux pratiques encadrée par M. Etoundi, surveillant d'atelier. Deux élèves de ta classe ont actionné une fraiseuse et un tour à métaux sans autorisation et sans équipement de protection, alors que l'affiche de sécurité placée au-dessus des machines l'interdisait formellement. Le proviseur, M. Nkoulou, te demande de rédiger un rapport d'incident complet et objectif, à remettre au plus tard le vendredi 17 octobre 2025.

Voici le dossier dont tu disposes :

\textbf{Document 1 : témoignage d'Amina, élève de la classe.} « J'étais près de la perceuse. J'ai vu Rodrigue et Junior allumer la fraiseuse alors que M. Etoundi était sorti chercher des outils au magasin. Ils riaient et n'avaient ni lunettes ni blouse. La machine a fait un bruit bizarre, puis il y a eu de la fumée. »

\textbf{Document 2 : témoignage de Rodrigue, l'un des élèves mis en cause.} « On voulait juste essayer la machine comme on avait vu faire en cours. On n'a pas lu l'affiche, on pensait que c'était facile. Le tour a aussi été abîmé quand Junior a voulu l'arrêter en panique. »

\textbf{Document 3 : constat du surveillant d'atelier, M. Etoundi.} « À mon retour à 10 h 40, la fraiseuse dégageait de la fumée et le tour à métaux était bloqué. J'ai immédiatement coupé l'alimentation électrique générale de l'atelier. Aucun élève n'est blessé. Les deux machines sont hors service. »

\textbf{Document 4 : données chiffrées fournies par l'administration.} Deux machines hors service (une fraiseuse et un tour à métaux) ; dommages estimés par le technicien à 1 500 000 FCFA ; aucun blessé ; cours de TP suspendus pour toute la semaine, soit 12 heures de cours perdues pour 3 classes.

\textbf{Document 5 : affiche de sécurité affichée au-dessus des machines.} Un panneau rectangulaire à fond jaune, bordé de rouge, portant le pictogramme d'un engrenage barré et la mention en lettres capitales : « DÉFENSE D'UTILISER LES MACHINES SANS AUTORISATION ET SANS ÉQUIPEMENT DE PROTECTION ».

\courssub{Consignes}

\begin{enumerate}
\item Lis attentivement l'ensemble du dossier. Dans un tableau à deux colonnes, trie les informations : d'une part les faits certains (vérifiables), d'autre part les propos rapportés (témoignages). Justifie ton classement.
\item Identifie le destinataire, l'expéditeur, l'objet, le lieu et la date du rapport. Rédige l'en-tête complet (structure externe).
\item Analyse le document 5 (affiche de sécurité) : décris ses éléments visuels (couleurs, pictogramme, typographie) et explique sa fonction. En quoi constitue-t-il un élément de preuve dans le rapport ?
\item Établis le plan détaillé du corps du rapport : circonstances de l'incident, déroulement des faits, dégâts constatés, causes probables. Classe les informations du dossier dans chaque partie.
\item Rédige l'introduction du rapport (présentation du contexte, annonce de l'objet) et la conclusion (bilan, propositions de mesures pour éviter qu'un tel incident ne se reproduise).
\item Rédige le corps du rapport en employant les procédés du texte explicatif : définition des termes techniques, explication des causes, reformulation des témoignages en style indirect.
\item Relis ton texte : souligne en rouge trois mots du lexique spécialisé (domaine technique ou administratif) et en bleu trois mots du lexique commun. Vérifie que le lexique spécialisé est justifié par le sujet.
\item Échange ta production avec celle d'un camarade. À l'aide de la grille d'évaluation fournie (structure externe : 4 points ; structure interne : 6 points ; exactitude des faits et des chiffres : 4 points ; lexique et correction de la langue : 6 points), note sa copie et propose deux améliorations concrètes.
\item Remanie ton rapport en tenant compte des remarques reçues, puis participe au compte rendu collectif : quelles erreurs reviennent le plus souvent dans la classe ? Comment les corriger ?
\end{enumerate}

\courssub{Ressources à mobiliser}

\begin{aretenir}[Ce qu'il faut se rappeler]
\begin{itemize}
\item \textbf{Structure externe du rapport} : nom de l'institution, lieu et date, destinataire, expéditeur, objet, références éventuelles, signature.
\item \textbf{Structure interne} : une introduction qui situe le contexte et annonce l'objet ; un corps organisé en parties logiques (faits, constats, causes) ; une conclusion qui fait le bilan et propose des solutions.
\item \textbf{Le texte explicatif} : il vise à faire comprendre. Ses procédés sont la définition, l'explication, l'illustration par l'exemple, la reformulation, la comparaison. Il emploie des connecteurs logiques : \emph{parce que, en effet, c'est-à-dire, par conséquent, ainsi}.
\item \textbf{Lexique commun et lexique spécialisé} : le lexique commun est compris de tous (\emph{machine, atelier, accident}) ; le lexique spécialisé appartient à un domaine précis (\emph{fraiseuse, tour à métaux, alimentation électrique, dommages}). Dans un rapport technique, on utilise le lexique spécialisé avec exactitude, en définissant les termes si nécessaire.
\item \textbf{Texte iconique} : une affiche de sécurité communique par l'image (pictogramme, couleurs conventionnelles : rouge = interdiction, jaune = avertissement) et par un texte bref et impératif.
\item \textbf{Qualités du rapport} : objectivité (faits, pas d'opinions), précision (chiffres exacts, heures, noms), clarté (phrases courtes), ordre chronologique ou logique, style indirect pour rapporter les paroles.
\end{itemize}
\end{aretenir}

\courssub{Démarche et corrigé commenté}

\begin{exoresolu}[Production d'un rapport d'incident complet]
Énoncé : à partir du dossier fourni, rédiger le rapport d'incident adressé au proviseur du lycée technique de Nkongsamba, en respectant la structure externe et interne, en mobilisant le texte explicatif et un lexique adapté.
\tcblower
\textbf{Étape 1 : tri des informations.} On distingue les faits certains (constats du surveillant, données chiffrées de l'administration, affiche) des propos rapportés (témoignages des élèves). Les faits certains fondent le corps du rapport ; les témoignages sont cités au style indirect et attribués à leur auteur (« selon Amina... », « Rodrigue déclare que... »). C'est une exigence d'objectivité : un rapport ne présente jamais une rumeur comme un fait.

\textbf{Étape 2 : structure externe.}
\begin{quote}
Lycée technique de Nkongsamba\\
Nkongsamba, le 17 octobre 2025\\
Le délégué de la classe de Première F4\\
À Monsieur le Proviseur\\
\textbf{Objet : rapport d'incident survenu à l'atelier de mécanique générale le mardi 14 octobre 2025.}
\end{quote}

\textbf{Étape 3 : introduction.} Elle situe le contexte et annonce l'objet : « Monsieur le Proviseur, j'ai l'honneur de vous présenter le rapport relatif à l'incident survenu le mardi 14 octobre 2025, vers 10 h 30, dans l'atelier de mécanique générale, pendant la séance de travaux pratiques de la classe de Première F4. Ce rapport expose les circonstances de l'incident, les dégâts constatés, ses causes probables, puis formule des propositions. »

\textbf{Étape 4 : corps du rapport.} Première partie, les faits, dans l'ordre chronologique : absence momentanée du surveillant, mise en marche non autorisée de la fraiseuse par deux élèves, dégagement de fumée, blocage du tour à métaux, retour du surveillant à 10 h 40 et coupure de l'alimentation électrique. Deuxième partie, les constats, avec les chiffres exacts : deux machines hors service (une fraiseuse et un tour à métaux), dommages estimés à 1 500 000 FCFA, aucun blessé, 12 heures de TP perdues pour 3 classes. Troisième partie, les causes : non-respect de l'affiche de sécurité (document iconique à fond jaune bordé de rouge, portant la mention « DÉFENSE D'UTILISER LES MACHINES SANS AUTORISATION ET SANS ÉQUIPEMENT DE PROTECTION »), absence d'équipement de protection, utilisation des machines sans encadrement. Le texte explicatif est mobilisé : « Une fraiseuse, c'est-à-dire une machine-outil destinée à usiner des pièces par enlèvement de matière, exige une formation préalable. En effet, sa mise en marche par un utilisateur non formé peut provoquer une surchauffe du moteur, ce qui explique le dégagement de fumée observé. »

\textbf{Étape 5 : conclusion.} Bilan et propositions : « En conclusion, cet incident, heureusement sans victime, a causé d'importants dégâts matériels et perturbé l'organisation des cours. Afin d'éviter qu'il ne se reproduise, nous proposons : le rappel systématique des consignes de sécurité en début de chaque séance, le verrouillage des machines en l'absence du surveillant, et une sensibilisation de tous les élèves à la lecture des affiches de sécurité. » Suivent la formule de politesse et la signature.

\textbf{Étape 6 : contrôle du lexique.} Lexique spécialisé justifié : \emph{fraiseuse, tour à métaux, alimentation électrique, dommages}. Lexique commun : \emph{machine, élève, atelier, fumée}. Le terme technique est défini à sa première occurrence pour garantir la clarté.

\textbf{Étape 7 : relecture croisée et amélioration.} La grille permet de vérifier : présence de tous les éléments de l'en-tête, plan en trois parties, exactitude des chiffres (1 500 000 FCFA, aucun blessé), emploi du style indirect pour les témoignages, correction orthographique et syntaxique. Les erreurs les plus fréquentes à corriger en remédiation : confusion entre faits et opinions, oubli de l'objet, chiffres inexacts, témoignages recopiés au style direct.
\end{exoresolu}

\courssub{Bilan de l'activité}

Cette activité a permis de mettre en évidence que le rapport est un \cle{texte fonctionnel} régi par des conventions précises : une structure externe codifiée et une structure interne en trois moments. Les élèves ont compris que la qualité d'un rapport repose sur l'\cle{objectivité} (distinction entre faits et témoignages), la \cle{précision} (chiffres, dates, noms exacts) et la \cle{clarté} (texte explicatif, lexique spécialisé défini et justifié). L'analyse de l'affiche de sécurité a montré qu'un texte iconique peut servir d'élément de preuve dans un écrit administratif. Enfin, la relecture croisée avec grille d'évaluation a développé l'autonomie des élèves : évaluer, justifier et améliorer une production font partie intégrante du savoir-faire rédactionnel attendu au Probatoire technique.

\newpage

\coursec{Méthodes et exercices résolus}

\coursec{Produire un rapport}

\courssub{Découverte du rapport : définition et caractéristiques}

\begin{definition}[Le rapport]
Le \cle{rapport} est un texte administratif ou technique à visée informative et décisionnelle. Il rend compte d'une mission, d'un stage, d'un incident, d'une visite ou d'une enquête. Il s'adresse à un destinataire précis (hiérarchie, administration, client) et doit permettre une action ou une décision.
\end{definition}

\begin{propriete}[Caractéristiques fondamentales]
\begin{enumerate}
\item \cle{Objectivité} : faits vérifiables, pas d'opinions personnelles dans le corps (sauf partie « observations » identifiée).
\item \cle{Concision} : phrases courtes, vocabulaire précis, pas de répétitions.
\item \cle{Structure rigide} : en-tête normalisé, introduction codifiée, corps numéroté, conclusion avec recommandations, signature administrative.
\item \cle{Lexique adapté} : termes techniques exacts, définis si le destinataire n'est pas spécialiste.
\item \cle{Présent de vérité} pour les explications générales ; \cle{passé composé} pour les faits datés.
\end{enumerate}
\end{propriete}

\begin{exemplebox}[Différence rapport / lettre administrative]
\begin{center}
\begin{tabularx}{\linewidth}{|l|X|X|}\hline
\rowcolor{popblueL} \textbf{Critère} & \textbf{Rapport} & \textbf{Lettre administrative} \\ \hline
Formule d'appel & \textbf{Absente} & \textbf{Présente} (\emph{Monsieur le Directeur,}) \\ \hline
Formule de politesse finale & \textbf{Absente} (remplacée par \emph{Fait à..., le... Signature}) & \textbf{Présente} (\emph{Veuillez agréer...}) \\ \hline
Structure & Parties numérotées (I, II, III ou 1., 2., 3.) & Paragraphes suivis, sans numérotation obligatoire \\ \hline
Objet & Annoncé dans l'introduction (plan) & Ligne « Objet : » sous la référence \\ \hline
\end{tabularx}
\end{center}
\legende{Le rapport n'est pas une lettre : pas de formules d'appel ni de politesse, structure en parties numérotées, signature administrative finale.}
\end{exemplebox}

\courssub{La structure externe du rapport}

\begin{definition}[Structure externe (mise en page administrative)]
Le rapport respecte un formalisme immuable, noté à l'examen :
\begin{enumerate}
\item \cle{En-tête} (bandeau) : Logo ou nom de l'institution (ex. \emph{RÉPUBLIQUE DU CAMEROUN / MINISTÈRE DE L'ÉDUCATION NATIONALE / LYCÉE TECHNIQUE DE...}), titre centré en majuscules (\emph{RAPPORT DE STAGE}, \emph{RAPPORT D'INCIDENT}, \emph{RAPPORT DE VISITE}), références (numéro, service émetteur), date de rédaction.
\item \cle{Corps du texte} : Paginé, parties titrées et numérotées (I, II, III ou 1., 2., 3.), paragraphes aérés (un saut de ligne entre chaque paragraphe).
\item \cle{Formule finale} : \emph{Fait à [Lieu], le [Date]} (aligné à droite), suivi de la \cle{signature} (nom, prénom, qualité/fonction de l'auteur).
\item \cle{Annexes} (si nécessaire) : Mentionnées en bas de page ou en fin de rapport (\emph{Pièce jointe : photos, procès-verbal, plan...}).
\end{enumerate}
\end{definition}

\courssub{La structure interne : introduction, corps, conclusion}

\begin{aretenir}[Squelette interne type]
\begin{itemize}
\item \textbf{Introduction} (5--8 lignes) : \cle{Qui ?} -- \cle{Pourquoi ?} -- \cle{Quand et où ?} -- \cle{Quoi ?} (objet + annonce du plan).
\item \textbf{Corps} (2 ou 3 parties numérotées) : Faits ordonnés (chronologique ou importance), explications techniques, données chiffrées.
\item \textbf{Conclusion} (3 temps) : Bilan des faits -- Appréciation générale -- Recommandations concrètes.
\item \textbf{Signature administrative} : Lieu, date, signature.
\end{itemize}
\end{aretenir}

\courssub{Rédiger l'introduction d'un rapport}

\begin{methode}[Rédiger l'introduction d'un rapport]
L'introduction d'un rapport est une partie brève (5 à 8 lignes) qui situe le document. Elle obéit à un plan fixe en quatre temps, facile à mémoriser.
\begin{enumerate}
\item \cle{Qui ?} Présenter l'auteur du rapport (nom, fonction, service) et, si besoin, le destinataire (la personne ou l'institution à qui le rapport est adressé).
\item \cle{Pourquoi ?} Indiquer l'origine du rapport : une demande, une mission, un incident, une visite, un stage. C'est la \textbf{cause} de la rédaction.
\item \cle{Quand et où ?} Préciser la date, la période et le lieu des faits rapportés.
\item \cle{Quoi ?} Annoncer l'objet du rapport et le plan qui sera suivi (une phrase du type : « Ce rapport présente d'abord..., ensuite..., enfin... »).
\end{enumerate}
\textbf{Réflexes :} écrire à la première personne si l'on est l'auteur ; employer le passé composé ou le présent de vérité ; rester sobre et objectif ; ne jamais donner son opinion dans l'introduction.
\end{methode}

\begin{exoresolu}[Introduction d'un rapport de stage]
\textbf{Énoncé.} Tu es élève en Première technique au Lycée technique de Douala. Du 3 au 28 février, tu as effectué un stage d'observation à l'entreprise SOCATER (Société camerounaise de travaux et de réseaux), à Bonabéri. À ton retour, ton chef de département, M. Nkoulou, te demande un rapport de stage. Rédige l'introduction de ce rapport (6 à 8 lignes).
\tcblower
\textbf{Solution détaillée.}

\textbf{Étape 1 — Identifier les quatre informations obligatoires.}
\begin{itemize}
\item Qui ? L'élève stagiaire, classe de Première technique, Lycée technique de Douala ; destinataire : M. Nkoulou, chef de département.
\item Pourquoi ? Stage d'observation demandé par l'établissement ; rapport exigé par le chef de département.
\item Quand et où ? Du 3 au 28 février, entreprise SOCATER, Bonabéri (Douala).
\item Quoi ? Déroulement du stage, activités observées, enseignements tirés.
\end{itemize}

\textbf{Étape 2 — Rédiger en suivant le plan en quatre temps.}

\emph{« Dans le cadre de ma formation en classe de Première technique au Lycée technique de Douala, j'ai effectué, du 3 au 28 février, un stage d'observation au sein de l'entreprise SOCATER, située à Bonabéri. Conformément aux instructions de M. Nkoulou, chef de département, je présente dans ce rapport le déroulement de ce stage. J'exposerai d'abord la présentation de l'entreprise, ensuite les activités auxquelles j'ai participé, enfin les enseignements que j'ai tirés de cette expérience. »}

\textbf{Étape 3 — Vérifier.} Les quatre temps sont présents (qui, pourquoi, quand/où, quoi), le ton est objectif, aucune opinion personnelle n'apparaît, le plan est annoncé. L'introduction est conforme.

\textbf{Conclusion.} Une introduction de rapport réussie répond aux questions \cle{qui ? pourquoi ? quand et où ? quoi ?} et annonce le plan, en moins de huit lignes.
\end{exoresolu}

\courssub{Rédiger le corps et la conclusion d'un rapport}

\begin{methode}[Construire le corps et la conclusion]
\begin{enumerate}
\item \cle{Diviser} le corps en parties numérotées et titrées (I, II, III ou 1., 2., 3.), chacune correspondant à un aspect du sujet : constat, faits, causes, résultats.
\item \cle{Exposer les faits} dans un ordre logique (chronologique ou du plus important au moins important), avec des données précises : dates, chiffres, noms, lieux.
\item \cle{Rester objectif} : phrases courtes, vocabulaire précis, pas de sentiments personnels ; si une opinion est demandée, la placer dans une partie « observations » ou « suggestions » clairement séparée.
\item \cle{Rédiger la conclusion} en trois temps : bilan bref des faits, appréciation générale, propositions ou recommandations concrètes.
\item \cle{Terminer} par la formule administrative : lieu et date, signature de l'auteur.
\end{enumerate}
\textbf{Formules à retenir :} « Il ressort de ce qui précède que... », « En conséquence, il est recommandé de... », « Fait à Douala, le... ».
\end{methode}

\begin{exoresolu}[Corps et conclusion d'un rapport d'incident]
\textbf{Énoncé.} En tant que délégué de classe, tu rédiges un rapport sur la panne du groupe électrogène du lycée survenue le 12 mars. On te demande de rédiger le corps (deux parties : constat, causes) et la conclusion avec une recommandation.
\tcblower
\textbf{Solution détaillée.}

\textbf{Étape 1 — Organiser le corps en parties numérotées.}

\emph{« 1. Constat. Le lundi 12 mars, à 8 h 15, le groupe électrogène du lycée s'est arrêté brutalement, cinq minutes après son démarrage. Les travaux pratiques d'électrotechnique, prévus de 8 h à 10 h, n'ont pas pu avoir lieu. L'examen visuel montre une fuite de carburant au niveau du réservoir.}

\emph{2. Causes. D'après le technicien de service, la fuite provient d'un joint usé qui n'a pas été remplacé lors de la dernière révision, effectuée il y a plus de huit mois. L'absence d'entretien régulier explique donc cette panne. »}

\textbf{Étape 2 — Rédiger la conclusion en trois temps} (bilan, appréciation, recommandation).

\emph{« En conclusion, la panne du 12 mars a perturbé le déroulement des travaux pratiques. Elle est due à un défaut d'entretien du groupe électrogène. En conséquence, il est recommandé d'établir un calendrier de révision trimestrielle et de remplacer immédiatement le joint défectueux. Fait à Douala, le 13 mars. Le délégué de classe. »}

\textbf{Étape 3 — Vérifier.} Faits précis (date, heure), causes distinguées du constat, objectivité respectée, recommandation concrète, formule finale administrative.

\textbf{Conclusion.} Le corps expose des faits vérifiables en parties numérotées ; la conclusion fait le bilan et propose des solutions datées et réalistes.
\end{exoresolu}

\courssub{Employer le texte explicatif dans un rapport}

\begin{methode}[Expliquer clairement un fait ou un fonctionnement]
Le rapport utilise souvent le \cle{texte explicatif} : il s'agit de faire comprendre un phénomène, une cause ou un fonctionnement.
\begin{enumerate}
\item \cle{Identifier} ce qu'on explique (le \emph{thème}) et ce qu'on dit dessus (l'explication).
\item \cle{Utiliser les marques de l'explication} : connecteurs de cause et de conséquence (\emph{parce que, car, puisque, c'est pourquoi, donc, ainsi, en effet}), présentatifs (\emph{il s'agit de, c'est, cela consiste à}).
\item \cle{Employer le présent de vérité} (présent intemporel) pour les explications générales, et le vocabulaire précis du domaine.
\item \cle{Ordonner} l'explication : de la cause vers la conséquence, ou du général vers le particulier ; définir les termes techniques au premier emploi.
\end{enumerate}
\textbf{Réflexe :} un paragraphe explicatif = une idée + sa cause + sa conséquence, reliées par des connecteurs.
\end{methode}

\begin{exoresolu}[Transformer un récit en texte explicatif]
\textbf{Énoncé.} Voici un passage de rapport mal rédigé : « La machine est tombée en panne mardi. L'huile n'était plus bonne. Le moteur a chauffé. » Réécris ce passage en un paragraphe explicatif cohérent, en employant au moins trois marques de l'explication.
\tcblower
\textbf{Solution détaillée.}

\textbf{Étape 1 — Repérer la chaîne cause $\rightarrow$ conséquence.}
Cause première : huile dégradée. Conséquence intermédiaire : mauvaise lubrification, donc chauffe du moteur. Conséquence finale : panne de la machine mardi.

\textbf{Étape 2 — Choisir les connecteurs.} \emph{parce que} (cause), \emph{en effet} (confirmation), \emph{c'est pourquoi} (conséquence).

\textbf{Étape 3 — Rédiger au présent de vérité pour l'explication, passé composé pour le fait daté.}

\emph{« La machine est tombée en panne mardi parce que l'huile du moteur était dégradée. En effet, une huile usée ne lubrifie plus correctement les pièces en mouvement : le frottement augmente et le moteur chauffe. C'est pourquoi le dispositif de sécurité a provoqué l'arrêt automatique de la machine. »}

\textbf{Étape 4 — Vérifier.} Trois marques explicatives (\emph{parce que, en effet, c'est pourquoi}), présent de vérité pour le fonctionnement général, passé composé pour l'événement, ordre logique cause $\rightarrow$ conséquence.

\textbf{Conclusion.} Un texte explicatif relie les faits par des connecteurs de cause et de conséquence et distingue le temps de l'événement (passé) du temps de l'explication (présent de vérité).
\end{exoresolu}

\courssub{Distinguer et employer lexique commun et lexique spécialisé}

\begin{methode}[Choisir le bon niveau de lexique]
\begin{enumerate}
\item \cle{Identifier le lexique commun} : mots compris de tous (\emph{machine, outil, eau, réparer}).
\item \cle{Identifier le lexique spécialisé} : termes propres à un métier ou à une science (\emph{disjoncteur, culasse, dosage, tension nominale}).
\item \cle{Adapter au destinataire} : dans un rapport technique destiné à des spécialistes, employer le lexique spécialisé ; pour un destinataire non spécialiste, définir chaque terme technique à sa première apparition (entre parenthèses ou dans une courte phrase de définition).
\item \cle{Éviter} le mélange des registres : ne pas remplacer un terme technique exact par un mot vague (\emph{truc, machin, bidule}).
\end{enumerate}
\textbf{Réflexe :} terme spécialisé + définition brève = précision sans obscurité.
\end{methode}

\begin{exoresolu}[Lexique commun ou spécialisé ?]
\textbf{Énoncé.} 1) Dans la phrase suivante, souligne les termes du lexique spécialisé : « Le technicien a vérifié le disjoncteur, remplacé le fusible fondu et mesuré la tension aux bornes du moteur. » 2) Réécris cette phrase pour un destinataire non spécialiste en définissant un des termes techniques.
\tcblower
\textbf{Solution détaillée.}

\textbf{Étape 1 — Repérage.} Les termes du lexique spécialisé (domaine de l'électricité) sont : \textbf{disjoncteur}, \textbf{fusible}, \textbf{tension}, \textbf{bornes}. Les mots \emph{technicien, vérifié, remplacé, mesuré, moteur} appartiennent au lexique commun.

\textbf{Étape 2 — Justifier.} Ces quatre termes désignent des réalités précises du métier d'électricien ; ils ne sont pas employés dans la conversation ordinaire avec ce sens exact.

\textbf{Étape 3 — Réécriture avec définition.}

\emph{« Le technicien a vérifié le disjoncteur, c'est-à-dire l'appareil qui coupe automatiquement le courant en cas de défaut, puis il a remplacé le fusible fondu et mesuré la tension aux bornes du moteur. »}

\textbf{Étape 4 — Vérifier.} Le terme technique est conservé (précision) et suivi d'une définition en lexique commun (clarté) introduite par \emph{c'est-à-dire}.

\textbf{Conclusion.} Dans un rapport, on garde le terme spécialisé exact et on le définit brièvement lorsque le destinataire n'est pas un spécialiste.
\end{exoresolu}

\courssub{Lire et exploiter un texte fonctionnel ou une affiche de sécurité}

\begin{methode}[Analyser un texte fonctionnel et une affiche de sécurité]
\begin{enumerate}
\item \cle{Identifier la nature} du document : notice, consigne, affiche, panneau ; puis son \textbf{émetteur} (qui parle ?) et son \textbf{destinataire} (à qui ?).
\item \cle{Dégager la fonction} : informer, interdire, obliger, avertir, guider. Sur une affiche de sécurité, repérer le code couleur normalisé (ISO 7010) : \textbf{rouge} = interdiction ou danger (rond barré ou triangle), \textbf{bleu} = obligation (rond plein), \textbf{jaune} = avertissement/attention (triangle), \textbf{vert} = issue, secours, matériel de sécurité (carré/rectangle).
\item \cle{Repérer les marques linguistiques} : infinitifs de consigne (\emph{porter, ne pas toucher}), impératifs, phrases nominales, pictogrammes, caractères gras et majuscules.
\item \cle{Dégager le message} en une phrase complète : qui doit faire quoi, où, et pourquoi.
\item \cle{Exploiter} ces informations dans un rapport : citer la consigne exacte, indiquer si elle était respectée ou non.
\end{enumerate}
\end{methode}

\begin{exoresolu}[Exploiter une affiche de sécurité dans un rapport]
\textbf{Énoncé.} Dans l'atelier de ton lycée figure l'affiche : « PORT DES LUNETTES DE PROTECTION OBLIGATOIRE », accompagnée d'un pictogramme bleu rond. 1) Analyse cette affiche (nature, fonction, marques). 2) Rédige deux phrases de rapport signalant que deux élèves ne respectaient pas cette consigne le 5 mai.
\tcblower
\textbf{Solution détaillée.}

\textbf{Étape 1 — Analyse de l'affiche.}
\begin{itemize}
\item Nature : affiche d'indication sécuritaire (texte iconique : texte + image).
\item Émetteur : l'administration du lycée ; destinataires : élèves et personnel de l'atelier.
\item Fonction : obliger. Le fond bleu et le pictogramme rond signalent une \textbf{obligation} (port d'EPI).
\item Marques linguistiques : phrase nominale sans verbe conjugué, majuscules, adjectif \emph{obligatoire} ; marque iconique : pictogramme représentant des lunettes.
\item Message : toute personne dans l'atelier doit porter des lunettes de protection pour éviter les projections.
\end{itemize}

\textbf{Étape 2 — Rédaction dans le rapport.}

\emph{« Le 5 mai, au cours de la séance de travaux pratiques, deux élèves effectuaient un perçage sans lunettes de protection, alors que l'affiche "Port des lunettes de protection obligatoire" est apposée à l'entrée de l'atelier. Cette consigne de sécurité, rappelée par un panneau bleu d'obligation, n'a donc pas été respectée. »}

\textbf{Étape 3 — Vérifier.} La consigne est citée exactement, les faits sont datés et localisés, le constat reste objectif.

\textbf{Conclusion.} Lire une affiche de sécurité, c'est identifier sa fonction (ici l'obligation, signalée par le bleu) ; dans un rapport, on cite la consigne et on rapporte les faits avec précision.
\end{exoresolu}

\courssub{Produire et améliorer un rapport complet}

\begin{methode}[Rédiger puis améliorer un rapport complet]
\begin{enumerate}
\item \cle{Avant d'écrire} : rassembler les faits (notes, dates, chiffres), identifier le destinataire et l'objet, dresser le plan (introduction, 2 ou 3 parties, conclusion).
\item \cle{Structure externe} : en-tête (logo ou nom de l'institution, titre « RAPPORT DE... », références, date), corps paginé avec parties numérotées, signature.
\item \cle{Structure interne} : introduction (qui, pourquoi, quand/où, quoi), corps (faits ordonnés, explications, données), conclusion (bilan + recommandations).
\item \cle{Réviser} en trois relectures : contenu (faits exacts et complets ?), organisation (plan logique, titres numérotés ?), langue (orthographe, accords, ponctuation, connecteurs, lexique adapté).
\item \cle{Améliorer} : remplacer les mots vagues par des termes précis, supprimer les répétitions et les opinions hors de la partie prévue, vérifier la concordance des temps.
\end{enumerate}
\end{methode}

\begin{exoresolu}[Améliorer un extrait de rapport]
\textbf{Énoncé.} Voici un extrait de rapport à améliorer : « Moi je pense que la visite s'est bien passée. On a vu des trucs intéressants. Le directeur il était gentil. Je recommande de refaire ça. » Réécris ce passage selon les normes du rapport (objectivité, lexique, syntaxe).
\tcblower
\textbf{Solution détaillée.}

\textbf{Étape 1 — Relever les défauts.}
\begin{itemize}
\item Opinion personnelle mal placée : « Moi je pense ».
\item Lexique vague : « des trucs », « ça ».
\item Syntaxe orale : « Le directeur il était gentil » (dislocation).
\item Imprécision : aucune date, aucun fait concret.
\end{itemize}

\textbf{Étape 2 — Corriger point par point.} Supprimer « je » au profit d'un constat ; remplacer « trucs » par un nom précis ; supprimer la dislocation ; transformer l'appréciation en constat mesurable ; formuler une recommandation précise.

\textbf{Étape 3 — Version améliorée.}

\emph{« La visite de l'entreprise s'est déroulée dans de bonnes conditions. Les élèves ont observé le fonctionnement de la chaîne de production et le poste de contrôle qualité. Le directeur a répondu à l'ensemble des questions posées. Au vu de l'intérêt pédagogique de cette visite, il est recommandé de la renouveler chaque année pour les classes de Première technique. »}

\textbf{Étape 4 — Vérifier.} Ton objectif, lexique précis, phrases correctes, recommandation concrète et datée dans son principe.

\textbf{Conclusion.} Améliorer un rapport consiste à remplacer l'oral par l'écrit administratif : objectivité, précision du lexique, correction syntaxique et recommandations exploitables.
\end{exoresolu}

\courssub{Compte rendu et remédiation}

\begin{aretenir}[Check-list finale avant de rendre un rapport (Probatoire)]
\begin{itemize}
\item [ ] \textbf{En-tête} : Institution, Titre (RAPPORT DE...), Références, Date.
\item [ ] \textbf{Introduction} : 4 temps (Qui, Pourquoi, Quand/Où, Quoi + Plan) en 5--8 lignes.
\item [ ] \textbf{Corps} : 2 ou 3 parties numérotées (I, II, III), titres clairs, faits datés/chiffrés, objectivité, connecteurs logiques.
\item [ ] \textbf{Explications} : Présent de vérité + connecteurs cause/conséquence, termes techniques définis si besoin.
\item [ ] \textbf{Conclusion} : 3 temps (Bilan, Appréciation, Recommandations concrètes).
\item [ ] \textbf{Signature} : \emph{Fait à [Lieu], le [Date]} + Signature (Nom, Prénom, Qualité).
\item [ ] \textbf{Langue} : Zéro \emph{je} (sauf intro), zéro \emph{truc/machin}, zéro dislocation, accords/orthographe vérifiés.
\end{itemize}
\end{aretenir}

\begin{attention}[Les 5 erreurs qui coûtent des points au Bac / Probatoire]
\begin{enumerate}
\item \textbf{Confondre rapport et lettre} : le rapport n'a ni formule d'appel (\emph{Monsieur,}) ni formule de politesse finale (\emph{Veuillez agréer...}). Il se termine par \emph{Fait à..., le...} suivi de la signature. Mélanger les deux genres fait perdre des points de présentation.
\item \textbf{Donner son opinion partout} : le rapport est objectif ; les appréciations personnelles ne sont admises que dans la conclusion ou une partie « suggestions » clairement identifiée. Un « je pense que » dans le corps du rapport est sanctionné.
\item \textbf{Oublier les quatre temps de l'introduction} : qui, pourquoi, quand/où, quoi (avec annonce du plan). Une introduction qui commence directement par les faits est incomplète.
\item \textbf{Employer un lexique vague ou oral} : \emph{truc, machin, super, gentil} sont à bannir ; chaque terme technique doit être exact, et défini si le destinataire n'est pas spécialiste.
\item \textbf{Négliger la structure externe} : absence de titre (\emph{Rapport de...}), de date, de parties numérotées ou de signature. À l'examen, la mise en forme du rapport est notée : titre en évidence, parties titrées et numérotées, paragraphes aérés.
\end{enumerate}
\end{attention}

\newpage

\coursec{Exercices}

\courssub{Je vérifie mes connaissances}

\begin{exercice}[QCM : les caractéristiques du rapport]
Pour chaque question, choisis la bonne réponse.
\begin{enumerate}
\item Le rapport est un texte qui rend compte :
\begin{enumerate}
\item d'événements imaginaires ;
\item de faits réels observés ou vécus ;
\item d'opinions personnelles ;
\item de rêves.
\end{enumerate}
\item L'introduction d'un rapport présente :
\begin{enumerate}
\item la bibliographie ;
\item le contexte, l'objet et le plan du rapport ;
\item les remerciements ;
\item la signature du rédacteur.
\end{enumerate}
\item Dans un rapport, le temps de verbe dominant est :
\begin{enumerate}
\item le futur simple ;
\item le passé composé ou l'imparfait selon la situation ;
\item le subjonctif ;
\item l'impératif.
\end{enumerate}
\item Le lexique spécialisé est :
\begin{enumerate}
\item le vocabulaire de tous les jours ;
\item le vocabulaire propre à un domaine d'activité ;
\item le vocabulaire familier ;
\item le vocabulaire des insultes.
\end{enumerate}
\end{enumerate}
\end{exercice}

\begin{exercice}[Vrai ou faux]
Dis si chaque affirmation est vraie ou fausse, puis justifie ta réponse en une phrase.
\begin{enumerate}
\item La conclusion d'un rapport peut introduire des faits nouveaux.
\item Un rapport doit être objectif : on évite les marques de subjectivité (\emph{je pense que}, \emph{à mon avis}).
\item L'objet du rapport est facultatif dans la structure externe.
\item Une affiche de sécurité combine souvent un pictogramme, une couleur et une consigne.
\end{enumerate}
\end{exercice}

\begin{exercice}[Définitions]
Définis en deux ou trois phrases chacun des termes suivants :
\begin{enumerate}
\item le rapport ;
\item le texte explicatif ;
\item le lexique commun ;
\item la structure externe d'un rapport.
\end{enumerate}
\end{exercice}

\begin{exercice}[Repérage rapide]
Voici un extrait : \og Le 12 mars, le technicien a constaté une fuite au niveau de la canalisation principale. Il a immédiatement fermé la vanne d'arrêt. \fg
\begin{enumerate}
\item Relève deux éléments qui montrent qu'il s'agit d'un compte rendu objectif.
\item Indique le temps de verbe employé et justifie ce choix.
\item Propose un titre d'objet pour ce passage s'il figurait dans un rapport.
\end{enumerate}
\end{exercice}

\courssub{Je m'entraîne}

\begin{exercice}[Lexique commun ou spécialisé ?]
Voici 15 termes extraits d'un rapport de chantier BTP : \emph{béton, maison, coffrage, mur, ferraillage, toit, plan, dalle, fenêtre, mortier, porte, étaiement, gravats, échafaudage, chambre}.
\begin{enumerate}
\item Classe ces termes en deux colonnes : lexique commun / lexique spécialisé.
\item Choisis trois termes spécialisés et définis-les pour un lecteur non spécialiste.
\item Explique en quelques lignes pourquoi un rapport technique emploie des termes spécialisés.
\end{enumerate}
\end{exercice}

\begin{exercice}[De la subjectivité à l'objectivité]
Voici le récit subjectif d'un accident d'atelier : \og Franchement, je trouve que c'était horrible ! Le pauvre Jean s'est coupé le doigt avec la scie, j'ai eu très peur, c'était vraiment de sa faute, il ne fait jamais attention. \fg
\begin{enumerate}
\item Relève toutes les marques de subjectivité et de jugement.
\item Réécris ce passage en un paragraphe de rapport objectif (faits, temps appropriés, vocabulaire neutre).
\item Ajoute une phrase de recommandation de sécurité à la fin du paragraphe.
\end{enumerate}
\end{exercice}

\begin{exercice}[Analyse d'une affiche de sécurité]
Dans l'atelier de ton lycée, une affiche représente un cercle rouge barré sur une silhouette qui court, avec la mention : \og Ne pas courir dans l'atelier \fg.
\begin{enumerate}
\item Décris les trois composantes de l'affiche : pictogramme, couleur, consigne.
\item Explique le rôle de chaque composante dans la transmission du message.
\item Rédige un court paragraphe explicatif (8 à 10 lignes) sur le fonctionnement de cette affiche.
\end{enumerate}
\end{exercice}

\begin{exercice}[Remédiation : corriger un rapport]
Le rapport d'un élève contient les huit erreurs suivantes : objet absent ; date manquante ; faits nouveaux en conclusion ; deux termes de jargon non définis ; connecteurs logiques absents ; marques de subjectivité ; paragraphes non articulés ; signature oubliée.
\begin{enumerate}
\item Pour chaque erreur, indique la règle du rapport qu'elle viole.
\item Propose une correction concrète pour chacune des huit erreurs.
\item Rédige l'objet et la phrase de conclusion corrigés de ce rapport (thème : visite d'une usine de transformation de manioc).
\end{enumerate}
\end{exercice}

\courssub{Je me prépare au Bac}

\begin{exercice}[Production écrite : rapport de stage]
Suite à ton stage d'observation de deux semaines dans une entreprise de ta localité (menuiserie, garage, exploitation agricole, cybercafé, atelier de couture...), ton professeur de filière te demande de rédiger le rapport de stage.
\begin{enumerate}
\item Rédige la structure externe complète : nom du rédacteur, classe, destinataire, lieu, date, objet.
\item Rédige l'introduction (contexte du stage, objectifs, annonce du plan) en 10 à 15 lignes.
\item Rédige le corps du rapport en trois parties articulées par des connecteurs logiques : présentation de l'entreprise, activités observées et réalisées, difficultés rencontrées.
\item Rédige la conclusion (bilan du stage, acquis, remerciements) sans introduire de fait nouveau.
\item Souligne dans ton texte trois termes de lexique spécialisé et définis-les en note ou entre parenthèses.
\end{enumerate}
\end{exercice}

\begin{exercice}[Texte explicatif et rapport d'incident]
Tu es délégué de classe. Un incendie de petit ampleur s'est déclaré dans le laboratoire de sciences physiques de ton lycée technique. Le proviseur te demande un rapport circonstancié.
\begin{enumerate}
\item Rédige l'en-tête du rapport (structure externe complète).
\item Rédige le corps du rapport en expliquant objectivement : le déroulement des faits (heure, lieu, causes probables), les mesures prises immédiatement, les dégâts constatés.
\item Insère dans ton rapport un paragraphe explicatif qui décrit le fonctionnement de l'extincteur utilisé, en employant les connecteurs d'explication (\emph{d'abord, ensuite, enfin, c'est-à-dire}).
\item Propose en conclusion deux recommandations de sécurité, rédigées sous forme de consignes comme sur une affiche (verbe à l'infinitif).
\end{enumerate}
\end{exercice}

\begin{exercice}[Réception et production : visite d'entreprise]
Ta classe a visité une cimenterie de la région du Littoral. Voici des notes prises par un élève : \og visite le 5 juin ; accueil par le chef de production ; silos de stockage ; broyeur ; four rotatif à 1450 degrés ; ensachage ; 200 ouvriers ; port du casque obligatoire ; bruit très fort ; poussière ; production : 5000 sacs par jour. \fg
\begin{enumerate}
\item À partir de ces notes, rédige un rapport de visite complet (structure externe, introduction, corps en deux ou trois parties, conclusion) en 25 à 30 lignes.
\item Relève dans tes notes quatre termes de lexique spécialisé et définis-les pour un lecteur non spécialiste.
\item Rédige le texte d'une affiche de sécurité destinée aux visiteurs de la cimenterie : pictogramme décrit, couleur choisie, consigne rédigée. Justifie tes choix en trois lignes.
\item Explique en un paragraphe en quoi le rapport que tu as rédigé est un texte objectif et fonctionnel.
\end{enumerate}
\end{exercice}

\newpage

\coursec{Corrigés des exercices}

\coursec{VI : Produire un rapport}

\courssub{Découverte du rapport : définition et caractéristiques}

\begin{definition}[Le rapport]
C'est un \textbf{texte fonctionnel}, \textbf{objectif} et \textbf{structuré}, qui rend compte par écrit de \textbf{faits réels} (observés, vécus, mesurés) dans un cadre professionnel ou scolaire, à destination d'un destinataire précis, en vue d'informer, de justifier ou d'aider à la décision.
\end{definition}

\begin{aretenir}[Caractéristiques principales]
\begin{itemize}
\item \textbf{Objectivité :} Pas de « je », pas d'opinion, pas d'affect. Usage de l'impersonnel (voix passive, « on », tournures infinitives), verbes d'action factuelle, données chiffrées, localisations précises.
\item \textbf{Structure codifiée :} Structure externe (en-tête, pied de page) et structure interne (introduction, corps, conclusion) obligatoires.
\item \textbf{Temps de la narration :} \textbf{Passé composé} (actions ponctuelles, achevées, datées) et \textbf{Imparfait} (états, descriptions, actions habituelles, contexte).
\item \textbf{Lexique :} Précis, souvent \textbf{spécialisé} (vocabulaire technique du métier), défini à sa première occurrence.
\item \textbf{Finalité :} Outil de travail, de communication professionnelle, de traçabilité et d'aide à la décision.
\end{itemize}
\end{aretenir}

\courssub{La structure externe du rapport}

\begin{methode}[Rédiger l'en-tête et le pied de page (structure externe)]
\begin{enumerate}
\item \textbf{En-tête (mentions obligatoires) :}
\begin{itemize}
\item \textbf{Rédacteur :} Nom, Prénom, Classe/Spécialité, Fonction (stagiaire, délégué...).
\item \textbf{Destinataire :} Titre, Nom, Structure (Proviseur, Chef d'entreprise, Professeur...).
\item \textbf{Lieu :} Ville de rédaction.
\item \textbf{Date :} Date de rédaction du rapport (pas seulement la date des faits).
\item \textbf{Objet :} Formule concise : « Rapport de [nature] : [sujet précis] — [période/date] ». \cle{Objet obligatoire}
\item \textbf{Références (si besoin) :} N° de commande, de stage, de bon d'intervention...
\end{itemize}
\item \textbf{Pied de page (mentions de validité) :}
\begin{itemize}
\item \textbf{Signature :} Manuscrit du rédacteur (nom lisible + qualité).
\item \textbf{Visa (si hiérarchie) :} Signature du tuteur / supérieur validant le contenu.
\item \textbf{Pièces jointes (PJ) :} Liste numérotée (photos, plans, fiches techniques, PV...).
\item \textbf{Diffusion :} Liste des destinataires en copie (CC).
\end{itemize}
\end{enumerate}
\end{methode}

\begin{exemplebox}[Modèle d'en-tête type]
\begin{tabularx}{\linewidth}{lX}
\textbf{Rédacteur :} & NOM Prénom, Élève 1re T [Spécialité] \\
\textbf{Destinataire :} & Monsieur le [Titre/Fonction], [Établissement/Entreprise] \\
\textbf{Lieu :} & [Ville] \\
\textbf{Date :} & [Jour Mois Année] \\
\textbf{Objet :} & Rapport de [stage / visite / incident / expérience] : [Sujet précis] — [Période]
\end{tabularx}
\end{exemplebox}

\courssub{La structure interne : introduction, corps, conclusion}

\begin{methode}[Rédiger l'introduction (3 mouvements obligatoires)]
\begin{enumerate}
\item \textbf{Le contexte (Cadre) :} Cadre officiel (module, stage, mission), dates, lieu, organisme d'accueil, tuteur.
\item \textbf{L'objet et les objectifs :} Sujet précis du rapport + buts visés (découvrir, observer, réaliser, diagnostiquer...).
\item \textbf{L'annonce du plan :} « Ce rapport présente successivement... », « Il s'articule en trois parties : 1... 2... 3... ».
\end{enumerate}
\end{methode}

\begin{methode}[Rédiger le corps du rapport]
\begin{enumerate}
\item \textbf{Découpage en parties équilibrées :} 2 à 3 parties principales (ex. : 1. Présentation structure / 2. Activités techniques / 3. Analyse / Difficultés).
\item \textbf{Un paragraphe = une idée force :} Amorce (connecteur), développement factuel, terme (transition ou synthèse partielle).
\item \textbf{Connecteurs logiques variés :} Chronologie (\og D'abord, Ensuite, Enfin \fg), Cause/Conséquence (\og Par conséquent, C'est pourquoi \fg), Opposition (\og Cependant, Toutefois \fg), Addition (\og En outre, Par ailleurs \fg).
\item \textbf{Objectivité stricte :} Passé composé / Imparfait, 3e personne, lexique spécialisé défini, données chiffrées (\SI{5000}{sacs/jour}, \SI{1450}{\celsius}).
\item \textbf{Illustrations (optionnel mais recommandé) :} Schémas, tableaux, photos légendés, insérés au fil du texte ou en annexe (PJ).
\end{enumerate}
\end{methode}

\begin{methode}[Rédiger la conclusion (3 mouvements obligatoires)]
\begin{enumerate}
\item \textbf{Bilan / Synthèse des résultats :} Rappel des faits marquants et acquis techniques (sans fait nouveau).
\item \textbf{Enseignements / Recommandations :} Apports pour la formation, orientation, propositions d'amélioration, consignes de sécurité.
\item \textbf{Remerciements :} Destinataire, tuteur, personnel, direction (formule administrative standard).
\end{enumerate}
\begin{attention}[Piège fatal] La conclusion \textbf{n'introduit jamais de faits nouveaux} (ni données, ni événements, ni anecdotes non traités dans le corps).
\end{attention}
\end{methode}

\courssub{Le texte explicatif au service du rapport}

\begin{definition}[Le texte explicatif]
C'est un type de texte qui répond à la question \og Comment ? \fg ou \og Pourquoi ? \fg en décrivant un phénomène, un mécanisme, une procédure ou un processus. Il utilise des \textbf{connecteurs logiques} (cause, conséquence, chronologie, finalité) et un \textbf{lexique précis}, souvent spécialisé.
\end{definition}

\begin{aretenir}[Connecteurs de l'explication]
\begin{center}
\begin{tabularx}{\linewidth}{|l|X|}\hline
\rowcolor{popblueL}
\textbf{Relation logique} & \textbf{Connecteurs / Marqueurs} \\ \hline
Chronologie / Procédure & D'abord, Ensuite, Puis, Enfin, Étape 1 / 2... \\ \hline
Cause & Parce que, Comme, Puisque, Du fait de, Grâce à, En raison de \\ \hline
Conséquence & Par conséquent, C'est pourquoi, Donc, Si bien que, De sorte que \\ \hline
Finalité & Pour, Afin de, Dans le but de, De manière à + infinitif \\ \hline
Définition / Précision & C'est-à-dire, Autrement dit, Il s'agit de, Notamment \\ \hline
\end{tabularx}
\end{center}
\end{aretenir}

\begin{exemplebox}[Paragraphe explicatif type : fonctionnement extincteur poudre ABC]
L'extincteur utilisé est un modèle \textbf{à poudre polyvalente ABC (phosphate d'ammonium)}, efficace sur feux de classe A (solides), B (liquides) et C (gaz). \textbf{D'abord,} on retire la \textbf{goupille de sécurité} qui verrouille la poignée de commande, ce qui arme l'appareil. \textbf{Ensuite,} on saisit la lance (ou le tuyau) et on la dirige \textbf{vers la base des flammes} (siège du feu), non pas vers le haut du panache. \textbf{Enfin,} on appuie fermement sur la \textbf{poignée de commande} (ou le levier) : cela ouvre la vanne interne, la pression du gaz propulseur (azote ou air comprimé) chasse la poudre hors du réservoir à travers le tube plongeur et la lance. \textbf{C'est-à-dire} que la poudre forme un nuage qui \textbf{étouffe le feu par inhibition chimique} : les particules fines interceptent les radicaux libres (H*, OH*) de la réaction en chaîne de la combustion, stoppant l'exothermicité, tout en formant une barrière isolante sur le combustible.
\end{exemplebox}

\courssub{Lexique commun et lexique spécialisé}

\begin{definition}[Lexique commun vs Lexique spécialisé]
\begin{itemize}
\item \textbf{Lexique commun :} Mots du langage courant, compris de tous, sans appartenance technique (ex. : \og maison \fg, \og manger \fg, \og grand \fg, \og mur \fg).
\item \textbf{Lexique spécialisé (vocabulaire technique) :} Termes propres à un domaine professionnel (BTP, Mécanique, Informatique, Santé...), à sens univoque, garantissant précision, concision et communication experte (ex. : \og coffrage \fg, \og ferraillage \fg, \og valise de diagnostic \fg, \og clinker \fg).
\end{itemize}
\end{definition}

\begin{methode}[Gérer le lexique spécialisé dans un rapport]
\begin{enumerate}
\item \textbf{Identifier} les termes techniques indispensables à la précision.
\item \textbf{Définir} chaque terme à sa \textbf{première occurrence} (entre parenthèses, en note de bas de page, ou en glossaire final).
\item \textbf{Souligner} ou mettre en \textbf{gras} les termes définis pour la correction / l'évaluation.
\item \textbf{Éviter} le jargon abusif (termes anglais non francisés, acronymes non définis, argot de métier).
\end{enumerate}
\end{methode}

\courssub{Lecture de textes fonctionnels : affiches d'indications sécuritaires}

\begin{definition}[Affiche de sécurité normalisée (ISO 3864 / ISO 7010)]
Dispositif de prévention primaire combinant \textbf{trois composantes redondantes} pour une efficacité maximale (bruit, stress, barrière linguistique, non-lecteurs) :
\end{definition}

\begin{aretenir}[Les 4 codes couleur ISO 3864 et leurs pictogrammes]
\begin{center}
\begin{tabularx}{\linewidth}{|c|X|X|X|}\hline
\rowcolor{popblueL}
\textbf{Couleur} & \textbf{Sens} & \textbf{Forme géométrique} & \textbf{Exemple} \\ \hline
\textcolor{popred}{\textbf{Rouge}} & \textbf{Interdiction / Arrêt / Danger immédiat} & Cercle barré (diagonale 45°) + Pictogramme noir & \og Ne pas courir \fg, \og Interdiction de fumer \fg \\ \hline
\textcolor{popblue}{\textbf{Bleu}} & \textbf{Obligation / Port EPI} & Cercle plein + Pictogramme blanc & \og Port du casque obligatoire \fg, \og Port des gants obligatoire \fg \\ \hline
\textcolor{popyellow}{\textbf{Jaune}} & \textbf{Avertissement / Danger / Attention} & Triangle pointe en haut + Pictogramme noir & \og Attention danger électrique \fg, \og Risque de chute \fg \\ \hline
\textcolor{popgreen}{\textbf{Vert}} & \textbf{Secours / Sauvetage / Issue} & Carré/Rectangle + Pictogramme blanc & \og Issue de secours \fg, \og Poste de premiers secours \fg \\ \hline
\end{tabularx}
\end{center}
\end{aretenir}

\begin{methode}[Analyser / Produire une affiche de sécurité]
\begin{enumerate}
\item \textbf{Identifier le type} (Interdiction, Obligation, Avertissement, Secours) $\rightarrow$ déduit la couleur et la forme.
\item \textbf{Décrire le pictogramme} : Symbole graphique simple, universel, centré.
\item \textbf{Formuler la consigne} : \textbf{Infinitif} (règle générale, impersonnelle, intemporelle), courte, affirmative (obligation) ou négative (interdiction), lieu précisé si besoin.
\item \textbf{Justifier l'efficacité} : Redondance codée (Image attire + Couleur alerte + Texte précise) = Communication multimodale fiable.
\end{enumerate}
\end{methode}

\courssub{Production et amélioration d'un rapport complet}

\begin{methode}[Démarche complète de rédaction (7 étapes)]
\begin{enumerate}
\item \textbf{Collecte \& Tri :} Notes de terrain, photos, mesures, témoignages, documents techniques. Trier par chronologie / thème.
\item \textbf{Planification :} Construire le plan détaillé (Intro, 2-3 parties corps, Conclu) avec idées forces et connecteurs.
\item \textbf{Rédaction de l'en-tête :} Remplir rigoureusement la structure externe.
\item \textbf{Rédaction du corps :} Une partie après l'autre. Style impersonnel, passé/imparfait, connecteurs, lexique spécialisé \textbf{défini}.
\item \textbf{Rédaction Intro / Conclu :} Respecter les 3 mouvements chacun. Vérifier l'adéquation Intro $\leftrightarrow$ Plan $\leftrightarrow$ Conclu.
\item \textbf{Relecture croisée (Check-list) :}
\begin{itemize}
\item [$\square$] Structure externe complète (Objet, Date, Signature, PJ) ?
\item [$\square$] Objectivité totale (pas de « je », pas d'adjectifs subjectifs) ?
\item [$\square$] Chronologie / Logique respectée ? Connecteurs présents ?
\item [$\square$] Lexique spécialisé défini (3 minimum au Bac) ?
\item [$\square$] Chiffres, unités (\SI{...}{...}), localisations précis ?
\item [$\square$] Conclusion sans fait nouveau ? Recommandations à l'infinitif ?
\item [$\square$] Orthographe, grammaire, mise en page (alinéas, titres) ?
\end{itemize}
\item \textbf{Finalisation :} Signature, Visa, PJ, Diffusion. Impression / Numérisation.
\end{enumerate}
\end{methode}

\begin{attention}[Erreurs fréquentes au Probatoire / Bac Tech]
\begin{itemize}
\item Oublier l'\textbf{Objet} ou la \textbf{Date} en en-tête (0 pt structure externe).
\item Écrire « Je pense que... », « C'était bien », « Heureusement » (subjectivité = sanction).
\item Introduire des faits nouveaux en conclusion (incohérence structurelle).
\item Utiliser le futur ou le conditionnel pour narrer des faits passés.
\item Oublier de définir les 3 termes spécialisés exigés (perte de 3 pts faciles).
\item Consignes de sécurité au lieu de l'infinitif (ex. « Vous devez porter... » au lieu de « Port du casque obligatoire »).
\end{itemize}
\end{attention}

% ============================================================
% EXERCICES ET CORRIGÉS
% ============================================================

\courssub{Je vérifie mes connaissances}

\begin{corrige}[Exercice 1 — QCM : les caractéristiques du rapport]
\begin{enumerate}
\item \textbf{b} — Le rapport est un texte fonctionnel qui rend compte de \textbf{faits réels observés ou vécus} (stage, visite, incident, expérience). Il exclut l'imaginaire, l'opinion pure et le rêve.
\item \textbf{b} — L'introduction doit obligatoirement présenter le \textbf{contexte} (cadre, date, commanditaire), l'\textbf{objet} (sujet précis) et annoncer le \textbf{plan} (articulation des parties). La bibliographie, les remerciements et la signature ne figurent pas en introduction.
\item \textbf{b} — Le temps dominant est le \textbf{passé composé} (actions ponctuelles, achevées, datées : \og Le technicien a constaté \fg) ou l'\textbf{imparfait} (états, descriptions, actions habituelles : \og La machine fonctionnait normalement \fg). Le futur, le subjonctif et l'impératif ne sont pas des temps de narration factuelle.
\item \textbf{b} — Le \textbf{lexique spécialisé} (ou vocabulaire technique) regroupe les termes propres à un domaine professionnel (BTP, mécanique, informatique, santé...), distincts du lexique commun partagé par tous.
\end{enumerate}
\end{corrige}

\begin{corrige}[Exercice 2 — Vrai ou faux]
\begin{enumerate}
\item \textbf{Faux.} La conclusion \textbf{synthétise} les résultats, tire les enseignements et formule des recommandations ; elle \textbf{n'introduit jamais de faits nouveaux} (ni données, ni événements non traités dans le corps).
\item \textbf{Vrai.} L'\textbf{objectivité} est la règle d'or du rapport : on utilise des formes impersonnelles (\og Il a été constaté que...\fg), le pronom \og on \fg ou la voix passive, et on proscrit \og je pense que \fg, \og à mon avis \fg, \og heureusement \fg, etc.
\item \textbf{Faux.} L'\textbf{objet} est une mention \textbf{obligatoire} de la structure externe (en-tête). Il permet au destinataire d'identifier immédiatement le sujet du document.
\item \textbf{Vrai.} Une affiche de sécurité normalisée (ISO 7010) combine trois composantes : un \textbf{pictogramme} (symbole graphique), un \textbf{code couleur} (rouge = interdiction/arrêt, bleu = obligation, jaune = avertissement, vert = secours) et une \textbf{consigne textuelle} courte à l'infinitif.
\end{enumerate}
\end{corrige}

\begin{corrige}[Exercice 3 — Définitions]
\begin{enumerate}
\item \textbf{Le rapport :} C'est un texte fonctionnel, objectif et structuré, qui rend compte par écrit de faits réels (observés, vécus, mesurés) dans un cadre professionnel ou scolaire, à destination d'un destinataire précis, en vue d'informer, de justifier ou d'aider à la décision.
\item \textbf{Le texte explicatif :} C'est un type de texte qui répond à la question \og Comment ? \fg ou \og Pourquoi ? \fg en décrivant un phénomène, un mécanisme, une procédure ou un processus. Il utilise des connecteurs logiques (cause, conséquence, chronologie, finalité) et un lexique précis, souvent spécialisé.
\item \textbf{Le lexique commun :} C'est l'ensemble des mots du langage courant, compris et utilisés par l'ensemble des locuteurs d'une langue dans la vie quotidienne, sans appartenir à un domaine technique ou scientifique spécifique (ex. : \og maison \fg, \og manger \fg, \og grand \fg).
\item \textbf{La structure externe d'un rapport :} C'est l'ensemble des mentions d'identification obligatoires placées en en-tête (rédacteur, destinataire, lieu, date, objet, références) et en pied de page (signature, visa, pièces jointes, diffusion) qui encadrent le contenu du rapport.
\end{enumerate}
\end{corrige}

\begin{corrige}[Exercice 4 — Repérage rapide]
\begin{enumerate}
\item \textbf{Éléments d'objectivité :}
\begin{itemize}
\item L'usage de la troisième personne (\og Le technicien \fg) et non du \og je \fg, ce qui efface le sujet énonciateur.
\item Les verbes d'action factuelle (\og a constaté \fg, \og a fermé \fg) sans adjectifs qualificatifs subjectifs (pas \og grave \fg, \og inquiétante \fg).
\item La précision spatio-temporelle (\og Le 12 mars \fg, \og au niveau de la canalisation principale \fg) et technique (\og vanne d'arrêt \fg).
\end{itemize}
\item \textbf{Temps :} \textbf{Passé composé} (\og a constaté \fg, \og a fermé \fg). \textbf{Justification :} Ce sont des actions ponctuelles, brèves, achevées et datées (\og Le 12 mars \fg) qui font avancer la narration chronologique des faits.
\item \textbf{Objet proposé :} \og Constat et intervention sur fuite de la canalisation principale — 12 mars \fg (ou \og Rapport d'incident : fuite canalisation principale \fg).
\end{enumerate}
\end{corrige}

\courssub{Je m'entraîne}

\begin{corrige}[Exercice 5 — Lexique commun ou spécialisé ?]
\begin{enumerate}
\item \textbf{Classement :}
\begin{center}
\begin{tabularx}{\linewidth}{|l|X|X|}\hline
\rowcolor{popblueL}
\textbf{Lexique commun} & \textbf{Lexique spécialisé (BTP)} \\ \hline
maison, mur, toit, fenêtre, porte, chambre, gravats, plan (sens courant) & béton, coffrage, ferraillage, mortier, étaiement, échafaudage, dalle \\ \hline
\end{tabularx}
\end{center}
\item \textbf{Définitions de trois termes spécialisés :}
\begin{itemize}
\item \textbf{Coffrage :} Structure temporaire (bois, métal) servant de moule pour maintenir le béton frais en place jusqu'à son durcissement et lui donner sa forme définitive.
\item \textbf{Ferraillage :} Ensemble de barres d'acier (armatures) disposées à l'intérieur du béton pour absorber les efforts de traction que le béton ne supporte pas seul.
\item \textbf{Étaiement :} Dispositif de soutien provisoire (étais, tours d'étayage) qui porte la charge d'un ouvrage (dalle, poutre) pendant la prise du béton, avant qu'il n'atteigne sa résistance mécanique.
\end{itemize}
\item \textbf{Intérêt du lexique spécialisé :} Il garantit la \textbf{précision technique} (un seul mot pour une réalité complexe), la \textbf{concision} (évite les périphrases), l'\textbf{univocité} (pas d'ambiguïté entre professionnels) et marque l'\textbf{appartenance professionnelle} (identité du métier). Il transforme le rapport en outil de travail fiable pour les décideurs et les pairs.
\end{enumerate}
\end{corrige}

\begin{corrige}[Exercice 6 — De la subjectivité à l'objectivité]
\begin{enumerate}
\item \textbf{Marques de subjectivité et de jugement :}
\begin{itemize}
\item \og Franchement \fg (marqueur énonciatif familier), \og je trouve que \fg (opinion personnelle).
\item \og horrible \fg, \og pauvre Jean \fg (jugement de valeur, pathos, compassion).
\item \og j'ai eu très peur \fg (émotion personnelle, sans intérêt factuel).
\item \og c'était vraiment de sa faute \fg (jugement moral, accusation sans preuve).
\item \og il ne fait jamais attention \fg (généralisation abusive, \og jamais \fg).
\end{itemize}
\item \textbf{Réécriture objective (style rapport) :}
\og Le 15 octobre 2025, vers 10 h 15, l'apprenti Jean Dupont, affecté au poste de sciage, s'est sectionné l'index de la main droite au contact de la lame de la scie circulaire stationnaire n° 3. Le témoin direct, M. Kouam, a actionné l'arrêt d'urgence et alerté l'infirmier de l'atelier. Les premiers secours (compression, pansement stérile) ont été prodigués dans les deux minutes. L'apprenti a été évacué vers le centre de santé du district pour suture. L'analyse préliminaire révèle l'absence de port des gants anti-coupure (EPI obligatoire) et le non-respect de la distance de sécurité par rapport à la zone de coupe. \fg
\item \textbf{Recommandation de sécurité :} \og Il est impératif de porter systématiquement les gants anti-coupure et de maintenir les deux mains hors de la zone de danger lors de l'utilisation de la scie circulaire. \fg
\end{enumerate}
\end{corrige}

\begin{corrige}[Exercice 7 — Analyse d'une affiche de sécurité]
\begin{enumerate}
\item \textbf{Description des trois composantes :}
\begin{itemize}
\item \textbf{Pictogramme :} Silhouette humaine en position de course, inscrite dans un cercle, barrée d'une diagonale rouge (symbole d'interdiction ISO 7010).
\item \textbf{Couleur :} \textbf{Rouge} (cercle et barre) sur fond blanc. Le rouge code l'interdiction, l'arrêt, le danger immédiat.
\item \textbf{Consigne :} \og Ne pas courir dans l'atelier \fg (verbe à l'infinitif négatif, phrase courte, impérative).
\end{itemize}
\item \textbf{Rôle de chaque composante :}
\begin{itemize}
\item \textbf{Pictogramme :} Assure une \textbf{compréhension immédiate et universelle} (translingue, accessible aux non-lecteurs, aux malvoyants, dans le bruit). Il montre \textbf{l'action interdite}.
\item \textbf{Couleur rouge :} Déclenche une \textbf{alerte visuelle forte} (salience perceptive), signale \textbf{l'interdiction formelle} (code couleur normalisé : rouge = stop/interdiction).
\item \textbf{Consigne textuelle :} \textbf{Ancre le sens} (fonction de relais), lève toute ambiguïté sur l'action (\og courir \fg) et le lieu (\og dans l'atelier \fg), donne la forme linguistique de l'obligation (infinitif = consigne générale).
\end{itemize}
\item \textbf{Paragraphe explicatif (fonctionnement de l'affiche) :}
Cette affiche relève de la signalisation de sécurité normalisée (ISO 3864 / ISO 7010). Son fonctionnement repose sur la redondance codée de trois signaux. Le \textbf{pictogramme} (silhouette courant barrée) figure l'action prohibée de manière iconique, compréhensible sans lecture. Le \textbf{code couleur} (rouge sur blanc) active une réponse attentionnelle automatique : le rouge signale universellement l'interdiction et le danger, conditionnant une obéissance rapide. La \textbf{consigne verbale} (\og Ne pas courir dans l'atelier \fg) explicite le message : l'infinitif négatif formule une règle générale, impersonnelle et intemporelle ; la préposition \og dans \fg circonscrit l'espace concerné. L'ensemble forme un \textbf{message multimodal} où l'image attire, la couleur alerte, le texte précise. Cette redondance garantit l'efficacité communicationnelle même en situation de stress, de bruit ou de barrière linguistique, transformant l'affiche en un \textbf{dispositif de prévention primaire} fiable.
\end{enumerate}
\end{corrige}

\begin{corrige}[Exercice 8 — Remédiation : corriger un rapport]
\begin{enumerate}
\item \textbf{Erreurs, règles violées et corrections :}
\begin{center}
\begin{tabularx}{\linewidth}{|l|X|X|}\hline
\rowcolor{popblueL}
\textbf{Erreur} & \textbf{Règle violée} & \textbf{Correction concrète} \\ \hline
Objet absent & Structure externe : l'objet est obligatoire en en-tête. & Ajouter : \og Objet : Rapport de visite de l'usine de transformation de manioc de [Localité] — Stage du [dates] \fg. \\ \hline
Date manquante & Structure externe : la date de rédaction est obligatoire. & Ajouter en en-tête : \og Fait à [Ville], le [Date du jour] \fg. \\ \hline
Faits nouveaux en conclusion & Conclusion : synthèse uniquement, pas de nouveaux faits. & Supprimer les faits nouveaux de la conclusion ; les intégrer dans le corps (partie activités) ou les supprimer. \\ \hline
Deux termes de jargon non définis & Lisibilité : tout terme spécialisé doit être défini à sa 1ère occurrence. & Ajouter définition entre parenthèses ou note : \og broyeur (machine qui réduit le manioc en pâte) \fg, \og four rotatif (four tournant pour cuisson/séchage) \fg. \\ \hline
Connecteurs logiques absents & Cohérence textuelle : articuler les idées (chronologie, cause, conséquence). & Insérer : \og D'abord \fg, \og Ensuite \fg, \og Par conséquent \fg, \og Cependant \fg, \og En outre \fg entre paragraphes et phrases. \\ \hline
Marques de subjectivité & Objectivité : style impersonnel, pas d'opinion. & Remplacer \og Je pense que c'est bien \fg par \og L'organisation s'avère efficace \fg ; supprimer adjectifs évaluatifs (\og super \fg, \og nul \fg). \\ \hline
Paragraphes non articulés & Structure interne : 1 idée = 1 paragraphe, liens logiques. & Restructurer en parties distinctes (Introduction, 1. Présentation usine, 2. Procédé transformation, 3. Conditions travail, Conclusion) avec alinéas et amorces. \\ \hline
Signature oubliée & Validité juridique/administrative : signature obligatoire. & Ajouter en bas à droite : \og Signature : [Prénom Nom], Élève 1re T [Spécialité] \fg. \\ \hline
\end{tabularx}
\end{center}
\item \textbf{Objet corrigé :} \og Objet : Rapport de visite de l'usine de transformation de manioc de [Localité] — Stage d'observation du [date début] au [date fin] \fg
\item \textbf{Conclusion corrigée :}
\og Cette visite a permis de comprendre le processus industriel complet de transformation du manioc en farine (gari) et en amidon, depuis la réception et le lavage des tubercules jusqu'à l'ensachage du produit fini. Les acquis techniques majeurs portent sur la conduite du four rotatif, la maintenance de premier niveau du broyeur et le respect strict des normes HSE (Hygiène, Sécurité, Environnement). Le stage a également développé des compétences transverses : travail en équipe 3x8, rigueur des consignes, communication professionnelle. Nous adressons nos vifs remerciements à M. [Nom], Directeur de l'usine, à M. [Nom], Chef de production (tuteur), et à l'ensemble du personnel pour leur accueil, leur disponibilité pédagogique et la qualité de l'encadrement. \fg
\end{enumerate}
\end{corrige}

\courssub{Je me prépare au Bac}

\begin{corrige}[Exercice 9 — Production écrite : rapport de stage]
\textbf{Barème indicatif (sur 20) :} Structure externe (3 pts) — Introduction (3 pts) — Corps : 3 parties articulées, lexique spécialisé, objectivité (8 pts) — Conclusion (3 pts) — 3 termes définis (3 pts). \textbf{Points de vigilance :} Respect scrupuleux de la structure externe ; introduction = contexte + objectifs + plan annoncé ; corps = 3 parties équilibrées, connecteurs logiques variés, \textbf{zéro subjectivité} ; conclusion = bilan + acquis + remerciements \textbf{sans fait nouveau} ; 3 termes spécialisés \textbf{soulignés} et définis (note ou parenthèse).
\end{corrige}

\begin{exemplebox}[Modèle de réponse — Rapport de stage (Garage automobile)]
\textbf{Structure externe :}
\begin{tabularx}{\linewidth}{lX}
\textbf{Nom du rédacteur :} & MBARGA Jean-Paul \\
\textbf{Classe :} & 1re T Maintenance des Véhicules Automobiles (MVA) \\
\textbf{Destinataire :} & Monsieur le Professeur de Technologie MVA, Lycée Technique de [Ville] \\
\textbf{Lieu :} & [Ville] \\
\textbf{Date :} & 15 juin 2026 \\
\textbf{Objet :} & Rapport de stage d'observation effectué du 2 au 13 juin 2026 au Garage \og Auto-Service \fg, [Ville]
\end{tabularx}

\textbf{Introduction :}
Dans le cadre de la formation en Première Technicien Maintenance des Véhicules Automobiles, j'ai effectué un stage d'observation de deux semaines, du 2 au 13 juin 2026, au garage \og Auto-Service \fg, spécialisé dans la réparation multimarque et l'entretien courant. Ce stage visait trois objectifs : confronter les connaissances théoriques acquises en atelier pédagogique à la réalité professionnelle, découvrir l'organisation du travail en entreprise (gestion des rendez-vous, relation client, gestion des stocks) et observer les interventions sur systèmes électroniques embarqués (diagnostic valise, injection, ABS/ESP). Ce rapport présente successivement la structure de l'entreprise, les activités techniques observées et réalisées, puis les difficultés rencontrées et les solutions apportées.

\textbf{Corps du rapport :}

\textbf{1. Présentation de l'entreprise}
\textbf{D'abord,} le garage \og Auto-Service \fg, situé zone industrielle de [Ville], emploie 12 personnes : un gérant, un chef d'atelier, 5 techniciens (dont 2 spécialisés électricité/électronique), 2 apprentis, 1 magasinier et 1 réceptionnaire. L'atelier dispose de 4 ponts élévateurs, d'un banc de géométrie, d'une station de climatisation et de 3 valises de diagnostic multimarque (Bosch, Delphi, Launch). L'organisation suit le processus standard : réception client $\rightarrow$ ordre de réparation (OR) $\rightarrow$ diagnostic $\rightarrow$ devis $\rightarrow$ intervention $\rightarrow$ contrôle qualité $\rightarrow$ facturation. Le magasin gère un stock de \SI{1200}{références} (filtres, huiles, plaquettes, capteurs).

\textbf{2. Activités observées et réalisées}
\textbf{Ensuite,} j'ai participé activement aux opérations d'entretien courant : vidanges (huile 5W30, 5W40), remplacement de filtres à huile/air/habitacle, contrôle des niveaux (LDR, lave-glace, direction assistée), vérification pression pneumatiques. J'ai assisté à des interventions plus complexes : \textbf{diagnostic électronique} (lecture/éffacement codes défaut sur calculateur moteur, boîtier ABS, airbag) via valise \textbf{Bosch KTS 570} (\emph{outil de diagnostic multimarque permettant la lecture des paramètres temps réel et l'actionnement d'actionneurs}) ; remplacement d'un \textbf{kit distribution + pompe à eau} sur moteur 1.5 dCi (respect des repères de calage, purge circuit refroidissement) ; \textbf{réfection freinage} (disques, plaquettes, purge circuit hydraulique, rodage). J'ai également géré deux réceptions clients : écoute de la plainte, rédaction de l'OR, essai routier pour reproduire le symptôme.

\textbf{3. Difficultés rencontrées}
\textbf{Cependant,} trois difficultés majeures sont apparues. \textbf{Premièrement,} la rapidité d'exécution exigée (temps barème constructeur) contraste avec le rythme scolaire : j'ai dû m'adapter en préparant outils et pièces en amont (\og mise en place \fg). \textbf{Deuxièmement,} le diagnostic d'une panne intermittente (voyant moteur allumé aléatoirement sur Peugeot 308) a nécessité une analyse approfondie des \textbf{paramètres temps réel} (\emph{valeurs instantanées des capteurs : pression rail, débitmètre, température eau}) sur plusieurs cycles de conduite, sans reproduction du défaut en atelier. \textbf{Troisièmement,} la communication avec des clients non techniques a parfois été délicate pour expliquer la nature d'une réparation (ex. : \og pourquoi changer le kit distribution si la courroie n'est pas cassée ? \fg). Mon tuteur m'a conseillé d'utiliser des analogies simples et de montrer les pièces usées.

\textbf{Conclusion :}
Ce stage a été une expérience formatrice déterminante. Sur le plan technique, j'ai consolidé mes gestes professionnels (serrage au couple, purge freinage, diagnostic valise) et découvert la réalité du diagnostic électronique embarqué, bien plus complexe que les cas d'école. Sur le plan humain, j'ai appris la rigueur de la gestion du temps, l'importance de la propreté du poste de travail et la relation de confiance avec le client. Les acquis (autonomie partielle sur entretien courant, initiation diagnostic valise, respect procédures sécurité) confirment mon orientation vers le BTS Maintenance des Véhicules. Je remercie chaleureusement M. FOUDJI, gérant, M. NGONO, chef d'atelier (mon tuteur), et toute l'équipe pour leur patience, leur pédagogie et la confiance qu'ils m'ont accordée.

\textbf{Signature :} MBARGA Jean-Paul

\textbf{Trois termes de lexique spécialisé (soulignés dans le texte) :}
\begin{enumerate}
\item \textbf{Valise de diagnostic} (souligné) : Boîtier électronique portable (ex. Bosch KTS 570) qui se connecte à la prise OBD du véhicule pour dialoguer avec les calculateurs (lecture codes défaut, paramètres temps réel, actionnement actionneurs, codage).
\item \textbf{Kit distribution} (souligné) : Ensemble comprenant la courroie de distribution, les galets tendeurs et enrouleurs, et souvent la pompe à eau ; assure la synchronisation vilebrequin/arbre à cames. Remplacement préventif selon préconisations constructeur (km ou années).
\item \textbf{Paramètres temps réel} (souligné) : Données instantanées transmises par les capteurs au calculateur (pression rail injection, débit d'air, température eau, position papillon, lambda) ; indispensables pour diagnostiquer une panne sans code défaut stocké.
\end{enumerate}
\end{exemplebox}

\begin{corrige}[Exercice 10 — Texte explicatif et rapport d'incident]
\textbf{Barème indicatif (sur 20) :} Structure externe (3 pts) — Corps : faits (heure, lieu, causes), mesures, dégâts — objectivité, chronologie, lexique technique (6 pts) — Paragraphe explicatif extincteur : connecteurs d'explication, exactitude technique (5 pts) — Conclusion : 2 consignes à l'infinitif, pertinentes (3 pts) — Qualité de langue, mise en page (3 pts). \textbf{Points de vigilance :} En-tête complet (délégué, proviseur, date, objet) ; corps = narration factuelle au passé (passé composé/imparfait), pas d'hypothèse gratuite sur causes (\og causes probables \fg) ; paragraphe explicatif = \textbf{d'abord, ensuite, enfin, c'est-à-dire} obligatoires ; conclusion = \textbf{consignes à l'infinitif} (style affiche).
\end{corrige}

\begin{exemplebox}[Modèle de réponse — Rapport d'incendie laboratoire]
\textbf{Structure externe :}
\begin{tabularx}{\linewidth}{lX}
\textbf{Rédacteur :} & KAMGA Pierre, Délégué de la classe de 1re T Sciences Physiques \\
\textbf{Destinataire :} & Monsieur le Proviseur, Lycée Technique de [Ville] \\
\textbf{Lieu :} & [Ville] \\
\textbf{Date :} & 20 juin 2026 \\
\textbf{Objet :} & Rapport d'incident — Début d'incendie au laboratoire de Sciences Physiques (Salle P3) le 20/06/2026
\end{tabularx}

\textbf{Corps du rapport :}

\textbf{1. Déroulement des faits}
Le 20 juin 2026, à 14 h 15, lors de la séance de TP de chimie (manipulation : distillation de l'éthanol), un début d'incendie s'est déclaré sur le poste de travail n° 4 (binôme : NGOUMOU / ESSOMBA). \textbf{D'après les témoins}, l'origine probable est le contact entre la flamme du bec Bunsen (mal réglé, flamme jaune instable) et des vapeurs d'éthanol échappées d'un ballon de distillation mal emboîté (joint défectueux). L'incendie a débuté vers 14 h 15 min 30 s sur la paillasse en bois verni, se propageant rapidement aux papiers de brouillon et au manchon en caoutchouc du condenseur. Aucune manipulation de produits toxiques (acide concentré, solvant chloré) n'était en cours sur ce poste.

\textbf{2. Mesures prises immédiatement}
\textbf{D'abord,} l'élève NGOUMOU a actionné le \textbf{manœuvre d'arrêt d'urgence gaz} (vanne quart de tour murale) à 14 h 15 min 40 s, coupant l'alimentation de tous les becs Bunsen de la salle. \textbf{Ensuite,} l'enseignant, M. TCHINDA, a saisi l'\textbf{extincteur à poudre ABC} (positionné à l'entrée, conforme NF EN 3) situé à 3 m du foyer, a retiré la goupille, dirigé la lance vers la base des flammes et déclenché la projection. \textbf{Enfin,} la classe a été évacuée calmement vers le point de rassemblement (cour arrière) selon le plan d'évacuation affiché ; l'appel a confirmé la présence de tous les 28 élèves. Les secours (sapeurs-pompiers) ont été alertés par le standardiste à 14 h 18 ; ils sont arrivés à 14 h 25, ont vérifié l'extinction complète et aéré le local.

\textbf{3. Dégâts constatés}
Les dégâts se limitent à la paillasse du poste n° 4 (surface brûlée $\approx$ 40 cm x 30 cm, bois carbonisé), au manchon du condenseur (fondu), aux cahiers de TP des deux élèves (brûlés) et à la projection de poudre ABC sur \SI{4}{\square\meter} environ (paillasses, sol, matériel verrerie). \textbf{Aucun blessé} n'est à déplorer (ni brûlure, ni inhalation). Le matériel verrerie (ballons, condenseur, béchers) est récupérable après nettoyage. Le détecteur de fumée de la salle a fonctionné normalement (alarme déclenchée à 14 h 16).

\textbf{4. Paragraphe explicatif : fonctionnement de l'extincteur à poudre ABC utilisé}
L'extincteur utilisé est un modèle \textbf{à poudre polyvalente ABC (phosphate d'ammonium)}, efficace sur feux de classe A (solides), B (liquides) et C (gaz). \textbf{D'abord,} on retire la \textbf{goupille de sécurité} qui verrouille la poignée de commande, ce qui arme l'appareil. \textbf{Ensuite,} on saisit la lance (ou le tuyau) et on la dirige \textbf{vers la base des flammes} (siège du feu), non pas vers le haut du panache. \textbf{Enfin,} on appuie fermement sur la \textbf{poignée de commande} (ou le levier) : cela ouvre la vanne interne, la pression du gaz propulseur (azote ou air comprimé) chasse la poudre hors du réservoir à travers le tube plongeur et la lance. \textbf{C'est-à-dire} que la poudre forme un nuage qui \textbf{étouffe le feu par inhibition chimique} : les particules fines interceptent les radicaux libres (H*, OH*) de la réaction en chaîne de la combustion, stoppant l'exothermicité, tout en formant une barrière isolante sur le combustible.

\textbf{Conclusion — Recommandations de sécurité (consignes affichage) :}
\begin{itemize}
\item \textbf{Vérifier systématiquement l'étanchéité des montages verrerie (emboîtements, joints) avant tout allumage de bec Bunsen.}
\item \textbf{Maintenir les issues de secours et l'accès aux extincteurs dégagés en permanence.}
\end{itemize}

\textbf{Signature :} KAMGA Pierre, Délégué 1re T SP
\end{exemplebox}

\begin{corrige}[Exercice 11 — Réception et production : visite d'entreprise]
\textbf{Barème indicatif (sur 20) :} Rapport complet structuré (ext/intro/corps/concl) 25-30 lignes (6 pts) — 4 termes spécialisés définis (4 pts) — Affiche sécurité : pictogramme décrit, couleur justifiée, consigne infinitif, justification 3 lignes (5 pts) — Paragraphe objectivité/fonctionnalité (5 pts). \textbf{Points de vigilance :} Rapport = style impersonnel, passé composé/imparfait, connecteurs, chiffres précis (dates, effectifs, production) ; termes spécialisés = définis pour non-spécialiste ; affiche = code couleur ISO (bleu obligation) ; paragraphe final = méta-langage (analyse du propre texte).
\end{corrige}

\begin{exemplebox}[Modèle de réponse — Visite cimenterie]
\textbf{1. Rapport de visite complet :}

\textbf{Structure externe :}
\begin{tabularx}{\linewidth}{lX}
\textbf{Rédacteur :} & ETOUNDI Marie, Élève 1re T Génie Civil \\
\textbf{Destinataire :} & Monsieur le Professeur de Technologie BTP, Lycée Technique de [Ville] \\
\textbf{Lieu :} & [Ville] \\
\textbf{Date :} & 10 juin 2026 \\
\textbf{Objet :} & Rapport de visite pédagogique à la Cimenterie du Littoral (CIMLIT) — 5 juin 2026
\end{tabularx}

\textbf{Introduction :}
Dans le cadre du module \og Matériaux de construction \fg, la classe de 1re T Génie Civil a effectué une visite pédagogique le 5 juin 2026 à la Cimenterie du Littoral (CIMLIT), située à [Localité], principal producteur de ciment CEM II 42,5 de la région. Accueillis par M. DJEUKAM, Chef de production, nous visions à observer le processus industriel complet de fabrication du ciment, de l'extraction du calcaire à l'ensachage, et à identifier les mesures de prévention des risques industriels. Ce rapport décrit successivement le processus de fabrication observé et les conditions de sécurité en milieu cimentier.

\textbf{Corps du rapport :}

\textbf{1. Processus de fabrication du ciment}
\textbf{D'abord,} la visite a débuté par la carrière de calcaire (extraction à l'explosif, chargement par pelles hydrauliques \SI{120}{t}, transport par tombereaux \SI{60}{t} vers le concasseur primaire). \textbf{Ensuite,} nous avons suivi la ligne de production : \textbf{broyeur vertical} (\emph{moulin à galets qui réduit le cru en poudre fine « farine » sous \SI{90}{\micro\meter}}) pour le mélange calcaire/argile/sable correctif ; \textbf{silos de stockage/homogénéisation} (capacité \SI{15000}{t}, air fluidifiant pour éviter la ségrégation) ; \textbf{four rotatif} (\emph{four incliné tournant 3-4 tr/min, longueur \SI{60}{m}, température de cuisson \SI{1450}{\celsius}, brûleur multi-combustibles : charbon/pétcoke/CSR}) où la farine se transforme en \textbf{clinker} (nodules vitrifiés) ; \textbf{refroidisseur à grilles} (récupération chaleur pour préchauffage) ; \textbf{broyeur à boulets ciment} (clinker + 5\% gypse + additions calcaires) ; \textbf{ensachage} automatique (sacs \SI{50}{kg}, paletteuse, stockage expédition). La production journalière atteint \textbf{\SI{5000}{sacs}} (\SI{250}{t/jour}) avec \SI{200}{ouvriers} (3x8).

\textbf{2. Sécurité et conditions de travail}
\textbf{Par ailleurs,} la prévention des risques est omniprésente : \textbf{port du casque obligatoire} (zone dure), protections auditives (bruit > \SI{85}{\decibel} au broyeur/four), masques P3 (poussières de ciment/silice), gants, chaussures de sécurité. La signalisation (panneaux ISO 7010, marquage au sol, chemins piétons) est rigoureuse. Nous avons noté la présence de \textbf{dépoussiéreurs à manches} (filtration gaz four/broyeur, rejet < \SI{30}{\milli\gram\per\cubic\meter}) et de bassins de décantation eaux pluviales. La visite s'est terminée par une séance questions-réponses sur la maintenance préventive (arrêt annuel 3 semaines) et la qualité (laboratoire interne : résistance 2/28 j, prise de vue, chimie).

\textbf{Conclusion :}
Cette immersion industrielle a concrétisé nos cours sur les liants hydrauliques : la maîtrise de la chimie du clinker (modules de saturation, silice, alumine), la thermodynamique du four, et l'enjeu énergétique/environnemental (substitution combustibles, valorisation chaleur fatale). Les acquis techniques (process continu, automatismes, instrumentation) et comportementaux (rigueur EPI, respect consignes, travail en équipe) sont précieux pour notre future insertion. Nous remercions la Direction de CIMLIT, M. DJEUKAM et les encadrants pour la qualité de l'accueil et la transparence technique.

\textbf{Signature :} ETOUNDI Marie

\textbf{2. Quatre termes de lexique spécialisé définis :}
\begin{enumerate}
\item \textbf{Broyeur vertical (ou moulin à galets) :} Machine qui broie le cru (mélange calcaire/argile) par compression entre une table tournante et 2 à 4 galets presseurs ; plus économe en énergie qu'un broyeur à boulets.
\item \textbf{Four rotatif :} Cœur de la cimenterie : tube en acier incliné (3-4\%), revêtu de briques réfractaires, tournant sur lui-même (3-4 tr/min) ; la farine y descend par gravité pendant que la flamme (\SI{1450}{\celsius}) remonte en contre-courant, réalisant la clinkérisation (réactions chimiques solides).
\item \textbf{Clinker :} Produit intermédiaire du four, nodules gris foncé (1-3 cm) résultant de la cuisson à \SI{1450}{\celsius} ; composé de 4 phases minérales principales (alite, bélite, aluminate, ferro-aluminate). Le ciment = clinker broyé + gypse (+ additions).
\item \textbf{Ensachage :} Poste final automatisé : la poudre de ciment est dosée (\SI{50}{kg} $\pm$ 0,5\%), mise en sacs papier kraft (souvent 3 plis, valve), les sacs sont fermés, palettisés (\SI{1.5}{t/palette}), filmés et stockés avant expédition camion/wagon.
\end{enumerate}

\textbf{3. Affiche de sécurité pour visiteurs de la cimenterie :}
\begin{itemize}
\item \textbf{Pictogramme :} Tête humaine de profil portant un casque de chantier (blanc), inscrite dans un \textbf{cercle bleu} (fond blanc) — symbole ISO 7010 M002 « Port du casque obligatoire ».
\item \textbf{Couleur :} \textbf{Bleu} (RAL 5005 / Bleu sécurité) pour le cercle ; fond blanc. Le bleu code l'\textbf{obligation} (contrairement au rouge = interdiction).
\item \textbf{Consigne :} \og Port du casque obligatoire \fg (infinitif, impersonnel, court).
\end{itemize}
\textbf{Justification (3 lignes) :} Le bleu signale universellement une obligation (ISO 3864) ; le pictogramme casque identifie sans ambiguïté l'EPI requis contre le risque de chute d'objets (grenailles, poussière, outils) omniprésent en cimenterie ; l'infinitif formule une règle générale, impérative et intemporelle, compréhensible instantanément par tout visiteur, même non francophone.

\textbf{4. Objectivité et fonctionnalité du rapport rédigé :}
Le rapport produit est \textbf{objectif} car il relate des faits datés (5 juin), chiffrés (\SI{5000}{sacs/jour}, \SI{200}{ouvriers}, \SI{1450}{\celsius}, \SI{85}{\decibel}, \SI{30}{\milli\gram\per\cubic\meter}), localisés (carrière, broyeur, four, ensachage), sans aucun jugement de valeur, adjectif subjectif (\og beau \fg, \og sale \fg) ni marque de la première personne (\og j'ai vu \fg remplacé par \og la visite a débuté \fg, \og nous avons suivi \fg). Il est \textbf{fonctionnel} car il répond à une commande précise (compte rendu pédagogique), structuré selon la norme (externe/interne), utilisable par le destinataire (professeur) pour évaluer les acquis, archiver la sortie, et servir de base à un exposé ou un examen ; son lexique spécialisé défini le rend accessible tout en restant techniquement rigoureux.
\end{exemplebox}

\newpage

\coursec{Fiche bilan}

\begin{popfigure}
\begin{center}\tcbox[colback=popblue,colframe=popblue,coltext=white,arc=9pt,boxsep=4pt,left=6pt,right=6pt]{\bfseries\small Produire un rapport}\end{center}
\vspace{-2pt}
\begin{tcbraster}[raster columns=3,raster equal height=rows,raster column skip=6pt,raster row skip=6pt,enhanced,blankest]
\begin{tcolorbox}[enhanced,colback=poporange,colframe=poporange,coltext=white,boxrule=0.9pt,arc=3pt,left=4pt,right=4pt,top=3pt,bottom=3pt,boxsep=1pt,halign=left]\bfseries\scriptsize D\'efinition : texte fonctionnel, objectif, fid\`ele\end{tcolorbox}
\begin{tcolorbox}[enhanced,colback=popgreen,colframe=popgreen,coltext=white,boxrule=0.9pt,arc=3pt,left=4pt,right=4pt,top=3pt,bottom=3pt,boxsep=1pt,halign=left]\bfseries\scriptsize Structure externe : en-t\^ete, titre, date, signature\end{tcolorbox}
\begin{tcolorbox}[enhanced,colback=poppurple,colframe=poppurple,coltext=white,boxrule=0.9pt,arc=3pt,left=4pt,right=4pt,top=3pt,bottom=3pt,boxsep=1pt,halign=left]\bfseries\scriptsize Structure interne : introduction, corps, conclusion\end{tcolorbox}
\begin{tcolorbox}[enhanced,colback=poppink,colframe=poppink,coltext=white,boxrule=0.9pt,arc=3pt,left=4pt,right=4pt,top=3pt,bottom=3pt,boxsep=1pt,halign=left]\bfseries\scriptsize Texte explicatif : d\'efinir, d\'ecrire, expliquer\end{tcolorbox}
\begin{tcolorbox}[enhanced,colback=popgold,colframe=popgold,coltext=white,boxrule=0.9pt,arc=3pt,left=4pt,right=4pt,top=3pt,bottom=3pt,boxsep=1pt,halign=left]\bfseries\scriptsize Lexique commun / lexique sp\'ecialis\'e\end{tcolorbox}
\begin{tcolorbox}[enhanced,colback=popblueL!80!black,colframe=popblueL!80!black,coltext=white,boxrule=0.9pt,arc=3pt,left=4pt,right=4pt,top=3pt,bottom=3pt,boxsep=1pt,halign=left]\bfseries\scriptsize Textes fonctionnels et affiches de s\'ecurit\'e\end{tcolorbox}
\begin{tcolorbox}[enhanced,colback=popgreenL!80!black,colframe=popgreenL!80!black,coltext=white,boxrule=0.9pt,arc=3pt,left=4pt,right=4pt,top=3pt,bottom=3pt,boxsep=1pt,halign=left]\bfseries\scriptsize R\'edaction, relecture, am\'elioration\end{tcolorbox}
\end{tcbraster}

\legende{Carte mentale de la le\c{c}on : les sept id\'ees forces pour produire un rapport r\'eussi.}
\end{popfigure}

\begin{bilan}
\textbf{1. D\'efinitions cl\'es}
\begin{itemize}
  \item \cle{Rapport} : texte fonctionnel, \'ecrit \`a la demande d'un sup\'erieur ou d'une autorit\'e, qui rend compte de fa\c{c}on \cle{objective}, exacte et ordonn\'ee d'une mission, d'une visite, d'un stage ou d'un incident.
  \item \cle{Texte explicatif} : texte qui r\'epond aux questions \emph{comment ?} et \emph{pourquoi ?} ; il d\'efinit, d\'ecrit et explique sans donner son opinion.
  \item \cle{Lexique commun} : mots compris de tous (\emph{machine, travail, r\'eparer}).
  \item \cle{Lexique sp\'ecialis\'e} : mots propres \`a un m\'etier ou \`a une technique (\emph{disjoncteur, usinage, culasse}).
  \item \cle{Texte fonctionnel} : texte utilitaire de la vie courante (affiche, consigne, note de service) qui informe ou prescrit.
\end{itemize}

\textbf{2. Les structures \`a conna\^itre par c\oe ur}
\begin{center}
\begin{tabularx}{\linewidth}{|l|X|}
\hline
\rowcolor{popblueL} \textbf{\'El\'ement} & \textbf{Contenu obligatoire} \\
\hline
Structure externe & en-t\^ete (lieu, date), destinataire, titre/objet, nom et signature du r\'edacteur \\
\hline
Introduction & contexte, objet du rapport, circonstances (qui ? quand ? o\`u ? pourquoi ?), annonce du plan \\
\hline
Corps & faits expos\'es dans l'ordre (chronologique ou th\'ematique), observations, r\'esultats chiffr\'es \\
\hline
Conclusion & bilan, appr\'eciation mesur\'ee, suggestions ou recommandations \\
\hline
\end{tabularx}
\end{center}

\textbf{3. M\'ethodes express}
\begin{itemize}
  \item Avant d'\'ecrire : collecter les faits, prendre des notes, classer les informations.
  \item Pendant la r\'edaction : phrases courtes, pr\'esent de v\'erit\'e ou pass\'e compos\'e, connecteurs logiques (\emph{d'abord, ensuite, enfin, par cons\'equent}), aucun jugement personnel.
  \item Apr\`es la r\'edaction : relire, corriger l'orthographe et la grammaire, v\'erifier la coh\'erence, am\'eliorer le lexique (remplacer les mots vagues par le terme technique exact).
  \item Lire une affiche de s\'ecurit\'e : identifier le pictogramme, la couleur (rouge = interdiction, jaune = danger, vert = issue/secours), la consigne.
\end{itemize}
\end{bilan}

\begin{aretenir}[Checklist avant le Bac]
\begin{itemize}
  \item $\square$ Je sais d\'efinir le rapport et citer ses qualit\'es : objectivit\'e, exactitude, clart\'e.
  \item $\square$ Je connais les \'el\'ements de la structure externe (en-t\^ete, titre, date, signature).
  \item $\square$ Je sais r\'ediger une introduction qui pr\'esente le contexte et annonce le plan.
  \item $\square$ Je sais organiser le corps du rapport en parties ordonn\'ees avec des connecteurs.
  \item $\square$ Je sais conclure par un bilan et des suggestions.
  \item $\square$ Je distingue lexique commun et lexique sp\'ecialis\'e et je choisis le terme juste.
  \item $\square$ Je reconna\^is les marques du texte explicatif (d\'efinition, description, explication).
  \item $\square$ Je sais lire un texte fonctionnel et d\'ecoder une affiche de s\'ecurit\'e.
  \item $\square$ J'\'evite toute opinion personnelle et toute exag\'eration dans un rapport.
  \item $\square$ Je relis et j'am\'eliore mon texte (orthographe, ponctuation, paragraphes).
  \item $\square$ Je respecte la consigne, le nombre de mots et le temps imparti.
  \item $\square$ Je justifie ma d\'emarche et mes choix si l'\'epreuve le demande.
\end{itemize}
\end{aretenir}

\findecours

\end{document}
